\documentclass[12pt, a4paper]{article}

% ── Idioma y codificación ──────────────────────────────────────────────────
\usepackage[utf8]{inputenc}
\usepackage[T1]{fontenc}
\usepackage[spanish, es-nodecimaldot]{babel}

% ── Matemáticas ───────────────────────────────────────────────────────────
\usepackage{amsmath, amssymb, amsthm}
\usepackage{mathtools}
\usepackage{physics}

% ── Fuente ────────────────────────────────────────────────────────────────
\usepackage{lmodern}
\usepackage{microtype}

% ── Geometría de página ───────────────────────────────────────────────────
\usepackage[
  top=2.5cm, bottom=2.5cm,
  left=2.8cm, right=2.8cm
]{geometry}

% ── Encabezado y pie ──────────────────────────────────────────────────────
\usepackage{fancyhdr}
\setlength{\headheight}{15pt}
\pagestyle{fancy}
\fancyhf{}
\fancyhead[L]{\small\nouppercase{\leftmark}}
\fancyhead[R]{\small\thepage}
\renewcommand{\headrulewidth}{0.4pt}

% Portada e índice sin encabezado
\fancypagestyle{plain}{%
  \fancyhf{}%
  \renewcommand{\headrulewidth}{0pt}%
}

% ── Figuras y diagramas ───────────────────────────────────────────────────
\usepackage{tikz}
\usetikzlibrary{arrows.meta, calc, angles, quotes, decorations.pathreplacing, babel}
\usepackage{pgfplots}
\pgfplotsset{compat=1.18}
\usepackage{float}
\usepackage{caption}
\usepackage{array}
\usepackage{tabularx}
\usepackage{booktabs}
\usepackage{setspace}
\usepackage{enumitem}

% ── Estilo visual global ──────────────────────────────────────────────────
\definecolor{brandblue}{HTML}{000000}
\definecolor{softblue}{HTML}{FFFFFF}
\definecolor{linegray}{HTML}{000000}
\colorlet{blue}{black}
\colorlet{red}{black}
\colorlet{green}{black}
\colorlet{orange}{black}

\captionsetup{
  font=small,
  labelfont={bf,color=brandblue},
  textfont=it,
  justification=centering
}

\newcolumntype{Y}{>{\raggedright\arraybackslash}X}
\newcolumntype{C}{>{\centering\arraybackslash}X}
\newcolumntype{L}[1]{>{\raggedright\arraybackslash}p{#1}}
\newcolumntype{B}[1]{>{\bfseries\raggedright\arraybackslash}p{#1}}
\setlength{\tabcolsep}{7pt}
\renewcommand{\arraystretch}{1.18}

\tikzset{
  every picture/.style={line cap=round, line join=round, >=Stealth},
  every path/.append style={draw=black},
  every node/.append style={
    font=\footnotesize,
    text=black,
    fill=white,
    fill opacity=1,
    text opacity=1,
    inner sep=1.2pt
  },
  opticslabel/.style={
    fill=white,
    fill opacity=0.9,
    text opacity=1,
    rounded corners=1pt,
    inner sep=1.2pt,
    font=\footnotesize
  }
}

% ── Hipervínculos ─────────────────────────────────────────────────────────
\usepackage[hidelinks]{hyperref}

% ── Cajas para resultados destacados ──────────────────────────────────────
\usepackage{mdframed}
\usepackage[skins,breakable]{tcolorbox}
\mdfdefinestyle{resultado}{%
  linecolor=black,
  backgroundcolor=white,
  linewidth=0.8pt,
  innertopmargin=6pt,
  innerbottommargin=6pt,
  innerleftmargin=8pt,
  innerrightmargin=8pt,
  skipabove=8pt,
  skipbelow=8pt,
}

\tcbset{
  demostracionbox/.style={
    enhanced,
    breakable,
    frame hidden,
    interior hidden,
    borderline west={0.7pt}{0pt}{black},
    left=8pt, right=6pt,
    top=4pt, bottom=4pt,
    before skip=8pt, after skip=8pt,
  }
}

% ── Entornos matemáticos ──────────────────────────────────────────────────
\theoremstyle{plain}
\newtheorem{teorema}{Teorema}[section]
\newtheorem{proposicion}[teorema]{Proposición}
\newtheorem{lema}[teorema]{Lema}
\newtheorem{corolario}[teorema]{Corolario}

\theoremstyle{definition}
\newtheorem{definicion}[teorema]{Definición}
\newtheorem{ejemplo}{Ejemplo}[section]

\theoremstyle{remark}
\newtheorem*{nota}{Nota}
\newtheorem*{observacion}{Observación}

% ── Formato de secciones ──────────────────────────────────────────────────
\usepackage{titlesec}
\usepackage{etoolbox}
\titleformat{\section}
  {\large\bfseries\color{brandblue}}
  {\thesection.}{0.6em}{}
  [\vspace{2pt}\color{linegray}\titlerule\vspace{4pt}]

\titleformat{\subsection}
  {\normalsize\bfseries\color{brandblue!90!black}}
  {\thesubsection.}{0.5em}{}

\titleformat{\subsubsection}
  {\normalsize\bfseries\color{brandblue!75!black}}
  {\thesubsubsection.}{0.5em}{}

\titlespacing*{\section}{0pt}{0pt}{0.75\baselineskip}
\titlespacing*{\subsection}{0pt}{0.95\baselineskip}{0.45\baselineskip}
\titlespacing*{\subsubsection}{0pt}{0.75\baselineskip}{0.35\baselineskip}

% Nueva página automática al comenzar cada sección
\pretocmd{\section}{\clearpage}{}{}

% ── Macros de conveniencia ────────────────────────────────────────────────
\newcommand{\R}{\mathbb{R}}
\newcommand{\C}{\mathbb{C}}
\newcommand{\N}{\mathbb{N}}
\newcommand{\Z}{\mathbb{Z}}
\newcommand{\eps}{\varepsilon}
\newcommand{\ve}[1]{\vec{#1}}
\newcommand{\uvec}[1]{\hat{#1}}

% ── Separación vertical y ritmo de lectura ────────────────────────────────
\raggedbottom
\setstretch{1.08}
\setlength{\parskip}{4pt}
\setlength{\parindent}{0pt}
\setlength{\abovedisplayskip}{7pt}
\setlength{\belowdisplayskip}{7pt}
\setlength{\abovedisplayshortskip}{6pt}
\setlength{\belowdisplayshortskip}{6pt}

% ── Listas con márgenes homogéneos ────────────────────────────────────────
\setlist[itemize]{leftmargin=1.35em, itemsep=2pt, topsep=4pt, parsep=0pt, partopsep=0pt}
\setlist[enumerate]{leftmargin=1.55em, itemsep=2pt, topsep=4pt, parsep=0pt, partopsep=0pt}
\setlist[description]{style=nextline, leftmargin=0pt, labelwidth=\linewidth, labelindent=0pt, labelsep=0pt, itemindent=0pt, listparindent=0pt, itemsep=4pt, topsep=4pt, parsep=0pt, partopsep=0pt}

% ── Estilo visual de demostraciones (tipo paper) ──────────────────────────
\renewcommand{\proofname}{Demostración}
\renewcommand{\qedsymbol}{$\square$}
\BeforeBeginEnvironment{proof}{\begin{tcolorbox}[demostracionbox]}
\AfterEndEnvironment{proof}{\end{tcolorbox}}
\apptocmd{\proof}{\small\itshape}{}{}

% ── Utilidad para formulario final ────────────────────────────────────────
\newcommand{\formulaname}[1]{\textbf{#1}\par\vspace{2pt}}

% ─────────────────────────────────────────────────────────────────────────
\begin{document}

% ══════════════════════════════════════════════════════════════════════════
%  PORTADA GENERAL
% ══════════════════════════════════════════════════════════════════════════
\begin{titlepage}
  \thispagestyle{plain}
  \centering
  \vspace*{1.6cm}

  \begin{tikzpicture}[remember picture,overlay]
    \draw[line width=0.8pt]
      ($(current page.north west)+(1.6cm,-1.6cm)$)
      rectangle
      ($(current page.south east)+(-1.6cm,1.6cm)$);
  \end{tikzpicture}

  {\Large\scshape Óptica\par}
  \vspace{0.25cm}
  {\large Curso 2025--2026\par}

  \vspace{0.9cm}
  {\Huge\bfseries Apuntes Unificados\par}
  \vspace{0.4cm}
  {\large Teoría, demostraciones y formulario\par}

  \vspace{0.8cm}
  \rule{0.78\linewidth}{0.8pt}
  \vspace{0.6cm}
  {\normalsize
  Temas 1--6: ondas en vacío, interacción radiación-materia,\\
  medios materiales, refracción y reflexión, difracción e interferencia\par}
  \vspace{0.6cm}
  \rule{0.78\linewidth}{0.8pt}

  \vspace{1.2cm}
  \begin{tikzpicture}[>=Stealth, scale=1.25]
    \draw[->] (-0.3,0) -- (5.4,0) node[right] {$\ve{k}$};
    \draw[->] (0,-1.15) -- (0,1.15) node[above] {$\ve{E}$};
    \draw[domain=0:5, samples=120, thick] plot (\x, {sin(3.14159*\x r)});
    \draw[->] (0,0) -- (0,0,-1.0) node[below right] {$\ve{B}$};
    \node[opticslabel, anchor=west] at (3.55,0.68) {Onda plana armónica};
  \end{tikzpicture}

  \vfill
  \begin{tabular}{@{}rl@{}}
    \textbf{Autor:} & Luis López \\
    \textbf{Titulación:} & Grado en Física \\
    \textbf{Fecha:} & Marzo 2026
  \end{tabular}
\end{titlepage}

% ══════════════════════════════════════════════════════════════════════════
%  ÍNDICE
% ══════════════════════════════════════════════════════════════════════════
\tableofcontents

% ══════════════════════════════════════════════════════════════════════════
\section{Tema 1: Ondas Electromagnéticas en el Vacío}
% ══════════════════════════════════════════════════════════════════════════

\subsection{Introducción}

La óptica estudia la generación, propagación e interacción de la luz con la
materia. Aunque esta formulación es breve, engloba problemas de enorme
contenido físico y matemático, y ha sido una vía central de unificación en la
física moderna.

Según la escala espacial y el fenómeno de interés, se emplean tres niveles de
descripción jerárquicamente relacionados:

\begin{description}
  \item[\textbf{Óptica geométrica.}] Modela la luz como \emph{rayos} y usa el
        principio de Fermat. Es adecuada cuando $\lambda \ll a$ (longitud de
        onda mucho menor que la escala geométrica del sistema), de modo que la
        difracción y la interferencia pueden despreciarse en primera
        aproximación.

  \item[\textbf{Óptica electromagnética (u ondulatoria).}] Describe la luz como
        una \emph{onda electromagnética} que satisface las ecuaciones de
        Maxwell. Es imprescindible cuando $\lambda \sim a$, y explica de forma
        natural fenómenos como difracción, interferencia y polarización.

  \item[\textbf{Óptica cuántica.}] Introduce el carácter cuántico del campo y
        de la materia (fotones, cuantización, estadística cuántica) cuando la
        descripción clásica deja de ser suficiente, por ejemplo en emisión
        espontánea, efecto fotoeléctrico o estados no clásicos de la luz.
\end{description}

Estos marcos no compiten: la óptica geométrica es el límite de alta frecuencia
de la óptica ondulatoria, y esta aparece como límite clásico de la óptica
cuántica. En este tema adoptaremos la descripción electromagnética y fijaremos
el punto de partida formal en las ecuaciones de Maxwell.

% ══════════════════════════════════════════════════════════════════════════
\subsection{Ecuaciones de Maxwell}
% ══════════════════════════════════════════════════════════════════════════

\begin{definicion}[Ecuaciones de Maxwell]
Las ecuaciones de Maxwell, formuladas por James Clerk Maxwell entre 1861 y 1865,
constituyen el sistema completo de leyes que gobierna el comportamiento de los
campos electromagnéticos en presencia de cargas $\rho$ (densidad volumétrica de
carga, C/m$^{3}$) y corrientes $\ve{j}$ (densidad volumétrica de corriente,
A/m$^{2}$). En el Sistema Internacional y en notación diferencial, el sistema
se escribe:
\begin{align}
  \div \ve{E} &= \frac{\rho}{\varepsilon_0}, \label{eq:maxwell1}\\[4pt]
  \div \ve{B} &= 0, \label{eq:maxwell2}\\[4pt]
  \curl \ve{E} + \pdv{\ve{B}}{t} &= \ve{0}, \label{eq:maxwell3}\\[4pt]
  \frac{1}{\mu_0}\curl \ve{B} - \varepsilon_0 \pdv{\ve{E}}{t} &= \ve{j}.
  \label{eq:maxwell4}
\end{align}
Las constantes fundamentales del electromagnetismo en el vacío son la
\emph{permeabilidad magnética} $\mu_0 = 4\pi\times10^{-7}\ \mathrm{N\,s^2/C^2}$
y la \emph{permitividad eléctrica} $\varepsilon_0 = 1/(\mu_0 c^2) \approx
8{,}85\times10^{-12}\ \mathrm{F/m}$.
Los campos $\ve{E}(\ve{r},t)$ y $\ve{B}(\ve{r},t)$ son funciones vectoriales
del punto del espacio $\ve{r}$ y del tiempo $t$.

Cada ecuación cumple un papel físico específico y complementario:
\begin{itemize}
  \item \eqref{eq:maxwell1} (\textbf{Gauss eléctrica}): la divergencia de
        $\ve{E}$ queda fijada por la densidad de carga.
  \item \eqref{eq:maxwell2} (\textbf{Gauss magnética}): el campo magnético no
        tiene fuentes escalares aisladas (no hay monopolos magnéticos
        observados), de modo que $\div\ve{B}=0$.
  \item \eqref{eq:maxwell3} (\textbf{Faraday}): una variación temporal de
        $\ve{B}$ genera circulación de $\ve{E}$.
  \item \eqref{eq:maxwell4} (\textbf{Ampère-Maxwell}): corrientes de conducción
        y variaciones temporales de $\ve{E}$ generan circulación de $\ve{B}$.
\end{itemize}
\end{definicion}

\begin{definicion}[Fuerza de Maxwell-Lorentz]
La fuerza que ejercen los campos sobre una carga $q$ que se mueve con
velocidad $\ve{v}$ es:
\begin{equation}
  \ve{F} = q\ve{E} + q\ve{v}\times\ve{B} = \ve{F}_E + \ve{F}_B.
  \label{eq:lorentz}
\end{equation}
\end{definicion}

\begin{observacion}
La estructura acoplada de \eqref{eq:maxwell3} y \eqref{eq:maxwell4} es la clave
de la propagación electromagnética: una variación temporal de $\ve{B}$ induce
rotacional de $\ve{E}$ y, recíprocamente, una variación temporal de $\ve{E}$
induce rotacional de $\ve{B}$. En ausencia de fuentes locales, este acoplamiento
permite la propagación autosostenida de perturbaciones de campo.
\end{observacion}

\begin{nota}
James Clerk Maxwell (1831--1879) unificó electricidad, magnetismo y óptica
en una única teoría de campos, prediciendo teóricamente la existencia de ondas
electromagnéticas y calculando su velocidad como $c = 1/\sqrt{\mu_0\varepsilon_0}$,
que coincidía con el valor medido para la velocidad de la luz. Este fue uno de
los grandes triunfos predictivos de la física del siglo XIX. Las ecuaciones de
Maxwell son invariantes bajo transformaciones de Lorentz, lo que motivó a
Einstein el desarrollo de la relatividad especial.
\end{nota}

\subsubsection{Ecuaciones de Maxwell en el vacío}

En el vacío no hay cargas libres ni corrientes: $\rho=0$, $\ve{j}=\ve{0}$.
Este paso no es una simplificación arbitraria: separa con claridad el
\emph{problema de generación} (región con fuentes) del \emph{problema de
propagación} (región libre de fuentes). Lejos de antenas, dipolos o cargas en
movimiento, los campos satisfacen las ecuaciones homogéneas siguientes.

\begin{proposicion}[Ecuaciones de Maxwell en el vacío]
Con $\rho=0$ y $\ve{j}=\ve{0}$, las
ecuaciones~\eqref{eq:maxwell1}--\eqref{eq:maxwell4} se simplifican a:
\begin{mdframed}[style=resultado]
\begin{align}
  \div \ve{E} &= 0, \label{eq:vac1}\\[4pt]
  \div \ve{B} &= 0, \label{eq:vac2}\\[4pt]
  \curl \ve{E} + \pdv{\ve{B}}{t} &= \ve{0}, \label{eq:vac3}\\[4pt]
  \frac{1}{\mu_0}\curl \ve{B} &= \varepsilon_0 \pdv{\ve{E}}{t}. \label{eq:vac4}
\end{align}
\end{mdframed}
\end{proposicion}

% ══════════════════════════════════════════════════════════════════════════
\subsection{Ecuación de Onda}
% ══════════════════════════════════════════════════════════════════════════

\begin{definicion}[Ecuación de onda]
Una onda es una perturbación que se propaga en el espacio transportando energía
y momento sin transferencia neta de materia. El modelo matemático más sencillo
de propagación ondulatoria está dado por la \emph{ecuación de onda} o ecuación
de D'Alembert:
\begin{equation}
  \nabla^2\ve{A} = \frac{1}{v^2}\pdv[2]{\ve{A}}{t}, \label{eq:onda}
\end{equation}
donde $\ve{A}(\ve{r},t)$ es la magnitud física perturbada y $v > 0$ es la
\emph{velocidad de fase} o velocidad de propagación de la perturbación. Esta
ecuación admite como solución general cualquier función de la forma
$f(\ve{k}\cdot\ve{r} \pm \omega t)$ siempre que $\omega/|\ve{k}| = v$.
\end{definicion}

\begin{teorema}[La luz es una onda electromagnética]
Los campos $\ve{E}$ y $\ve{B}$ que satisfacen las ecuaciones de Maxwell en
el vacío~\eqref{eq:vac1}--\eqref{eq:vac4} cumplen la ecuación de
onda~\eqref{eq:onda} con velocidad $c = 1/\sqrt{\mu_0\varepsilon_0}
\approx 3\times10^8\ \mathrm{m/s}$:
\begin{mdframed}[style=resultado]
\begin{equation}
  \nabla^2\ve{E} = \frac{1}{c^2}\pdv[2]{\ve{E}}{t}, \qquad
  \nabla^2\ve{B} = \frac{1}{c^2}\pdv[2]{\ve{B}}{t}. \label{eq:ondaEM}
\end{equation}
\end{mdframed}
\end{teorema}

\begin{proof}
\textbf{Paso 1 (ecuación para $\ve{E}$).} Aplicamos $\curl$ a
\eqref{eq:vac3}:
\[
  \curl(\curl\ve{E})+\curl\!\left(\pdv{\ve{B}}{t}\right)=\ve{0}.
\]
Bajo hipótesis de regularidad, $\curl$ conmuta con $\partial/\partial t$:
\[
  \curl(\curl\ve{E})+\pdv{}{t}(\curl\ve{B})=\ve{0}.
\]
Usamos la identidad del doble rotacional
$\curl(\curl\ve{F})=\grad(\div\ve{F})-\nabla^2\ve{F}$:
\[
  \grad(\div\ve{E})-\nabla^2\ve{E}+\pdv{}{t}(\curl\ve{B})=\ve{0}.
\]
Como en vacío $\div\ve{E}=0$ por \eqref{eq:vac1}, desaparece el término de
gradiente. Sustituyendo $\curl\ve{B}=\mu_0\varepsilon_0\,\partial_t\ve{E}$ por
\eqref{eq:vac4}:
\[
  -\nabla^2\ve{E}+\mu_0\varepsilon_0\,\pdv[2]{\ve{E}}{t}= \ve{0}
  \quad\Longrightarrow\quad
  \nabla^2\ve{E}-\frac{1}{c^2}\pdv[2]{\ve{E}}{t}= \ve{0},
\]
donde $c^{-2}=\mu_0\varepsilon_0$.

\textbf{Paso 2 (ecuación para $\ve{B}$).} Partimos de \eqref{eq:vac4}:
\[
  \curl\ve{B}-\mu_0\varepsilon_0\pdv{\ve{E}}{t}=\ve{0}.
\]
Aplicamos $\curl$:
\[
  \curl(\curl\ve{B})-\mu_0\varepsilon_0\pdv{}{t}(\curl\ve{E})=\ve{0}.
\]
De nuevo, identidad del doble rotacional:
\[
  \grad(\div\ve{B})-\nabla^2\ve{B}
  -\mu_0\varepsilon_0\pdv{}{t}\!\left(\pdv{\ve{B}}{t}\right)=\ve{0},
\]
donde usamos \eqref{eq:vac3}, $\curl\ve{E}=-\partial_t\ve{B}$. Como
$\div\ve{B}=0$ por \eqref{eq:vac2}, resulta
\[
  \nabla^2\ve{B}-\frac{1}{c^2}\pdv[2]{\ve{B}}{t}= \ve{0}.
\]
\end{proof}

\begin{observacion}
Este teorema fija tres conclusiones estructurales para el resto del curso:
\begin{enumerate}
  \item La radiación electromagnética se propaga como onda.
  \item La propagación en vacío es posible sin medio material soporte.
  \item La onda es vectorial: la polarización forma parte de su descripción
        física esencial.
\end{enumerate}
Además, la velocidad
$c=1/\sqrt{\mu_0\varepsilon_0}\approx2{,}998\times10^8\ \mathrm{m/s}$ queda
determinada por constantes electromagnéticas, cerrando la identificación entre
luz y campo electromagnético.
\end{observacion}

% ══════════════════════════════════════════════════════════════════════════
\subsection{Ondas Armónicas o Monocromáticas}
% ══════════════════════════════════════════════════════════════════════════

\begin{definicion}[Onda armónica]
Una onda armónica (o monocromática) es aquella cuya dependencia temporal es
estrictamente periódica y sinusoidal con una única frecuencia bien definida.
Sus componentes cartesianas del campo son de la forma:
\begin{equation}
  E_j(\ve{r},t) = A_j(\ve{r})\cos\!\bigl(\omega t - \delta_j(\ve{r})\bigr),
  \qquad j = x,y,z, \label{eq:armonicacomp}
\end{equation}
donde $A_j(\ve{r}) \geq 0$ es la \emph{amplitud} (que puede depender de la
posición, por ejemplo en una onda esférica), $\delta_j(\ve{r})$ es la
\emph{fase inicial} en el punto $\ve{r}$, $\omega = 2\pi\nu$ la
\emph{frecuencia angular} (rad/s), $\nu = 1/T$ la \emph{frecuencia} (Hz) y
$T$ el \emph{periodo temporal}.

La monocromaticidad estricta es una idealización (fuente infinita en tiempo y
con ancho espectral nulo). En la práctica toda fuente posee banda finita, pero
la aproximación armónica es extremadamente útil cuando dicho ancho es pequeño
frente a las escalas espectrales del problema.
\end{definicion}

\begin{proposicion}[Condición de armonicidad]
Una onda es armónica de frecuencia $\omega$ si y solo si satisface:
\begin{equation}
  \pdv[2]{E}{t} = -\omega^2 E. \label{eq:armonicacond}
\end{equation}
Esta condición es la base del tratamiento frecuencial: en ecuaciones lineales
con dependencia temporal armónica, $\partial_t$ puede sustituirse por
$i\omega$ (o $-i\omega$, según convención), reduciendo el problema diferencial
a uno algebraico en amplitudes complejas.
\end{proposicion}

\begin{proof}
Derivando dos veces respecto al tiempo la expresión~\eqref{eq:armonicacomp}:
\[
  \pdv{E_j}{t} = -\omega A_j\sin(\omega t-\delta_j),
  \qquad
  \pdv[2]{E_j}{t}
  = -\omega^2 A_j\cos(\omega t-\delta_j)
  = -\omega^2 E_j.
\]
Esto demuestra la implicación directa.

Recíprocamente, si el campo satisface~\eqref{eq:armonicacond} para algún
$\omega > 0$, la dependencia temporal obedece la ecuación diferencial ordinaria:
\[
  \ddot E + \omega^2 E = 0.
\]
La ecuación característica asociada es:
\[
  r^2+\omega^2=0 \quad\Rightarrow\quad r=\pm i\omega,
\]
con raíces puramente imaginarias conjugadas. La solución general real es:
\[
  E(t)=C_1\cos(\omega t)+C_2\sin(\omega t)
  = A\cos(\omega t-\delta),
\]
donde $A = \sqrt{C_1^2+C_2^2}$ y $\tan\delta = C_2/C_1$, lo que confirma la
forma~\eqref{eq:armonicacomp}.
\end{proof}

\begin{observacion}
La relevancia de las ondas armónicas no es solo técnica: por el formalismo de
Fourier, señales físicamente admisibles pueden descomponerse en superposición
de componentes monocromáticas. Por linealidad, resolver el problema armónico
permite reconstruir la evolución de señales generales.
\end{observacion}

\subsubsection{Espectro electromagnético}

Las ondas electromagnéticas armónicas se clasifican según su frecuencia $\nu$
(o equivalentemente su longitud de onda en el vacío $\lambda = c/\nu$), y
conforman el \emph{espectro electromagnético}: desde las ondas de radio
($\nu \sim 10^6\ \mathrm{Hz}$, $\lambda \sim 300\ \mathrm{m}$) hasta los
rayos gamma ($\nu \sim 10^{21}\ \mathrm{Hz}$, $\lambda \sim 10^{-12}\ \mathrm{m}$).
La relación fundamental entre longitud de onda y frecuencia en el vacío es:
\[
  \lambda\nu = c.
\]

\begin{definicion}[Espectro visible]
La \emph{luz visible} es la pequeña fracción del espectro electromagnético
detectable por el ojo humano. Corresponde al intervalo de frecuencias:
\begin{equation}
  4\times10^{14}\ \mathrm{Hz} \leq \nu \leq 7{,}5\times10^{14}\ \mathrm{Hz},
  \qquad
  400\ \mathrm{nm} \leq \lambda \leq 750\ \mathrm{nm}. \label{eq:espectro}
\end{equation}
Las distintas longitudes de onda dentro de este rango son percibidas como
colores diferentes por el sistema visual humano, como se muestra en la
tabla siguiente.
\end{definicion}

\begin{center}
\begin{tabular}{@{}cc@{}}
  \toprule
  $\lambda$ (nm) & Color \\
  \midrule
  450 & Violeta \\
  490 & Azul    \\
  555 & Verde   \\
  589 & Amarillo \\
  630 & Rojo    \\
  \bottomrule
\end{tabular}
\end{center}

% ══════════════════════════════════════════════════════════════════════════
\subsection{Representación Compleja de una Onda Armónica}
% ══════════════════════════════════════════════════════════════════════════

\begin{definicion}[Representación compleja]
Por la identidad de Euler, toda onda armónica real puede representarse como
parte real de una cantidad compleja. Para el campo eléctrico:
\begin{equation}
  E_j(\ve{r},t)
  = \operatorname{Re}\!\left[A_j(\ve{r})\,e^{-i\delta_j(\ve{r})}\,e^{i\omega t}\right]
  = \operatorname{Re}\!\left[\widetilde{E}_j(\ve{r})\,e^{i\omega t}\right],
  \label{eq:compleja}
\end{equation}
donde:
\begin{itemize}
  \item $E_j(\ve{r},t)$ es el \textbf{campo real} medible;
  \item $\widetilde{E}_j(\ve{r})=A_j(\ve{r})e^{-i\delta_j(\ve{r})}$ es la
        \textbf{amplitud compleja} (fasor espacial);
  \item $\widetilde{E}_j(\ve{r})\,e^{i\omega t}$ es la representación compleja
        temporal completa.
\end{itemize}
Durante el cálculo se opera con la representación compleja y al final se toma
$\Re\{\cdot\}$ para recuperar la magnitud física.
\end{definicion}

\begin{proposicion}[Validez y limitaciones de la representación compleja]
\begin{itemize}
  \item[\checkmark] Es exacta para operadores lineales (suma, derivación,
        integración, ecuaciones lineales con coeficientes reales).
  \item[\checkmark] En régimen armónico, $\partial_t \mapsto i\omega$ simplifica
        de forma sistemática la resolución.
  \item[$\boldsymbol{\times}$] No puede aplicarse de forma ingenua en magnitudes
        no lineales (productos y módulos cuadrados): intensidad, energía media
        y vector de Poynting requieren fórmulas de promedio adecuadas o volver
        al campo real.
\end{itemize}
\end{proposicion}

\subsubsection{Repaso de números complejos}

Sea $z = a + ib$:
\[
  \bar{z} = a-ib,\quad
  |z|^2 = a^2+b^2,\quad
  \operatorname{Re}\{z\} = \frac{z+\bar{z}}{2},\quad
  \operatorname{Im}\{z\} = \frac{z-\bar{z}}{2i}.
\]
Forma polar: $z = |z|\,e^{i\alpha}$ con $\tan\alpha = b/a$. Fórmulas de Euler:
\[
  e^{i\alpha} = \cos\alpha + i\sin\alpha, \qquad
  \cos\alpha = \frac{e^{i\alpha}+e^{-i\alpha}}{2}, \qquad
  \sin\alpha = \frac{e^{i\alpha}-e^{-i\alpha}}{2i}.
\]

\subsubsection{Ejemplos de identificación de armonicidad}

\begin{ejemplo}
Determinar cuáles de las siguientes ondas son armónicas usando la
condición~\eqref{eq:armonicacond}: $\partial^2E/\partial t^2 = -\omega^2 E$.

\begin{enumerate}

  \item $E = A\cos(\omega_1 t - kz) + B\sin(\omega_2 t - kz)$

  Si $\omega_1=\omega_2\equiv\omega$:
  \[
    \pdv[2]{E}{t} = -\omega^2(A\cos + B\sin) = -\omega^2 E.
  \]
  \checkmark\ \textbf{Armónica.} \quad
  Si $\omega_1\neq\omega_2$: no existe $\omega$ común $\Rightarrow$ \textbf{No armónica.}

  \item $E = A\,e^{\cos(\omega t-kz)}$
  \[
    \pdv[2]{E}{t} = e^{\cos(\cdots)}\!\bigl(-\omega^2\cos + \omega^2\sin^2\bigr)
    \neq -\omega^2 E.
  \]
  \textbf{No armónica.}

  \item $E = \cos(\omega t)\cos(kz)$
  \[
    \pdv[2]{E}{t} = -\omega^2\cos(\omega t)\cos(kz) = -\omega^2 E.
  \]
  \checkmark\ \textbf{Armónica} (onda estacionaria).

  \item $E = A\sin(\omega t - \delta)$
  \[
    \pdv[2]{E}{t} = -\omega^2 A\sin(\omega t-\delta) = -\omega^2 E.
  \]
  \checkmark\ \textbf{Armónica.}

\end{enumerate}
\end{ejemplo}

% ══════════════════════════════════════════════════════════════════════════
\subsection{Frentes de Onda}
% ══════════════════════════════════════════════════════════════════════════

\begin{definicion}[Frente de onda]
El \emph{frente de onda} es el lugar geométrico de todos los puntos del espacio
que, en un instante $t_0$ fijo, presentan la misma fase instantánea. Para una
onda armónica de la forma~\eqref{eq:armonicacomp}, la condición de fase constante
es:
\begin{equation}
  \omega t_0 - \delta(\ve{r}) = \mathrm{cte}. \label{eq:frente}
\end{equation}
El frente de onda es, por tanto, una superficie de nivel de la función de fase
$\delta(\ve{r})$. La dirección de propagación de la onda es, en cada punto, la
perpendicular a este frente de onda (es decir, la dirección del gradiente
$\nabla\delta$).
\end{definicion}

Los tipos de frente de onda quedan completamente determinados por la forma
funcional de $\delta(\ve{r})$:

\medskip
\begin{center}
\begin{tabularx}{0.95\linewidth}{@{}L{2.1cm}Y Y@{}}
  \toprule
  Tipo & $\delta(\ve{r})$ & Campo \\
  \midrule
  Plano    & $kz$                     & $\ve{E}_0\,e^{i\omega t}e^{-ikz}$ \\
  Esférico & $k\sqrt{x^2+y^2+z^2}$   & $\dfrac{\ve{E}_0}{r}\,e^{i\omega t}e^{-ikr}$ \\
  General  & $\delta(\ve{r})$ arbitraria & cualquier superficie de fase cte \\
  \bottomrule
\end{tabularx}
\end{center}

\begin{figure}[H]
  \centering
  \begin{tikzpicture}[>=Stealth, scale=1.1]
    % Onda plana
    \begin{scope}[xshift=-3.8cm]
      \draw[->] (-0.3,0) -- (3.3,0) node[right] {$z$};
      \foreach \x in {0.4, 1.1, 1.8, 2.5}{
        \draw[thick] (\x,-1.1) -- (\x,1.1);
      }
      \draw[->, thick] (0.75,0.5) -- (1.45,0.5)
            node[midway, above=2pt] {$\uvec{k}$};
    \end{scope}
    % Onda esférica
    \begin{scope}[xshift=3.2cm]
      \filldraw (0,0) circle (2pt) node[below left=2pt] {$O$};
      \foreach \r in {0.55, 1.0, 1.5}{
        \draw[thick] (0,0) circle (\r);
      }
      \draw[->, thick] (0,0) -- (0.72,0.72)
            node[midway, above left] {$r$};
    \end{scope}
    \node[align=left, anchor=west] at (5.1,0.5)
      {\footnotesize Izquierda: frente plano\\
       \footnotesize Derecha: frente esférico};
  \end{tikzpicture}
  \caption{Frentes de onda plano (izquierda) y esférico (derecha).}
  \label{fig:frentes}
\end{figure}

% ══════════════════════════════════════════════════════════════════════════
\subsection{Ondas Planas}
% ══════════════════════════════════════════════════════════════════════════

\begin{definicion}[Onda plana]
Una onda plana es aquella en la que el campo tiene el mismo valor (misma amplitud
y misma fase) para todos los puntos de cualquier plano perpendicular a la
dirección de propagación $\uvec{k}$. La expresión compleja es:
\begin{equation}
  \ve{E}(\ve{r},t) = \ve{E}_0\,e^{i\omega t}\,e^{-i\ve{k}\cdot\ve{r}},
  \label{eq:ondaplana}
\end{equation}
donde $\ve{k}=(k_x,k_y,k_z)$ es el \emph{vector de onda} —cuya dirección es
la de propagación y cuyo módulo $k = |\ve{k}|$ es el \emph{número de onda}—,
y $\ve{E}_0$ es la amplitud vectorial, independiente de $\ve{r}$ y de $t$.
Los frentes de onda son los planos $\ve{k}\cdot\ve{r} = \mathrm{cte}$.

En la práctica, las ondas planas son una idealización (implicarían energía
infinita extendida por todo el espacio), pero constituyen la solución
matemáticamente más sencilla de la ecuación de onda y forman una base completa
para construir haces más realistas mediante la integral de Fourier.
\end{definicion}

\begin{definicion}[Velocidad de fase]
La \emph{velocidad de fase} es la velocidad a la que se propagan los frentes
de fase constante de la onda. Para una onda propagándose en la dirección $z$
con $\ve{k} = k\,\uvec{z}$:
\[
  \omega t_0 - kz = \mathrm{cte} \implies z = \frac{\omega}{k}\,t
  + \mathrm{cte},
\]
de modo que:
\begin{equation}
  v_f = \pdv{z}{t} = \frac{\omega}{k} = c \quad(\text{en el vacío}).
  \label{eq:vfase}
\end{equation}
En medios materiales, $v_f = c/n$, donde $n$ es el índice de refracción del
medio (véase el Tema~3). La velocidad de fase puede en principio superar $c$
sin contradecir la relatividad especial, pues no transporta información.
\end{definicion}

\subsubsection{Ondas homogéneas e inhomogéneas}

\begin{description}
  \item[\textbf{Homogénea.}] Los planos de fase constante y los de amplitud
        constante coinciden (son los mismos planos). La amplitud $\ve{E}_0$
        es uniforme en todo el frente de onda. Este es el caso de una onda
        plana propagándose en un medio sin pérdidas.

  \item[\textbf{Inhomogénea.}] Los planos de fase constante y los de amplitud
        constante son distintos entre sí:
        \[
          \ve{E}(\ve{r},t) = \ve{E}_0\,e^{i\omega t}\,e^{-ikz}\,e^{-ax},
        \]
        donde la amplitud efectiva $E_0\,e^{-ax}$ decae en la dirección $x$,
        mientras la fase evoluciona en la dirección $z$. Las ondas evanescentes
        (asociadas a la reflexión total interna) son el ejemplo más importante
        de este tipo.
\end{description}

% ══════════════════════════════════════════════════════════════════════════
\subsection{Ondas Armónicas Planas}
% ══════════════════════════════════════════════════════════════════════════

\begin{definicion}[Onda armónica plana (OAP)]
Una OAP es la combinación de una onda armónica y una onda plana:
\begin{equation}
  \ve{E}(\ve{r},t) = \ve{E}_0\,e^{i\omega t}\,e^{-i\ve{k}\cdot\ve{r}},
  \qquad
  \ve{B}(\ve{r},t) = \ve{B}_0\,e^{i\omega t}\,e^{-i\ve{k}\cdot\ve{r}},
  \label{eq:OAP}
\end{equation}
con $\ve{E}_0$ y $\ve{B}_0$ independientes de $\ve{r}$ y $t$.
\end{definicion}

\subsubsection{Condiciones impuestas por las ecuaciones de Maxwell}

Una vez propuesto el ansatz de onda plana armónica~\eqref{eq:OAP}, los
operadores diferenciales actúan de forma algebraica. Aplicando las reglas
$\partial/\partial t \to i\omega$ y $\nabla \to -i\ve{k}$ (que se siguen
directamente al diferenciar $e^{i\omega t - i\ve{k}\cdot\ve{r}}$):
\[
  \pdv{\ve{F}}{t} \;\to\; i\omega\ve{F}, \qquad
  \div\ve{F} \;\to\; -i\,\ve{k}\cdot\ve{F}, \qquad
  \curl\ve{F} \;\to\; -i\,\ve{k}\times\ve{F}.
\]
Esta sustitución es exacta para ondas planas y transforma el sistema de
ecuaciones en derivadas parciales de Maxwell en el sistema algebraico:
Sustituyendo en las ecuaciones~\eqref{eq:vac1}--\eqref{eq:vac4}:

\medskip
\begin{center}
\begin{tabularx}{0.95\linewidth}{@{}L{4.8cm}Y@{}}
  \toprule
  Ecuación de Maxwell & Resultado \\
  \midrule
  $\div\ve{E}=0$                         & $\ve{k}\cdot\ve{E}=0$ \\
  $\div\ve{B}=0$                         & $\ve{k}\cdot\ve{B}=0$ \\
  $\curl\ve{E}+\partial_t\ve{B}=\ve{0}$  & $\omega\ve{B}=\ve{k}\times\ve{E}$ \\
  $(1/\mu_0)\curl\ve{B}=\varepsilon_0\partial_t\ve{E}$ &
  $\ve{k}\times\ve{B}=\mu_0\varepsilon_0\omega\ve{E}$ \\
  \bottomrule
\end{tabularx}
\end{center}

\subsubsection{Propiedades fundamentales de las OAP}

\begin{proposicion}[Transversalidad]
En una OAP, los vectores $\ve{E}$, $\ve{B}$ y $\ve{k}$ forman un triedro
mutuamente ortogonal:
\begin{mdframed}[style=resultado]
\begin{equation}
  \ve{k}\perp\ve{E}, \quad \ve{k}\perp\ve{B}, \quad \ve{E}\perp\ve{B}.
  \label{eq:transversal}
\end{equation}
\end{mdframed}
La luz es una onda \textbf{transversal}: las oscilaciones del campo ocurren
en el plano perpendicular a la dirección de propagación, lo que es consecuencia
directa de las ecuaciones de Maxwell en el vacío.
\end{proposicion}

\begin{proof}
Las ecuaciones de divergencia~\eqref{eq:vac1} y~\eqref{eq:vac2} aplicadas
a la OAP dan, en forma algebraica:
\[
  \ve{k}\cdot\ve{E}=0,
  \qquad
  \ve{k}\cdot\ve{B}=0.
\]
Un vector cuyo producto escalar con $\ve{k}$ es nulo es, por definición,
ortogonal a $\ve{k}$. Por tanto, $\ve{k}\perp\ve{E}$ y $\ve{k}\perp\ve{B}$.

De la ley de Faraday en forma algebraica:
\[
  \omega\ve{B}=\ve{k}\times\ve{E}.
\]
El producto vectorial $\ve{k}\times\ve{E}$ es perpendicular a $\ve{E}$ (propiedad
del producto vectorial). Luego $\ve{B}\perp\ve{E}$.

Además, $\ve{B}$ y $\ve{k}$ son ya mutuamente perpendiculares, y $\ve{B}$ es
paralelo a $\ve{k}\times\ve{E}$, lo que confirma que los tres vectores forman
un triedro ortonormal (salvo factor de escala): $\ve{E}$, $\ve{B}$,
$\uvec{k} = \ve{k}/k$.
\end{proof}

\begin{observacion}
La transversalidad implica que la dinámica del campo eléctrico ocurre en el
plano perpendicular a $\ve{k}$; por eso la polarización se describe en ese
plano y no en la dirección de propagación.
\end{observacion}

\begin{proposicion}[Relación de dispersión]
El módulo del vector de onda de una OAP en el vacío satisface:
\begin{equation}
  k = \frac{\omega}{c} = \frac{2\pi}{\lambda}, \qquad \lambda\nu = c.
  \label{eq:dispersion}
\end{equation}
donde $k$ es el número de onda (rad/m), $\lambda$ la longitud de onda (m) y $c$
la velocidad de la luz en el vacío. Esta relación, denominada \emph{relación de
dispersión}, no es una hipótesis sino una consecuencia de las ecuaciones de
Maxwell: la onda plana armónica es solución si y solo si $k = \omega/c$.
\end{proposicion}

\begin{proof}
Partimos de:
\[
  \omega\ve{B}=\ve{k}\times\ve{E}.
\]
Aplicamos $\ve{k}\times$ y usamos:
\[
  \ve{k}\times\ve{B}=\mu_0\varepsilon_0\omega\ve{E}.
\]
\[
  \ve{k}\times(\ve{k}\times\ve{E})
  = \ve{k}(\ve{k}\cdot\ve{E})
    - \ve{E}(\ve{k}\cdot\ve{k})
  = \mu_0\varepsilon_0\omega^2\ve{E}.
  \tag*{\normalfont\footnotesize(en OAP, $\ve{k}\cdot\ve{E}=0$)}
\]
Luego:
\[
  -k^2\ve{E}=\mu_0\varepsilon_0\omega^2\ve{E}
  \quad\Rightarrow\quad
  k^2=\mu_0\varepsilon_0\omega^2.
\]
Usando:
\[
  c^{-2}=\mu_0\varepsilon_0,
\]
\[
  k^2=\frac{\omega^2}{c^2}
  \quad\Rightarrow\quad
  k=\frac{\omega}{c}.
\]
\end{proof}

\begin{observacion}
La ecuación \eqref{eq:dispersion} fija la cinemática de fase en el vacío:
frecuencia y número de onda no son independientes para una OAP física.
\end{observacion}

\begin{proposicion}[Relación entre amplitudes]
Las amplitudes de los campos eléctrico y magnético en una OAP satisfacen:
\begin{equation}
  E = c\,B. \label{eq:amplitudes}
\end{equation}
Dado que $c \approx 3\times10^8\ \mathrm{m/s}$, el campo magnético es unas
$3\times10^8$ veces más pequeño (en unidades SI) que el campo eléctrico. Por
este motivo, en muchos problemas de óptica se trabaja únicamente con $\ve{E}$,
que es también el responsable de la mayoría de los efectos en los detectores
(fotodíodos, retina, etc.).
\end{proposicion}

\begin{proof}
De la ley de Faraday en forma algebraica para la OAP:
\[
  \omega\ve{B}=\ve{k}\times\ve{E}.
\]
Como $\ve{k}\perp\ve{E}$, el módulo del producto vectorial es
$|\ve{k}\times\ve{E}| = k\,E\sin(90^\circ) = kE$, luego:
\[
  \omega B = kE.
\]
Usando la relación de dispersión $k = \omega/c$:
\[
  B = \frac{k}{\omega}E = \frac{E}{c},
\]
lo que equivale a~\eqref{eq:amplitudes}.
\end{proof}

\begin{observacion}
El cociente $E/B=c$ expresa que ambos campos contienen la misma información
ondulatoria; en óptica experimental suele bastar con medir $\ve{E}$ y deducir
$\ve{B}$ a partir de \eqref{eq:amplitudes}.
\end{observacion}

\begin{figure}[H]
  \centering
  \begin{tikzpicture}[>=Stealth, scale=1.4]
    \coordinate (O) at (0,0);
    \draw[->, very thick] (O) -- (2.8,0)  node[right]       {$\ve{k}$};
    \draw[->, very thick] (O) -- (0,1.7)  node[above]       {$\ve{E}$};
    \draw[->, very thick] (O) -- (-1.2,-0.8) node[below left] {$\ve{B}$};
    \draw (0.28,0) -- (0.28,0.28) -- (0,0.28);
    \filldraw (O) circle (1.5pt);
    \node[align=left, anchor=west] at (3.25,1.25)
      {\footnotesize $\ve{E}\perp\ve{k}$\\
       \footnotesize $\ve{B}\perp\ve{k}$\\
       \footnotesize $\ve{E}\perp\ve{B}$};
  \end{tikzpicture}
  \caption{Triedro ortogonal de una OAP.}
  \label{fig:triedro}
\end{figure}

% ══════════════════════════════════════════════════════════════════════════
\subsection{Polarización de las Ondas Armónicas Planas}
% ══════════════════════════════════════════════════════════════════════════

\begin{definicion}[Polarización]
La \emph{polarización} de una OAP es la trayectoria que describe el extremo de
$\ve{E}$ en un plano transversal fijo. Para propagación según $z$,
\eqref{eq:transversal} impone $E_z=0$, de modo que la trayectoria se estudia en
el plano $XY$. Tomando $\delta_x=0$ y $\delta=\delta_y-\delta_x$:
\begin{equation}
  E_x = E_{0x}\cos(\omega t - kz), \qquad
  E_y = E_{0y}\cos(\omega t - kz + \delta). \label{eq:polcomps}
\end{equation}
El estado de polarización queda determinado por el cociente
$E_{0y}/E_{0x}$ y la fase relativa $\delta\in(-\pi,\pi]$.
\end{definicion}

\subsubsection{Elipse de polarización}

\begin{proposicion}[Ecuación de la elipse]
Eliminando el argumento temporal de~\eqref{eq:polcomps}:
\begin{equation}
  \left(\frac{E_x}{E_{0x}}\right)^{\!2}
  + \left(\frac{E_y}{E_{0y}}\right)^{\!2}
  - 2\,\frac{E_x E_y}{E_{0x}E_{0y}}\cos\delta
  = \sin^2\!\delta. \label{eq:elipse}
\end{equation}
\end{proposicion}

\begin{proof}
Definimos $\phi=\omega t-kz$. De \eqref{eq:polcomps}:
\[
  \frac{E_x}{E_{0x}}=\cos\phi,
  \qquad
  \frac{E_y}{E_{0y}}=\cos(\phi+\delta)
  =\cos\phi\cos\delta-\sin\phi\sin\delta.
\]
\textbf{Sustitución y despeje:}
\[
  \frac{E_y}{E_{0y}}-\frac{E_x}{E_{0x}}\cos\delta
  =-\sin\phi\sin\delta.
\]
\[
  \sin^2\phi
  =\frac{1}{\sin^2\delta}
   \left(
     \frac{E_y}{E_{0y}}-\frac{E_x}{E_{0x}}\cos\delta
   \right)^2.
\]
\textbf{Identidad trigonométrica:}
\[
  \sin^2\phi+\cos^2\phi=1
  \quad\Longrightarrow\quad
  1-\left(\frac{E_x}{E_{0x}}\right)^2
  =\frac{1}{\sin^2\delta}
   \left(
     \frac{E_y}{E_{0y}}-\frac{E_x}{E_{0x}}\cos\delta
   \right)^2.
\]
Multiplicando por $\sin^2\delta$, expandiendo y usando
$\sin^2\delta+\cos^2\delta=1$, se obtiene exactamente
\eqref{eq:elipse}.
\end{proof}

\subsubsection{Vector de Jones}

\begin{definicion}[Vector de Jones]
Toda la información sobre el estado de polarización de una OAP monocromática
queda codificada en el \emph{vector de Jones}, que es un vector columna complejo
de dos componentes:
\begin{equation}
  \ve{J} = \begin{pmatrix}E_{0x}\\E_{0y}\,e^{i\delta}\end{pmatrix}.
  \label{eq:jones}
\end{equation}
Los módulos de las componentes dan las amplitudes en cada dirección transversal,
y la fase relativa $\delta$ determina la forma de la elipse de polarización.
Dos vectores de Jones que difieran solo en un factor de fase global $e^{i\phi}$
describen el mismo estado de polarización (son físicamente equivalentes).
\end{definicion}

\subsubsection{Tipos de polarización}

\paragraph{Polarización lineal}

\begin{proposicion}
La luz está linealmente polarizada cuando $\delta = m\pi$ ($m\in\Z$).
El campo vibra en una dirección fija de azimut $\alpha$:
\begin{equation}
  \tan\alpha = \frac{E_{0y}}{E_{0x}}, \qquad
  \ve{J} = \begin{pmatrix}\cos\alpha\\\pm\sin\alpha\end{pmatrix}.
  \label{eq:jones_lineal}
\end{equation}
Casos particulares normalizados:
\[
  \alpha=0^{\circ}:\;\ve{J}=\begin{pmatrix}1\\0\end{pmatrix},\quad
  \alpha=90^{\circ}:\;\ve{J}=\begin{pmatrix}0\\1\end{pmatrix},\quad
  \alpha=45^{\circ}:\;\ve{J}=\tfrac{1}{\sqrt{2}}\begin{pmatrix}1\\1\end{pmatrix}.
\]
\end{proposicion}

\begin{proof}
Tomamos la condición:
\[
  \delta=m\pi.
\]
Entonces, la ecuación~\eqref{eq:elipse} se reduce a:
\[
  \left(\frac{E_y}{E_{0y}} \mp \frac{E_x}{E_{0x}}\right)^2 = 0.
\]
De ahí:
\[
  \frac{E_y}{E_x} = \pm\frac{E_{0y}}{E_{0x}}.
\]
La pendiente de esa recta es:
\[
  \tan\alpha=\pm\frac{E_{0y}}{E_{0x}}.
\]
\end{proof}
Geométricamente, el extremo de $\ve{E}$ recorre una recta fija en el plano
transversal.

\paragraph{Polarización circular}

\begin{proposicion}
La luz está circularmente polarizada cuando $E_{0x}=E_{0y}=E_0$ y
$\delta=(2m+1)\pi/2$. La elipse degenera en una circunferencia:
\begin{equation}
  \text{Dextrógira }(\delta=+\tfrac{\pi}{2}):\;
    \ve{J}=E_0\!\begin{pmatrix}1\\i\end{pmatrix};\qquad
  \text{Levógira }(\delta=-\tfrac{\pi}{2}):\;
    \ve{J}=E_0\!\begin{pmatrix}1\\-i\end{pmatrix}.
  \label{eq:jones_circ}
\end{equation}
\end{proposicion}

\begin{proof}
Tomamos:
\[
  E_{0x}=E_{0y}=E_0,
  \qquad
  \cos\delta=0.
\]
Entonces, la ecuación~\eqref{eq:elipse} queda:
\[
  \frac{E_x^2}{E_0^2} + \frac{E_y^2}{E_0^2} = 1.
\]
Esta es una circunferencia de radio $E_0$.
\end{proof}
Geométricamente, el extremo de $\ve{E}$ recorre una circunferencia y el
sentido de giro viene dado por el signo de $\delta$.

\paragraph{Polarización elíptica}

\begin{proposicion}
La luz está elípticamente polarizada cuando $E_{0x}\neq E_{0y}$ y
$\delta=(2m+1)\pi/2$:
\begin{equation}
  \text{Dextrógira:}\;
    \ve{J}=\begin{pmatrix}E_{0x}\\iE_{0y}\end{pmatrix};\qquad
  \text{Levógira:}\;
    \ve{J}=\begin{pmatrix}E_{0x}\\-iE_{0y}\end{pmatrix}.
  \label{eq:jones_elip}
\end{equation}
\end{proposicion}
Geométricamente, el extremo de $\ve{E}$ recorre una elipse cuyos semiejes y
helicidad dependen del cociente de amplitudes y de la fase relativa.

\paragraph{Luz despolarizada}

\begin{definicion}[Luz natural]
La luz natural o despolarizada es aquella en la que las amplitudes $E_{0x}$,
$E_{0y}$ y la diferencia de fase $\delta(t)$ fluctúan rápida y aleatoriamente
en el tiempo, de forma que el promedio temporal de cualquier estado de
polarización definido es nulo. No puede describirse con un vector de Jones fijo.
Matemáticamente, requiere del formalismo de la \emph{matriz de coherencia} o
de los \emph{parámetros de Stokes}, que permiten caracterizar estadísticamente
el estado de polarización.

La \emph{luz parcialmente polarizada} es superposición incoherente de luz
completamente polarizada (descrita por un vector de Jones) y luz despolarizada;
el grado de polarización $p = I_\text{pol}/I_\text{total}$ (con $0 \leq p \leq 1$)
cuantifica la fracción polarizada.
\end{definicion}

\begin{figure}[H]
  \centering
  \begin{tikzpicture}[>=Stealth, scale=1.05]
    % Lineal
    \begin{scope}[xshift=-5.5cm]
      \draw[->] (-1.3,0)--(1.3,0) node[right]{$x$};
      \draw[->] (0,-1.3)--(0,1.3) node[above]{$y$};
      \draw[thick, <->] (-0.9,-0.9)--(0.9,0.9);
      \node[below=0.25cm] at (0,-1.3) {\footnotesize Lineal ($45^\circ$)};
    \end{scope}
    % Circular
    \begin{scope}[xshift=0cm]
      \draw[->] (-1.3,0)--(1.3,0) node[right]{$x$};
      \draw[->] (0,-1.3)--(0,1.3) node[above]{$y$};
      \draw[thick] (0,0) circle (0.95);
      \draw[->, thick] (0,0.95)
            arc[start angle=90, end angle=115, radius=0.95];
      \node[below=0.25cm] at (0,-1.3) {\footnotesize Circular};
    \end{scope}
    % Elíptica
    \begin{scope}[xshift=5.5cm]
      \draw[->] (-1.3,0)--(1.3,0) node[right]{$x$};
      \draw[->] (0,-1.3)--(0,1.3) node[above]{$y$};
      \draw[thick] (0,0) ellipse (0.95 and 0.55);
      \draw[->, thick] (0,0.55)
            arc[start angle=90, end angle=115, x radius=0.95, y radius=0.55];
      \node[below=0.25cm] at (0,-1.3) {\footnotesize Elíptica};
    \end{scope}
  \end{tikzpicture}
  \caption{Tipos de polarización: lineal, circular y elíptica. Las flechas
           indican el sentido de rotación dextrógiro.}
  \label{fig:polarizacion}
\end{figure}

\medskip
\begin{center}
\begin{tabularx}{0.95\linewidth}{@{}L{2.2cm}L{4.0cm}Y@{}}
  \toprule
  Tipo & Condición sobre $\delta$ & Condición sobre amplitudes \\
  \midrule
  Lineal    & $\delta = m\pi$              & cualquiera \\
  Circular  & $\delta = (2m+1)\pi/2$      & $E_{0x} = E_{0y}$ \\
  Elíptica  & $\delta = (2m+1)\pi/2$      & $E_{0x} \neq E_{0y}$ \\
  General   & $\delta$ arbitrario          & amplitudes arbitrarias \\
  \bottomrule
\end{tabularx}
\end{center}

% ══════════════════════════════════════════════════════════════════════════
\subsection{Intensidad de las Ondas Armónicas}
% ══════════════════════════════════════════════════════════════════════════

\begin{definicion}[Vector de Poynting]
El \emph{vector de Poynting} $\ve{S}$ describe la potencia electromagnética
transportada por unidad de área perpendicular a la dirección de propagación
(W/m$^2$) y su dirección de flujo. Su expresión en el vacío en términos de los
campos reales es:
\begin{equation}
  \ve{S} = \frac{1}{\mu_0}\,\ve{E}_r\times\ve{B}_r. \label{eq:poynting}
\end{equation}
donde el subíndice $r$ recuerda que deben usarse los campos reales, no las
amplitudes complejas. El vector de Poynting se obtiene del \emph{teorema de
Poynting}, que expresa la conservación local de energía electromagnética:
\[
  \pdv{u}{t} + \div\ve{S} = -\ve{j}\cdot\ve{E},
\]
donde $u = \frac{1}{2}(\varepsilon_0 E^2 + \frac{1}{\mu_0}B^2)$ es la densidad
de energía del campo y el miembro derecho es la potencia cedida por el campo
a las corrientes. En el vacío sin fuentes, $\ve{j}=\ve{0}$ y la energía se
conserva localmente.
\end{definicion}

\subsubsection{Intensidad de una onda armónica plana}

\begin{proposicion}
Para una OAP $\ve{E}=\ve{E}_0\cos(\omega t - \ve{k}\cdot\ve{r})$, la
intensidad (promedio temporal) es:
\begin{mdframed}[style=resultado]
\begin{equation}
  I = \langle S\rangle = \frac{E_0^2}{2\mu_0 c}. \label{eq:intensOAP}
\end{equation}
\end{mdframed}
donde $I$ es la intensidad (promedio temporal del módulo del vector de Poynting), $E_0$ la amplitud del campo eléctrico, $\mu_0$ la permeabilidad del vacío y $\langle\cdot\rangle$ denota el promedio temporal.
La energía se propaga en la dirección de $\ve{k}$ y es proporcional al
cuadrado de la amplitud del campo eléctrico.
\end{proposicion}

\begin{proof}
Partimos del campo real instantáneo, pues
\eqref{eq:poynting} es una magnitud cuadrática:
\[
  \ve{B}=\frac{1}{\omega}\,\ve{k}\times\ve{E}.
\]
Usando esta relación en el vector de Poynting:
\[
  \ve{S} = \frac{1}{\mu_0\omega}\,\ve{E}\times(\ve{k}\times\ve{E})
         = \frac{1}{\mu_0\omega}\!\left[
             \ve{k}E^2 - \ve{E}(\ve{k}\cdot\ve{E})
           \right]
         = \frac{E^2}{\mu_0\omega}\,\ve{k}.
\]
\noindent\normalfont\footnotesize(onda transversal: $\ve{k}\cdot\ve{E}=0$)\par\small\itshape
Para una OAP real:
\[
  E(t)=E_0\cos(\omega t-\ve{k}\cdot\ve{r})
  \;\Rightarrow\;
  E^2(t)=E_0^2\cos^2(\omega t-\ve{k}\cdot\ve{r}).
\]
Luego el módulo instantáneo del flujo es:
\[
  S(t)=\frac{k}{\mu_0\omega}E_0^2\cos^2(\omega t-\ve{k}\cdot\ve{r}).
\]
Al promediar en un período, $\langle\cos^2\rangle=1/2$:
\[
  \langle S\rangle = \frac{k}{\mu_0\omega}E_0^2\langle\cos^2\rangle
  = \frac{k}{\mu_0\omega}E_0^2\cdot\frac{1}{2}
  = \frac{E_0^2\,k}{2\mu_0\omega}
  = \frac{E_0^2}{2\mu_0 c}.
\]
\end{proof}

\subsubsection{Intensidad de una onda armónica no plana}

Para $\ve{E}=\ve{E}_0(\ve{r})\,e^{i\omega t}$, $\ve{B}=\ve{B}_0(\ve{r})\,e^{i\omega t}$,
tomando partes reales y promediando en el tiempo (los términos a $e^{\pm 2i\omega t}$
se anulan en el promedio):

\begin{mdframed}[style=resultado]
\begin{equation}
  \langle\ve{S}\rangle
  = \frac{1}{2\mu_0}\operatorname{Re}\!\left(\ve{E}_0\times\ve{B}_0^*\right).
  \label{eq:intensNP}
\end{equation}
\end{mdframed}

\begin{observacion}
La ecuación~\eqref{eq:intensNP} es el resultado general. Para una OAP,
$\ve{B}_0=(k/\omega)\,\uvec{k}\times\ve{E}_0$, y se recupera~\eqref{eq:intensOAP}.
\end{observacion}

% ══════════════════════════════════════════════════════════════════════════
\section{Tema 2: Interacción Radiación--Materia}
% ══════════════════════════════════════════════════════════════════════════

\subsection{Introducción}

El Tema 1 estableció que la luz es una onda electromagnética armónica plana
(OAP) que se propaga en el vacío satisfaciendo las ecuaciones de Maxwell. Sus
propiedades fundamentales —deducidas directamente de esas ecuaciones— son:

\begin{itemize}
  \item $\ve{B}\perp\ve{E}\perp\ve{k}$: onda estrictamente transversal.
  \item Relación de dispersión $k = \omega/c = 2\pi/\lambda$, con $\lambda\nu = c$.
  \item Relación entre amplitudes $E = cB$.
  \item Estado de polarización completamente caracterizado por el vector de Jones
        $\ve{J}=\bigl(E_{0x},\,E_{0y}\,e^{i\delta}\bigr)^T$.
\end{itemize}

Hasta aquí, el estudio ha sido puramente formal: hemos analizado la propagación
de la luz en el vacío sin preguntarnos cómo interacciona con la materia. Ahora
pasamos a estudiar la \textbf{interacción de la radiación electromagnética con
la materia} desde el punto de vista clásico microscópico (modelo de Lorentz,
1878). El mecanismo fundamental es:

\begin{enumerate}
  \item La onda incidente ejerce la fuerza de Lorentz sobre las cargas eléctricas
        del medio (electrones y núcleos).
  \item Estas cargas, al acelerarse de forma oscilatoria, actúan como dipolos
        oscilantes y emiten a su vez radiación electromagnética en todas las
        direcciones (\emph{scattering}).
  \item La competencia entre la frecuencia incidente $\omega$ y la frecuencia
        natural de oscilación del átomo $\omega_0$ determina la cantidad de
        energía intercambiada y los fenómenos observables (dispersión, absorción,
        índice de refracción).
\end{enumerate}

\begin{observacion}
En este tema trabajamos en el límite de átomos muy separados (medio
\emph{ópticamente diluido}): la onda interacciona con cada átomo de forma
independiente y los campos reemitidos por un átomo no afectan a sus vecinos.
Este régimen se supera en el Tema~3, donde los efectos colectivos generan
el índice de refracción del medio.
\end{observacion}

% ══════════════════════════════════════════════════════════════════════════
\subsection{Cargas Libres y Cargas Ligadas}
% ══════════════════════════════════════════════════════════════════════════

\begin{definicion}[Carga libre]
Una carga libre es aquella sin estructura interna relevante para el problema
(típicamente un electrón). Se caracteriza por su carga $q$, masa $m$ y
posición $\ve{r}(t)$. Puede moverse libremente bajo la acción de los campos.
\end{definicion}

\subsubsection{Aproximación para cargas libres}

La fuerza de Lorentz sobre una carga libre es:
\[
  \ve{F} = q\ve{E} + q\ve{v}\times\ve{B}.
\]
Usando $E = c B$ y estimando los órdenes de magnitud:
\[
  \left|\frac{q\ve{v}\times\ve{B}}{q\ve{E}}\right|
  = \frac{|\ve{v}||\ve{B}|\sin\alpha}{|\ve{E}|}
  = \frac{|\ve{v}|\sin\alpha}{c} \ll 1
  \quad \text{para } |\ve{v}| \ll c.
\]

\begin{proposicion}[Aproximación no relativista]
Para cargas cuya velocidad es no relativista ($|\ve{v}|\ll c$), el término
magnético es despreciable y la fuerza de Lorentz se reduce a:
\begin{equation}
  \ve{F} \approx q\ve{E}. \label{eq:fuerza_libre}
\end{equation}
\end{proposicion}

\begin{definicion}[Carga ligada]
Una carga ligada es aquella con estructura interna relevante (típicamente un
átomo). Se describe mediante la distancia $\ve{r}(t)$ entre los centros de
gravidad de las cargas positiva y negativa.

El campo eléctrico externo polariza el átomo, separando el núcleo del
electrón y generando un \emph{momento dipolar}:
\begin{equation}
  \mathbb{p} = q\,\ve{r}(t). \label{eq:dipolo}
\end{equation}
Una vez perturbado, el átomo tiende a volver a su posición de equilibrio
mediante una \textbf{fuerza recuperadora} proporcional al desplazamiento:
\begin{equation}
  \ve{F}_\text{rec} = -m\omega_0^2\,\ve{r}, \label{eq:frecup}
\end{equation}
análoga a la de un oscilador armónico con frecuencia natural $\omega_0$.
\end{definicion}

\begin{figure}[H]
  \centering
  \begin{tikzpicture}[>=Stealth, scale=1.2]
    % t=0: átomo en equilibrio
    \begin{scope}[xshift=-3.5cm]
      \node[circle, draw, thick, minimum size=14pt,
            inner sep=0pt] (n0) at (0,0) {$+$};
      \node[circle, draw, thick, fill=black!10,
            minimum size=14pt, inner sep=0pt] (e0) at (0,0) {};
    \end{scope}
    % Campo E aplicado
    \draw[->, very thick] (-1.0,-1.2) -- (1.0,-1.2)
          node[right] {$\ve{E}$};
    % t>0: átomo polarizado
    \begin{scope}[xshift=3.5cm]
      \node[circle, draw, thick, minimum size=14pt,
            inner sep=0pt] (np) at (-0.5,0) {$+$};
      \node[circle, draw, thick, fill=black!10,
            minimum size=14pt, inner sep=0pt] (ep) at (0.6,0) {$-$};
      \draw[->, thick] (-0.5,-0.6) -- (0.6,-0.6);
    \end{scope}
    \node[align=left, anchor=west] at (5.1,0.9)
      {\footnotesize Izquierda: equilibrio ($t=0$, $\ve{r}=0$)\\
       \footnotesize Derecha: átomo polarizado ($t>0$)\\
       \footnotesize Flecha inferior: $\mathbb{p}=q\ve{r}$};
  \end{tikzpicture}
  \caption{Polarización de un átomo bajo un campo eléctrico externo.
           En equilibrio ($t=0$) los centros de gravedad coinciden; bajo
           $\ve{E}$ se separan generando un dipolo $\mathbb{p}=q\ve{r}$.}
  \label{fig:dipolo}
\end{figure}

% ══════════════════════════════════════════════════════════════════════════
\subsection{Modelo Atómico de Lorentz}
% ══════════════════════════════════════════════════════════════════════════

El modelo de Lorentz formula la respuesta microscópica de un átomo como la de
un \textbf{oscilador armónico amortiguado y forzado}. Aunque su base es
clásica, reproduce con gran precisión fenomenológica la dispersión y la
absorción fuera de resonancias cuánticas finas, por lo que sigue siendo la
puerta de entrada estándar al comportamiento óptico de la materia.

\subsubsection{Hipótesis del modelo}

\begin{description}
  \item[\textbf{H1. Estado fundamental estable.}]
        Existe un estado natural en el que no se produce radiación. En
        ausencia de excitación, los centros de gravedad de las cargas
        positiva y negativa coinciden: $\ve{r}(t)=0$ y $\mathbb{p}=0$.

  \item[\textbf{H2. Fuerza recuperadora interna.}]
        Una fuerza lineal se opone a que el átomo sea desplazado de su
        estado fundamental:
        \[
          \ve{F}_\text{rec} = -m\omega_0^2\,\ve{r},
        \]
        donde $\omega_0$ es la \emph{frecuencia de resonancia natural} del
        átomo. Simboliza el carácter armónico del potencial interno.

  \item[\textbf{H3. Sistema abierto: amortiguamiento.}]
        Se introduce una fuerza de rozamiento que describe las pérdidas
        irreversibles de energía (radiación emitida u otros mecanismos):
        \[
          \ve{F}_\text{roz} = -m\gamma\,\dot{\ve{r}},
        \]
        con coeficiente de amortiguamiento:
        \begin{equation}
          \gamma = \frac{q^2\omega^2}{6\pi\varepsilon_0 m c^3}.
          \label{eq:gamma}
        \end{equation}
\end{description}

\begin{observacion}
De las ecuaciones de Maxwell sabemos que las cargas aceleradas emiten
ondas electromagnéticas, y estas transportan energía. El término de
amortiguamiento~\eqref{eq:gamma} es precisamente la tasa a la que el
sistema pierde esa energía.
\end{observacion}

\subsubsection{Ecuación de evolución}

Aplicando la segunda ley de Newton con las tres fuerzas descritas más la
fuerza de Lorentz del campo incidente (y utilizando la
aproximación~\eqref{eq:fuerza_libre}):

\begin{mdframed}[style=resultado]
\begin{equation}
  \ddot{\ve{r}} + \gamma\,\dot{\ve{r}} + \omega_0^2\,\ve{r}
  = \frac{q}{m}\ve{E}_0\,e^{i\omega t}.
  \label{eq:lorentz_atom}
\end{equation}
\end{mdframed}

Esta es la ecuación de un \textbf{oscilador armónico amortiguado y forzado},
en la que cada término tiene su interpretación física:
\begin{itemize}
  \item $+\omega_0^2\ve{r}$: término elástico (fuerza recuperadora, de signo
        opuesto al desplazamiento).
  \item $+\gamma\,\dot{\ve{r}}$: término disipativo (amortiguamiento por
        radiación emitida u otros mecanismos de pérdida).
  \item $(q/m)\ve{E}_0\,e^{i\omega t}$: forzamiento externo oscilante (campo
        electromagnético incidente).
\end{itemize}
La solución general es suma de la solución homogénea (transitoria, que decae
con constante de tiempo $\tau \sim 1/\gamma \sim 10^{-8}\ \mathrm{s}$) y la
solución particular estacionaria. En óptica, el transitorio es irrelevante y
se trabaja siempre con el régimen estacionario.

\begin{observacion}
Para \emph{cargas libres} ($\omega_0=0$, como los electrones de conducción en un metal), la ecuación~\eqref{eq:lorentz_atom} se reduce al oscilador forzado sin término elástico, con $\ve{r}$ como posición absoluta de la carga libre.
Para \emph{cargas ligadas} (electrones de valencia), $\ve{r}$ es el desplazamiento relativo entre el centro de gravedad del cortejo electrónico y el núcleo, y $\omega_0$ es la frecuencia de resonancia atómica.
\end{observacion}

\subsubsection{Aproximaciones para resolver la ecuación}

La ecuación~\eqref{eq:lorentz_atom} se simplifica con dos aproximaciones:

\begin{description}
  \item[\textbf{Ap. 1: Aproximación dipolar.}]
        Se ignora la dependencia espacial del campo eléctrico en la posición
        del electrón. Esto está justificado porque:
        \[
          k\cdot r \sim \frac{2\pi}{\lambda}\cdot r
          \approx \frac{1\,\text{\AA}}{400\,\text{nm}} \approx 2.5\times10^{-4} \ll 1,
        \]
        por lo que $e^{-i\ve{k}\cdot\ve{r}}\approx 1$ y el campo actúa
        uniformemente sobre el átomo: $\ve{E}(\ve{r},t)\approx\ve{E}_0 e^{i\omega t}$.

  \item[\textbf{Ap. 2: Despreciar el campo magnético.}]
        Usando $B=E/c$ y la aproximación no relativista, el término magnético
        de la fuerza de Lorentz es despreciable frente al eléctrico
        (ecuación~\eqref{eq:fuerza_libre}).
\end{description}

\subsubsection{Solución de la ecuación de evolución}

\begin{proposicion}[Solución estacionaria]
La solución particular de~\eqref{eq:lorentz_atom} en régimen estacionario es:
\begin{equation}
  \ve{r}(t) = \frac{q/m}{\omega_0^2 - \omega^2 + i\gamma\omega}
              \,\ve{E}_0\,e^{i\omega t}. \label{eq:sol_lorentz}
\end{equation}
\end{proposicion}

\begin{proof}
Proponemos la solución armónica:
\[
  \ve{r}(t) = \ve{r}_0\,e^{i\omega t}.
\]
Entonces:
\[
  \dot{\ve{r}}=i\omega\ve{r}_0e^{i\omega t},
  \qquad
  \ddot{\ve{r}}=-\omega^2\ve{r}_0e^{i\omega t}.
\]
Sustituyendo en~\eqref{eq:lorentz_atom}:
\[
  \bigl(-\omega^2 + i\gamma\omega + \omega_0^2\bigr)\ve{r}_0\,e^{i\omega t}
  = \frac{q}{m}\ve{E}_0\,e^{i\omega t}.
\]
Simplificando $e^{i\omega t}$ y despejando $\ve{r}_0$:
\[
  \ve{r}_0 = \frac{q/m}{\omega_0^2 - \omega^2 + i\gamma\omega}\,\ve{E}_0.
\]
\end{proof}

\subsubsection{Análisis de la solución: desfase}

Separando $r_0$ en partes real e imaginaria (con $E_0$ real):
\[
  r_0 = \frac{q}{m}\,E_0\left[
    \frac{\omega_0^2-\omega^2}{(\omega_0^2-\omega^2)^2+\gamma^2\omega^2}
    - i\frac{\gamma\omega}{(\omega_0^2-\omega^2)^2+\gamma^2\omega^2}
  \right] = |r_0|\,e^{-i\alpha},
\]
con el ángulo de desfase:
\begin{equation}
  \tan\alpha = \frac{\gamma\omega}{\omega_0^2-\omega^2}. \label{eq:desfase}
\end{equation}
El módulo de la amplitud de oscilación es:
\begin{equation}
  |r_0| = \frac{(q/m)\,E_0}
  {\sqrt{(\omega_0^2-\omega^2)^2+\gamma^2\omega^2}},
  \qquad \mathbb{p}_0 = q\,r_0.
  \label{eq:modulo_lorentz}
\end{equation}

La ecuación \eqref{eq:desfase} cuantifica el retraso entre excitación y
respuesta del dipolo: por debajo de resonancia domina la parte elástica (fase
pequeña), en resonancia domina la disipación (cuadratura) y por encima domina
la inercia (fase cercana a $\pi$).

\medskip
\begin{center}
\begin{tabularx}{0.98\linewidth}{@{}L{2.5cm}L{2.8cm}L{1.4cm}Y@{}}
  \toprule
  Régimen & Condición & $\alpha$ & Interpretación \\
  \midrule
  En fase       & $\omega \ll \omega_0$ & $0$      & Campo y movimiento en fase; el dipolo sigue adiabáticamente al campo \\
  Resonancia    & $\omega = \omega_0$   & $\pi/2$  & Velocidad en fase con el campo; transferencia máxima de energía \\
  Oposición     & $\omega \gg \omega_0$ & $\pi$    & Inercia domina; movimiento opuesto al campo \\
  Carga libre   & $\omega_0 = 0$        & $\pi$    & Igual que el régimen de oposición \\
  \bottomrule
\end{tabularx}
\end{center}

\begin{observacion}
La condición de transferencia máxima de energía a la resonancia ($\alpha = \pi/2$)
tiene un paralelo mecánico exacto: un columpio gana energía máximamente cuando se
empuja cuando está en el extremo de su recorrido (velocidad cero), es decir, en
cuadratura de fase con la excitación. Este criterio —velocidad en fase con la
fuerza— es universal para todos los osciladores forzados, independientemente del
mecanismo de pérdida. La anchura espectral de la resonancia es $\Delta\omega \approx
\gamma$, por lo que la calidad del resonador es $Q = \omega_0/\gamma$.
\end{observacion}

% ══════════════════════════════════════════════════════════════════════════
\subsection{Potencia Transferida a la Carga}
% ══════════════════════════════════════════════════════════════════════════

\begin{definicion}[Potencia transferida]
La potencia instantánea que el campo electromagnético transfiere a la carga
es el trabajo que la fuerza de Lorentz realiza por unidad de tiempo:
\[
  W_\text{inst} = \ve{F}\cdot\dot{\ve{r}} \approx q\ve{E}\cdot\dot{\ve{r}}.
\]
Al ser una cantidad oscilatoria, se toma el promedio temporal, análogamente
a como se hace con el vector de Poynting:
\[
  \langle\ve{S}\rangle = \frac{1}{2\mu_0}\operatorname{Re}(\ve{E}_0\times\ve{B}_0^*).
\]
\end{definicion}

\begin{proposicion}[Potencia media transferida]
La potencia media transferida a la carga es:
\begin{mdframed}[style=resultado]
\begin{equation}
  \langle W\rangle = \frac{q^2}{2m}
  \frac{\gamma\omega^2}{\bigl(\omega_0^2-\omega^2\bigr)^2+\gamma^2\omega^2}
  \,E_0^2. \label{eq:potencia_transferida}
\end{equation}
\end{mdframed}
\end{proposicion}

\begin{proof}
Usando la regla del promedio para productos de campos armónicos:
\[
  \langle W\rangle = \frac{1}{2}\operatorname{Re}\{q\ve{E}\cdot\dot{\ve{r}}^*\}.
\]
Tomamos:
\[
  \ve{E} = \ve{E}_0 e^{i\omega t},
  \qquad
  \dot{\ve{r}}^* = (-i\omega)\ve{r}_0^* e^{-i\omega t}.
\]
\[
  \langle W\rangle = \frac{1}{2}\operatorname{Re}\{q E_0\,e^{i\omega t}
  \cdot(-i\omega) r_0^* e^{-i\omega t}\}
  = \frac{q\omega E_0}{2}\operatorname{Im}\{r_0^*\}.
\]
Tomando la parte imaginaria de $r_0$:
\[
  r_0=\frac{q/m}{\omega_0^2-\omega^2+i\gamma\omega}E_0
  =\frac{q}{m}\frac{\omega_0^2-\omega^2-i\gamma\omega}
    {(\omega_0^2-\omega^2)^2+\gamma^2\omega^2}E_0,
\]
\[
  \operatorname{Im}\{r_0^*\}
  = \frac{q/m\cdot\gamma\omega}{(\omega_0^2-\omega^2)^2+\gamma^2\omega^2}E_0.
\]
Sustituyendo se obtiene~\eqref{eq:potencia_transferida}.
\end{proof}

\begin{observacion}
La potencia transferida tiene carácter dispersivo (de tipo Lorentziano) en
$\omega$: es máxima cuando $\omega\approx\omega_0$ (resonancia) y decrece
rápidamente fuera de ella con una anchura $\Delta\omega \approx \gamma$. Esta
absorción selectiva en frecuencia es la base física de las líneas espectrales
de absorción. Macroscópicamente, en un medio con $N$ átomos por unidad de
volumen, la potencia absorbida conduce a la \emph{ley de Beer-Lambert}: la
intensidad decae exponencialmente con la distancia recorrida en el medio,
$I(z) = I_0\,e^{-\alpha_\text{abs}\,z}$, donde el coeficiente de absorción
$\alpha_\text{abs} = 2k_0\kappa$ es proporcional a la parte imaginaria del
índice complejo.
\end{observacion}

% ══════════════════════════════════════════════════════════════════════════
\subsection{Potencia Esparcida por la Carga}
% ══════════════════════════════════════════════════════════════════════════

\subsubsection{Tipos de medios ópticos}

\begin{definicion}[Medio ópticamente diluido]
Un medio es \emph{ópticamente diluido} cuando el número de partículas por
volumen de longitud de onda es pequeño:
\[
  \frac{N_\text{átomos}}{\lambda^3} \leq 1.
\]
Las cargas están muy separadas entre sí y no interactúan mutuamente. Cada
átomo puede tratarse de forma independiente.
\end{definicion}

\begin{definicion}[Medio ópticamente denso]
Un medio es \emph{ópticamente denso} cuando:
\[
  \frac{N_\text{átomos}}{\lambda^3} > 1.
\]
Las interacciones entre cargas producen efectos constructivos o destructivos
(interferencias, dispersión, absorción).
\end{definicion}

\subsubsection{Esparcimiento (\emph{scattering})}

En un medio diluido, la potencia reemitida por la carga es igual a la
potencia que el campo le transfiere. Sustituyendo $\gamma$ de~\eqref{eq:gamma}
en~\eqref{eq:potencia_transferida}:

\begin{mdframed}[style=resultado]
\begin{equation}
  P_\text{esp} = A\,
  \frac{\omega^4}{\bigl(\omega_0^2-\omega^2\bigr)^2+\gamma^2\omega^2},
  \qquad
  A = \frac{q^4 E_0^2}{12\pi\varepsilon_0 m^2 c^3} = \mathrm{cte}.
  \label{eq:P_esp}
\end{equation}
\end{mdframed}

Este fenómeno —la reemisión de radiación en todas las direcciones por una
carga excitada— se denomina \textbf{esparcimiento} o \emph{scattering}.

\subsubsection{Casos particulares}

Analizamos tres regímenes de frecuencia:

\paragraph{Frecuencia muy alta: $\omega\gg\omega_0$}

\[
  \bigl(\omega_0^2-\omega^2\bigr)^2+\gamma^2\omega^2 \approx \omega^4,
\]
\begin{equation}
  P_\text{esp} \approx A = \mathrm{cte}. \label{eq:P_alta_freq}
\end{equation}
La potencia reemitida es constante e independiente de $\omega$. La carga
responde como libre ($\omega_0\to 0$).

\paragraph{Resonancia: $\omega=\omega_0$}

\[
  \bigl(\omega_0^2-\omega^2\bigr)^2+\gamma^2\omega^2 \approx \gamma^2\omega_0^2,
\]
\begin{equation}
  P_\text{esp} \approx A\,\frac{\omega_0^2}{\gamma^2}. \label{eq:P_resonancia}
\end{equation}
La potencia es muy grande (máximo), ya que $\gamma\ll\omega_0$ en general.
Su valor depende de la frecuencia propia del átomo.

\paragraph{Frecuencia muy baja: $\omega\ll\omega_0$ — Esparcimiento de Rayleigh}

\[
  \bigl(\omega_0^2-\omega^2\bigr)^2+\gamma^2\omega^2 \approx \omega_0^4,
\]
\begin{mdframed}[style=resultado]
\begin{equation}
  P_\text{esp} \approx A\,\frac{\omega^4}{\omega_0^4}
  \propto \frac{1}{\lambda^4}. \label{eq:rayleigh}
\end{equation}
\end{mdframed}
La potencia esparcida crece con la cuarta potencia de la frecuencia
(o inversamente con $\lambda^4$). Este es el \textbf{esparcimiento de
Rayleigh}.

\medskip
\begin{center}
\begin{tabularx}{0.98\linewidth}{@{}L{2.5cm}L{2.8cm}Y L{2.7cm}@{}}
  \toprule
  Régimen & Condición & $P_\text{esp}$ & Nombre \\
  \midrule
  Alta frecuencia & $\omega\gg\omega_0$ & $\approx A$              & -- \\
  Resonancia      & $\omega=\omega_0$   & $\approx A\omega_0^2/\gamma^2$ & Resonancia óptica \\
  Baja frecuencia & $\omega\ll\omega_0$ & $\propto\omega^4/\omega_0^4$   & Rayleigh \\
  \bottomrule
\end{tabularx}
\end{center}

% ══════════════════════════════════════════════════════════════════════════
\subsection{Esparcimiento del Cielo — Color del Cielo}
% ══════════════════════════════════════════════════════════════════════════

\subsubsection{Esparcimiento de Rayleigh en la atmósfera}

La atmósfera está compuesta principalmente por moléculas de $\mathrm{N_2}$,
$\mathrm{O_2}$, $\mathrm{H_2O}$, $\mathrm{CO_2}$, cuya frecuencia de
resonancia se encuentra en el ultravioleta:
\[
  \omega_0^\text{atm} \sim 3\times10^{15}\ \text{rad/s}
  \quad (\lambda_0\approx 300\ \text{nm}).
\]
Como la luz visible cumple $\omega_\text{vis}\ll\omega_0^\text{atm}$,
se aplica el régimen de Rayleigh~\eqref{eq:rayleigh}:
\[
  P_\text{esp} \propto \omega^4 \propto \frac{1}{\lambda^4}.
\]
A menor longitud de onda, mayor esparcimiento:

\begin{center}
\begin{tabularx}{0.9\linewidth}{@{}L{2.8cm}L{2.4cm}Y@{}}
  \toprule
  Color & $\lambda$ (nm) & $\omega$ relativo \\
  \midrule
  Violeta/Azul & 400--490 & mayor $\omega$ → mayor esparcimiento \\
  Verde        & 555       &  \\
  Amarillo     & 589       &  \\
  Rojo         & 630--750 & menor $\omega$ → menor esparcimiento \\
  \bottomrule
\end{tabularx}
\end{center}

\subsubsection{Los tres casos del cielo}

La dependencia $P_\text{esp} \propto 1/\lambda^4$ del esparcimiento de Rayleigh
implica que la luz de corta longitud de onda (violeta, azul) se esparce con
mucha más eficacia que la de larga longitud de onda (rojo, naranja). La
relación de intensidades esparcidas entre el azul ($\lambda\approx450$ nm) y
el rojo ($\lambda\approx650$ nm) es aproximadamente $(650/450)^4 \approx 4.4$:
el azul se esparce más de cuatro veces más que el rojo.

\begin{description}
  \item[\textbf{Caso 1: Observador mirando al Sol directamente.}]
        La luz solar (blanca) recorre un camino corto dentro de la atmósfera.
        El azul y el violeta son preferentemente esparcidos lateralmente,
        de modo que el haz que llega directamente al observador se empobrece
        en esas longitudes de onda y resulta enriquecido en amarillo y naranja:
        el Sol parece \textbf{amarillo-blanco}.

  \item[\textbf{Caso 2: Observador mirando en una dirección oblicua.}]
        La luz llega al observador \emph{indirectamente}, tras haber sido
        esparcida por las moléculas de la atmósfera en múltiples eventos de
        Rayleigh. Los fotones azules y violetas son los más esparcidos y los
        que con mayor probabilidad alcanzan al observador desde cualquier
        dirección lateral. Por eso el cielo aparece \textbf{azul}. El violeta
        es aún más esparcido que el azul, pero la sensibilidad del ojo humano
        es menor en ese extremo del espectro.

  \item[\textbf{Caso 3: Atardecer (luz rasante).}]
        En el horizonte, la luz solar debe atravesar un espesor de atmósfera
        unas 40 veces mayor que al mediodía. Tras este largo trayecto, los
        fotones azules y violetas han sido casi totalmente esparcidos y han
        abandonado el haz; solo los colores de baja frecuencia (rojo, naranja,
        amarillo) sobreviven el viaje. El cielo y el Sol aparecen
        \textbf{naranja-rojizo}.
\end{description}

\begin{figure}[H]
  \centering
  \begin{tikzpicture}[>=Stealth, scale=1.1]
    % Tierra
    \draw[thick] (-4.5, -0.5) -- (4.5, -0.5);

    % Atmósfera (banda)
    \draw[thick, dashed] (-4.5, 1.2) -- (4.5, 1.2);

    % Sol (derecha, arriba)
    \filldraw[black] (4.0, 2.5) circle (5pt);

    % Caso 1: rayo directo
    \draw[->, thick] (4.0, 2.5) -- (0.8, -0.5);

    % Caso 2: rayo esparcido
    \draw[->, thick] (4.0, 2.5) -- (0.0, 0.9);
    \draw[->, thick, dashed] (0.0, 0.9) -- (-1.5, -0.5);
    \filldraw (0.0, 0.9) circle (2pt);

    % Caso 3: rayo rasante (atardecer)
    \draw[->, thick] (-4.0, 2.0) -- (-2.5, 0.7);
    \filldraw (-2.5, 0.7) circle (2pt);
    \draw[->, thick] (-2.5, 0.7) -- (-1.0, -0.5);

    % Observadores
    \filldraw (0.8,-0.5) circle (2pt) node[below=3pt] {Obs. 1};
    \filldraw (-1.5,-0.5) circle (2pt) node[below=3pt] {Obs. 2};
    \filldraw (-1.0,-0.5) circle (2pt) node[below=3pt] {Obs. 3};
    \node[align=left, anchor=west] at (5.0,1.1)
      {\footnotesize Línea base: superficie terrestre\\
       \footnotesize Línea discontinua: atmósfera\\
       \footnotesize Caso 1 (Obs. 1): luz directa\\
       \footnotesize Caso 2 (Obs. 2): luz azul esparcida\\
       \footnotesize Caso 3 (Obs. 3): trayectoria rojiza};
  \end{tikzpicture}
  \caption{Esquema del esparcimiento de Rayleigh para los tres casos:
           visión directa del Sol (Obs.~1), cielo lateral (Obs.~2)
           y atardecer (Obs.~3).}
  \label{fig:rayleigh_cielo}
\end{figure}

% ══════════════════════════════════════════════════════════════════════════
\subsection{Polarización de la Luz Esparcida}
% ══════════════════════════════════════════════════════════════════════════

\begin{definicion}[Dipolo oscilante]
Cuando un campo eléctrico linealmente polarizado incide sobre una carga, la
excita a oscilar a lo largo de la dirección de $\ve{E}$. El electrón
acelerado actúa como un \emph{dipolo oscilante}, emitiendo radiación en
todas las direcciones salvo a lo largo del propio eje de oscilación.
\end{definicion}

\subsubsection{Distribución angular de la emisión}

La potencia emitida por un dipolo oscilante en función del ángulo $\theta$
(medido respecto al eje de oscilación) es:
\begin{equation}
  \frac{dP}{d\Omega} \propto \sin^2\theta. \label{eq:dipolo_angular}
\end{equation}
donde $\Omega$ es el ángulo sólido de observación y $\theta$ el ángulo medido respecto al eje de oscilación del dipolo.

Consecuencias:
\begin{itemize}
  \item En la dirección $\theta=0^{\circ}$ (eje de oscilación): $dP/d\Omega=0$.
        No hay emisión.
  \item En la dirección $\theta=90^{\circ}$ (perpendicular al eje): emisión
        máxima.
\end{itemize}

\subsubsection{Polarización del cielo}

Cuando luz solar no polarizada incide sobre las moléculas de la atmósfera,
cada molécula responde en el plano de $\ve{E}$. Para un observador que mire
en una dirección perpendicular al Sol:

\begin{proposicion}[Polarización de la luz esparcida]
La luz esparcida en dirección perpendicular al rayo solar está
\textbf{completamente polarizada linealmente}, con $\ve{E}$ perpendicular
al plano que contiene el Sol, la molécula y el observador.
\end{proposicion}

\begin{proof}
Descomponemos la luz no polarizada en dos componentes lineales ortogonales:
\[
  \ve{E}_\parallel,
  \qquad
  \ve{E}_\perp.
\]
La componente $\ve{E}_\parallel$ induce un dipolo oscilando en el plano
Sol--molécula--observador.
Para esa componente, la dirección hacia el observador coincide con el eje de
oscilación del dipolo, de modo que:
\[
  \theta=0^{\circ}
  \quad\Rightarrow\quad
  \left.\frac{dP}{d\Omega}\right|_{\parallel}\propto \sin^2\theta=0.
\]

La componente $\ve{E}_\perp$ induce un dipolo perpendicular a ese plano.
Hacia el observador, en ese caso:
\[
  \theta=90^{\circ},
\]
por tanto:
\[
  \left.\frac{dP}{d\Omega}\right|_{\perp}\propto \sin^2\theta=1.
\]
Así, la radiación observada procede solo de $\ve{E}_\perp$.
Por tanto, la luz observada queda polarizada en la dirección $\ve{E}_\perp$.
\end{proof}

\begin{figure}[H]
  \centering
  \begin{tikzpicture}[>=Stealth, scale=1.2]
    % Sol
    \filldraw (0,0) circle (4pt);

    % Molécula
    \filldraw (3.5,0) circle (3pt);

    % Observador (perpendicular)
    \draw[->, thick] (3.5, 0) -- (3.5, -1.8);

    % Rayo incidente
    \draw[->, thick] (0,0) -- (3.2,0);

    % E_perp (perpendicular al plano): doble flecha en la molécula
    \draw[<->, very thick] (3.5, 0.5) -- (3.5, -0.5);

    % E_parallel: no llega al obs.
    \draw[<->, thick, dashed] (3.0, 0) -- (4.0, 0);
    \node[align=left, anchor=west] at (5.0,1.0)
      {\footnotesize Punto izquierdo: Sol\\
       \footnotesize Punto central: molécula\\
       \footnotesize Flecha vertical: observador\\
       \footnotesize Segmento sólido: componente $\ve{E}_\perp$\\
       \footnotesize Segmento discontinuo: $\ve{E}_\parallel\to 0$};
  \end{tikzpicture}
  \caption{Polarización de la luz esparcida a $90^{\circ}$ del Sol.
           Solo la componente $\ve{E}_\perp$ llega al observador.}
  \label{fig:pol_esparcida}
\end{figure}

\begin{observacion}
Este efecto es observable experimentalmente: el cielo lateral al Sol
(a $90^{\circ}$) está parcialmente polarizado y puede detectarse con un
filtro polarizador. Las abejas y otros insectos usan esta polarización
para orientarse.
\end{observacion}

% ══════════════════════════════════════════════════════════════════════════
\section{Tema 3: Propagación de Ondas en Medios Materiales}
% ══════════════════════════════════════════════════════════════════════════

\subsection{Introducción}

Los temas anteriores estudiaron la propagación de la luz en el vacío y su
interacción microscópica con cargas aisladas. Ahora afrontamos la pregunta
físicamente más rica: \textbf{¿cómo se propaga la luz dentro de un medio
material?}

La clave conceptual está en que un medio material contiene un número enorme de
átomos por unidad de volumen. Cada átomo responde al campo incidente tal como
se estudió en el Tema~2, pero en un medio denso los campos reemitidos por todos
los átomos se superponen coherentemente, modificando el campo total. Este
efecto colectivo da lugar a fenómenos nuevos: la onda se propaga con una
velocidad diferente de $c$ (índice de refracción real $\eta$), puede
atenuarse exponencialmente (absorción, cuantificada por $\kappa$), y su
longitud de onda efectiva en el medio es $\lambda_\text{med}=\lambda_0/\eta$.

Los objetivos del tema son dos:
\begin{enumerate}
  \item Resolver las ecuaciones de Maxwell macroscópicas (promediadas
        espacialmente) en la materia, para obtener una \emph{permitividad
        dieléctrica generalizada} $\varepsilon_\text{gen}$ que incorpore
        simultáneamente los efectos de cargas libres (conductividad) y de
        cargas ligadas (polarización dieléctrica).
  \item Obtener el \emph{índice de refracción complejo} $n_c = \eta - i\kappa$,
        que determina completamente la propagación de la luz en el medio:
        velocidad de fase, atenuación (longitud de penetración), longitud de
        onda efectiva y reflectividad de las interfaces.
\end{enumerate}

\begin{observacion}
Las ecuaciones de Maxwell siguen siendo la base exacta, pero ahora las fuentes
$\rho$ y $\ve{j}$ incluyen contribuciones de cargas libres (corrientes de
conducción) y de cargas ligadas (polarización $\ve{P}$ y magnetización $\ve{M}$
del medio). El paso clave del Tema~3 es el tránsito de la descripción
microscópica a la descripción macroscópica mediante un promedio espacial
sobre una escala intermedia $a \ll \Delta V \ll \lambda^3$.
\end{observacion}

% ══════════════════════════════════════════════════════════════════════════
\subsection{Medios Ópticamente Densos}
% ══════════════════════════════════════════════════════════════════════════

\begin{definicion}[Medio ópticamente diluido]
Un medio es ópticamente diluido cuando el número de átomos por volumen de
longitud de onda es pequeño:
\[
  \frac{N_\text{átomos}}{\lambda^3} \leq 1.
\]
Las cargas están muy separadas y no interactúan entre sí. Cada átomo puede
tratarse de forma independiente (Tema~2).
\end{definicion}

\begin{definicion}[Medio ópticamente denso]
Un medio es ópticamente denso cuando:
\[
  \frac{N_\text{átomos}}{\lambda^3} > 1.
\]
Las interacciones entre cargas producen efectos constructivos o destructivos.
El campo total sobre cada átomo es la suma del campo incidente y los campos
reemitidos por todos los demás:
\begin{equation}
  \ve{E}_\text{tot} = \ve{E}_\text{inc} + \sum_{i=1}^{N} \ve{E}_i.
  \label{eq:campo_total}
\end{equation}
\end{definicion}

\begin{ejemplo}[Densidad óptica según $\lambda$]
Consideramos un sólido con densidad $10^{22}\ \text{átomos/cm}^3$.

\begin{itemize}
  \item Para $\lambda = 600\ \text{nm}$:
        $N/\lambda^3 = 10^{22}/(600\times10^{-7})^3\ \text{cm}^{-3}
        \approx 4.6\times10^{3} \gg 1$. \textbf{Medio denso.}
  \item Para $\lambda = 0{,}1\ \text{nm}$ (rayos X):
        $N/\lambda^3 = 10^{22}/(10^{-8})^3 = 10^{22}/10^{-24}
        = 10^{46}\ \text{cm}^{-3}$.
        Todos los átomos interactúan coherentemente.
\end{itemize}
\end{ejemplo}

% ══════════════════════════════════════════════════════════════════════════
\subsection{Ecuaciones de Maxwell en la Materia: Ecuaciones Macroscópicas}
% ══════════════════════════════════════════════════════════════════════════

\subsubsection{Del nivel microscópico al macroscópico}

Las ecuaciones microscópicas son exactas, pero su uso directo en un sólido
requiere seguir un número astronómico de cargas. El paso macroscópico consiste
en \textbf{promediar espacialmente} sobre un volumen $\Delta V$ tal que:
\[
  a^3 \ll \Delta V \ll \lambda^3,
\]
donde $a \sim 10^{-10}\ \mathrm{m}$ es el tamaño atómico. Este volumen
contiene muchos átomos (promedio estadístico bien definido) pero es mucho
menor que el cubo de la longitud de onda (resolución óptica preservada).

Un resultado clave es que, bajo estas hipótesis de escala, promedio y
derivación espacial conmutan:
\begin{equation}
  \left\langle \pdv{\ve{E}}{x}\right\rangle
  = \pdv{\langle\ve{E}\rangle}{x}. \label{eq:conmuta}
\end{equation}
Por tanto, las ecuaciones de Maxwell promediadas tienen la misma estructura
formal que las microscópicas, pero con $\rho$ y $\ve{j}$ reemplazados por
sus promedios macroscópicos. El resultado son las ecuaciones
macroscópicas~\eqref{eq:maxmat1}--\eqref{eq:maxmat4}, que se obtienen una
vez se descomponen esos promedios.

\begin{figure}[H]
  \centering
  \begin{tikzpicture}[>=Stealth, scale=1.0]
    \draw[rounded corners, fill=gray!10] (0,0) rectangle (2.8,1.2);
    \node at (1.4,0.8) {$a^3$};
    \node[below=3pt] at (1.4,0) {\footnotesize escala atómica};

    \draw[rounded corners, fill=gray!6] (4.2,0) rectangle (8.0,1.2);
    \node at (6.1,0.8) {$\Delta V$};
    \node[below=3pt] at (6.1,0) {\footnotesize volumen de promedio};

    \draw[rounded corners, fill=gray!3] (9.4,0) rectangle (13.2,1.2);
    \node at (11.3,0.8) {$\lambda^3$};
    \node[below=3pt] at (11.3,0) {\footnotesize escala óptica};

    \draw[->, thick] (2.9,0.6) -- (4.1,0.6) node[midway, above=3pt] {$\ll$};
    \draw[->, thick] (8.1,0.6) -- (9.3,0.6) node[midway, above=3pt] {$\ll$};
  \end{tikzpicture}
  \caption{Jerarquía de escalas usada para pasar de la descripción
           microscópica a la macroscópica: $a^3 \ll \Delta V \ll \lambda^3$.}
  \label{fig:escalas_promedio}
\end{figure}

\subsubsection{Descomposición de cargas y corrientes}

\begin{proposicion}[Densidad de carga promediada]
La densidad de carga total se descompone en cargas libres y ligadas:
\begin{equation}
  \langle\rho\rangle = \rho_\text{libre} - \div\ve{P}, \label{eq:rho_macro}
\end{equation}
donde $\ve{P}$ es la \emph{polarización eléctrica} (momento dipolar por
unidad de volumen):
\begin{equation}
  \ve{P} = \frac{1}{\Delta V}\sum_{j\in\Delta V} q_j\,\ve{r}_j.
  \label{eq:polarizacion}
\end{equation}
\end{proposicion}

\begin{proposicion}[Densidad de corriente promediada]
La densidad de corriente total se descompone en:
\begin{equation}
  \langle\ve{j}\rangle = \ve{j}_\text{libre}
  + \curl\ve{M} + \pdv{\ve{P}}{t}, \label{eq:j_macro}
\end{equation}
donde $\ve{M}$ es la \emph{magnetización} (momento dipolar magnético por
unidad de volumen).
\end{proposicion}

\subsubsection{Ecuaciones de Maxwell macroscópicas}

Sustituyendo~\eqref{eq:rho_macro} y~\eqref{eq:j_macro} en las ecuaciones
microscópicas y definiendo los vectores auxiliares:
\begin{equation}
  \ve{D} = \varepsilon_0\ve{E} + \ve{P}, \qquad
  \ve{H} = \frac{1}{\mu_0}\ve{B} - \ve{M}, \label{eq:DH}
\end{equation}
donde $\ve{D}$ es el desplazamiento eléctrico, $\ve{H}$ el campo magnético auxiliar, $\ve{P}$ la polarización eléctrica y $\ve{M}$ la magnetización del medio.
se obtienen las \textbf{ecuaciones de Maxwell macroscópicas}:

\begin{mdframed}[style=resultado]
\begin{align}
  \div\ve{D} &= \rho_\text{libre}, \label{eq:maxmat1}\\[4pt]
  \div\ve{B} &= 0, \label{eq:maxmat2}\\[4pt]
  \curl\ve{E} + \pdv{\ve{B}}{t} &= \ve{0}, \label{eq:maxmat3}\\[4pt]
  \curl\ve{H} &= \ve{j}_\text{libre} + \pdv{\ve{D}}{t}. \label{eq:maxmat4}
\end{align}
\end{mdframed}

% ══════════════════════════════════════════════════════════════════════════
\subsection{Relaciones de Constitución y Permitividad Generalizada}
% ══════════════════════════════════════════════════════════════════════════

Las ecuaciones macroscópicas introducen campos auxiliares ($\ve{D},\ve{H}$),
pero no cierran el problema sin ecuaciones de material. Las
\emph{relaciones de constitución} cumplen precisamente esa función: vinculan
campo aplicado y respuesta interna del medio.

\subsubsection{Primera relación: permitividad eléctrica}

La respuesta dieléctrica del medio es lineal con el campo:
$\ve{P} = \varepsilon_0\chi_e\ve{E}$, donde $\chi_e$ es la susceptibilidad
eléctrica. Entonces:
\begin{equation}
  \ve{D} = \varepsilon_0(1+\chi_e)\ve{E}
         = \varepsilon_0\varepsilon_r\ve{E}
         = \varepsilon\ve{E}, \qquad
  \varepsilon_r = 1+\chi_e. \label{eq:const1}
\end{equation}
donde $\varepsilon_r$ es la permitividad relativa del medio.

\subsubsection{Segunda relación: permeabilidad magnética}

Análogamente, $\ve{M} = \chi_m\ve{H}$:
\begin{equation}
  \ve{B} = \mu_0(1+\chi_m)\ve{H}
         = \mu_0\mu_r\ve{H}
         = \mu\ve{H}, \qquad
  \mu_r = 1+\chi_m. \label{eq:const2}
\end{equation}
donde $\mu_r$ es la permeabilidad relativa del medio.

\subsubsection{Tercera relación: conductividad (ley de Ohm)}

Para cargas libres, la ley de Ohm establece:
\begin{equation}
  \ve{j}_\text{libre}(\ve{r},t) = \sigma\ve{E}(\ve{r},t). \label{eq:ohm}
\end{equation}
donde $\sigma$ es la conductividad eléctrica del medio.
Aplicando conservación de carga y pasando a régimen armónico
($\partial/\partial t\to i\omega$):
\[
  \pdv{\rho_\text{libre}}{t} + \div\ve{j}_\text{libre} = 0
  \implies \rho_\text{libre} = \div\!\left(\frac{i\sigma}{\omega}\ve{E}\right).
\]

\subsubsection{Permitividad dieléctrica generalizada}

Sustituyendo las tres relaciones de constitución en la
ecuación~\eqref{eq:maxmat4} y trabajando con ondas armónicas, el término
de corriente libre y el de desplazamiento se pueden combinar:
\[
  \curl\ve{H} = \ve{j}_\text{libre} + \pdv{(\varepsilon\ve{E})}{t}
  = \sigma\ve{E} + i\omega\varepsilon\ve{E}
  = i\omega\!\left(\varepsilon - \frac{i\sigma}{\omega}\right)\ve{E}.
\]
Así, en frecuencia, corriente de conducción y corriente de desplazamiento se
reagrupan en una única permitividad compleja efectiva. Las ecuaciones formales
mantienen la estructura de Maxwell en vacío, pero con \emph{permitividad
dieléctrica generalizada}:

\begin{mdframed}[style=resultado]
\begin{equation}
  \varepsilon_\text{gen} = \varepsilon - \frac{i\sigma}{\omega}
  = \varepsilon' + i\varepsilon'' - \frac{i\sigma}{\omega},
  \label{eq:epsilongen}
\end{equation}
\end{mdframed}

donde $\varepsilon'$ y $\varepsilon''$ son las partes real e imaginaria de
$\varepsilon$, debidas a las cargas ligadas, y el término $-i\sigma/\omega$
proviene de las cargas libres.

Las ecuaciones de Maxwell en la materia quedan entonces:
\begin{align}
  \div(\varepsilon_\text{gen}\ve{E}) &= 0, \label{eq:mat1}\\[4pt]
  \div\ve{H} &= 0, \label{eq:mat2}\\[4pt]
  \curl\ve{E} &= -i\omega\mu\ve{H}, \label{eq:mat3}\\[4pt]
  \curl\ve{H} &= i\omega\varepsilon_\text{gen}\ve{E}. \label{eq:mat4}
\end{align}

El vector de Poynting en la materia es $\ve{S} = \ve{E}\times\ve{H}$.

% ══════════════════════════════════════════════════════════════════════════
\subsection{Clasificación Óptica de los Medios Materiales}
% ══════════════════════════════════════════════════════════════════════════

Toda la información sobre las propiedades ópticas del medio está contenida
en $\varepsilon_\text{gen}(\ve{r},\omega)$, que determina frecuencia,
dirección de propagación, velocidad, energía y polarización de la onda.

\begin{center}
\begin{tabularx}{0.95\linewidth}{@{}B{3.2cm}Y Y@{}}
  \toprule
  Clasificación & \textbf{Homogéneo} & \textbf{Inhomogéneo} \\
  \midrule
  Definición & $\varepsilon_\text{gen}$ no depende de $\ve{r}$ &
               $\varepsilon_\text{gen}=\varepsilon_\text{gen}(\ve{r})$ \\
  Ejemplo & Cristal, vidrio & Capas de aire a distinta $T$ \\
  \bottomrule
\end{tabularx}
\end{center}

\medskip
\begin{center}
\begin{tabularx}{0.95\linewidth}{@{}B{3.2cm}Y Y@{}}
  \toprule
  Clasificación & \textbf{Dispersivo} & \textbf{No dispersivo} \\
  \midrule
  Definición & $\varepsilon_\text{gen}=\varepsilon_\text{gen}(\omega)$ &
               $\varepsilon_\text{gen}$ no depende de $\omega$ \\
  Nota & Todos los medios reales son dispersivos & Aproximación en
         $\omega\gg\omega_0$ o $\omega\ll\omega_0$ \\
  \bottomrule
\end{tabularx}
\end{center}

\medskip
\begin{center}
\begin{tabularx}{0.95\linewidth}{@{}B{3.2cm}Y Y@{}}
  \toprule
  Clasificación & \textbf{Isótropo} & \textbf{Anisótropo} \\
  \midrule
  Definición & $\varepsilon_\text{gen}$ es un \textbf{escalar} &
               $\varepsilon_\text{gen}$ es un \textbf{tensor} \\
  Ejemplo & Vidrio, gas & Cristal birrefringente \\
  \bottomrule
\end{tabularx}
\end{center}

\medskip
\begin{center}
\begin{tabularx}{0.95\linewidth}{@{}B{3.2cm}Y Y@{}}
  \toprule
  Clasificación & \textbf{Conductor} & \textbf{No conductor} \\
  \midrule
  Definición & $\sigma\neq 0$ & $\sigma = 0$ \\
  Consecuencia & $\varepsilon_\text{gen}$ compleja & $\varepsilon_\text{gen}$
                 puede ser real \\
  \bottomrule
\end{tabularx}
\end{center}

\medskip
\begin{center}
\begin{tabularx}{0.95\linewidth}{@{}B{3.2cm}Y Y@{}}
  \toprule
  Clasificación & \textbf{Absorbente} & \textbf{Transparente} \\
  \midrule
  Definición & $\varepsilon_\text{gen}\notin\R$ & $\varepsilon_\text{gen}\in\R$ \\
  Origen & $\sigma\neq 0$ y/o $\varepsilon''\neq 0$ & $\sigma=0$ y $\varepsilon''=0$ \\
  \bottomrule
\end{tabularx}
\end{center}

\begin{ejemplo}[Clasificación de un medio]
Sea un medio con $\varepsilon_\text{gen} = \varepsilon_0\!\left(5 +
\dfrac{i\cdot10^{13}}{\omega}\right)$.

\medskip
\begin{itemize}
  \item $\varepsilon_\text{gen}\neq\varepsilon_\text{gen}(\ve{r})$:
        \textbf{homogéneo}.
  \item $\varepsilon_\text{gen}=\varepsilon_\text{gen}(\omega)$:
        \textbf{dispersivo}.
  \item $\varepsilon_\text{gen}$ escalar: \textbf{isótropo}.
  \item Parte imaginaria $\neq 0$: \textbf{absorbente} (tiene conductividad
        $\sigma = \varepsilon_0\cdot10^{13}\ \text{Hz}$).
\end{itemize}
\end{ejemplo}

% ══════════════════════════════════════════════════════════════════════════
\subsection{Propagación en Medios Homogéneos e Isótropos}
% ══════════════════════════════════════════════════════════════════════════

En un medio homogéneo ($\varepsilon_\text{gen}\neq\varepsilon_\text{gen}(\ve{r})$)
e isótropo ($\varepsilon_\text{gen}$ escalar), buscamos soluciones armónicas
de las ecuaciones de Maxwell.

\subsubsection{Ecuación de Helmholtz}

\begin{teorema}[Ecuación de Helmholtz]
En un medio homogéneo e isótropo, el campo eléctrico satisface:
\begin{mdframed}[style=resultado]
\begin{equation}
  \nabla^2\ve{E} + \frac{\omega^2}{c^2}n_c^2\,\ve{E} = 0,
  \label{eq:helmholtz}
\end{equation}
\end{mdframed}
donde el \emph{índice de refracción complejo} es:
\begin{equation}
  n_c = \sqrt{\mu_r\varepsilon_\text{gen}^r} = \eta - i\kappa,
  \label{eq:nc}
\end{equation}
con $\varepsilon_\text{gen}^r = \varepsilon_\text{gen}/\varepsilon_0$,
$\eta > 0$ la \emph{parte real} del índice (responsable de la velocidad de
fase $v_f = c/\eta$) y $\kappa \geq 0$ el \emph{coeficiente de extinción}
(responsable de la absorción). La convención de signo $n_c = \eta - i\kappa$
corresponde a la dependencia temporal $e^{+i\omega t}$ adoptada en estos apuntes.
En medios no absorbentes, $\kappa = 0$ y $n_c = \eta \in \R$ es el índice de
refracción real habitual.
\end{teorema}

\begin{proof}
Aplicamos el rotacional a~\eqref{eq:mat3} y usamos~\eqref{eq:mat4}:
\[
  \curl(\curl\ve{E})
  = -i\omega\mu(\curl\ve{H})
  = -i\omega\mu\cdot i\omega\varepsilon_\text{gen}\ve{E}
  = \omega^2\mu\varepsilon_\text{gen}\ve{E}.
\]
Además:
\[
  \div(\varepsilon_\text{gen}\ve{E})=0
  \quad\Rightarrow\quad
  \div\ve{E}=0.
\]
Entonces:
\[
  -\nabla^2\ve{E} = \omega^2\mu\varepsilon_\text{gen}\ve{E}.
\]
Con:
\[
  \mu = \mu_0\mu_r,
  \qquad
  \varepsilon_\text{gen} = \varepsilon_0\varepsilon_\text{gen}^r,
  \qquad
  \mu_0\varepsilon_0 = \frac{1}{c^2},
\]
y tomando $\mu_r\approx 1$ en frecuencias ópticas:
\[
  \nabla^2\ve{E} + \frac{\omega^2}{c^2}\varepsilon_\text{gen}^r\ve{E} = 0.
\]
Al identificar:
\[
  n_c^2 = \mu_r\varepsilon_\text{gen}^r \approx \varepsilon_\text{gen}^r,
\]
se obtiene~\eqref{eq:helmholtz}.
\end{proof}

\subsubsection{Solución: ondas armónicas planas con vector de onda complejo}

Proponemos una solución de onda plana con \textbf{vector de onda complejo}:
\[
  \ve{k}_c = \ve{k} - i\ve{a},
\]
donde $\ve{k}$ es el vector de fase y $\ve{a}$ es el vector de atenuación.
El campo resulta:
\begin{equation}
  \ve{E} = \ve{E}_0\,e^{i\omega t}\,e^{-i\ve{k}\cdot\ve{r}}\,e^{-\ve{a}\cdot\ve{r}}.
  \label{eq:onda_compleja}
\end{equation}
Los frentes de \emph{fase constante} vienen dados por $\omega t - \ve{k}\cdot\ve{r} = \mathrm{cte}$,
y los de \emph{amplitud constante} por $\ve{a}\cdot\ve{r}=\mathrm{cte}$.

Sustituyendo en las ecuaciones de Maxwell~\eqref{eq:mat1}--\eqref{eq:mat4}
con el reemplazo $\nabla\to -i\ve{k}_c$:

\begin{center}
\begin{tabularx}{0.95\linewidth}{@{}L{5.2cm}Y@{}}
  \toprule
  Ecuación & Resultado \\
  \midrule
  $\div(\varepsilon_\text{gen}\ve{E})=0$ & $\ve{k}_c\cdot\ve{E}=0$ \\
  $\div\ve{H}=0$                          & $\ve{k}_c\cdot\ve{H}=0$ \\
  $\curl\ve{E}=-i\omega\mu\ve{H}$         & $\ve{k}_c\times\ve{E}=\omega\mu\ve{H}$ \\
  $\curl\ve{H}=i\omega\varepsilon_\text{gen}\ve{E}$ &
    $\ve{k}_c\times\ve{H}=-\omega\varepsilon_\text{gen}\ve{E}$ \\
  \bottomrule
\end{tabularx}
\end{center}

\begin{proposicion}[Relaciones entre $k$, $a$, $\eta$ y $\kappa$]
Aplicando $\ve{k}_c\times(\ve{k}_c\times\ve{E})$ y usando la identidad
vectorial se obtienen:
\begin{mdframed}[style=resultado]
\begin{equation}
  k^2 - a^2 = \frac{\omega^2}{c^2}(\eta^2-\kappa^2), \qquad
  \ve{k}\cdot\ve{a} = \frac{\omega^2}{c^2}\,\eta\kappa.
  \label{eq:relaciones_ka}
\end{equation}
\end{mdframed}
\end{proposicion}

\begin{proof}
Partimos de:
\[
  \ve{k}_c\times\ve{E}=\omega\mu\ve{H}.
\]
Aplicando $\ve{k}_c\times$:
\[
  \ve{k}_c\times(\ve{k}_c\times\ve{E})
  = \omega\mu(\ve{k}_c\times\ve{H})
  = \omega\mu\cdot(-\omega\varepsilon_\text{gen})\ve{E}
  = -\omega^2\mu\varepsilon_\text{gen}\ve{E}.
\]
Usamos la identidad del doble producto vectorial:
\[
  \ve{A}\times(\ve{B}\times\ve{C})
  = \ve{B}(\ve{A}\cdot\ve{C})-\ve{C}(\ve{A}\cdot\ve{B}).
\]
Y la condición de transversalidad:
\[
  \ve{k}_c\cdot\ve{E}=0.
\]
\[
  -\ve{E}(\ve{k}_c\cdot\ve{k}_c) = -\omega^2\mu\varepsilon_\text{gen}\ve{E}
  \implies k_c^2 = \frac{\omega^2}{c^2}n_c^2.
\]
Expandimos:
\[
  k_c^2 = (k-ia)^2 = k^2-a^2-2i\ve{k}\cdot\ve{a},
\]
\[
  n_c^2=(\eta-i\kappa)^2=\eta^2-\kappa^2-2i\eta\kappa.
\]
Sustituyendo ambas expansiones en:
\[
  k_c^2=\frac{\omega^2}{c^2}\,n_c^2,
\]
\[
  k^2-a^2-2i\ve{k}\cdot\ve{a}
  =
  \frac{\omega^2}{c^2}\left(\eta^2-\kappa^2-2i\eta\kappa\right).
\]
Igualando parte real:
\[
  k^2-a^2=\frac{\omega^2}{c^2}(\eta^2-\kappa^2).
\]
Igualando parte imaginaria:
\[
  \ve{k}\cdot\ve{a}=\frac{\omega^2}{c^2}\eta\kappa.
\]
que coincide con~\eqref{eq:relaciones_ka}.
\end{proof}

\subsubsection{Clasificación de medios según $\kappa$}

\paragraph{Medios no absorbentes: $\kappa = 0$}

Las ecuaciones~\eqref{eq:relaciones_ka} se simplifican a
$k^2 - a^2 = \omega^2\eta^2/c^2$ y $\ve{k}\cdot\ve{a}=0$.

\begin{description}
  \item[\textbf{Caso 1: $\ve{a}=0$ (onda homogénea).}]
        No hay atenuación. La onda se propaga sin pérdidas con:
        \begin{equation}
          k = k_0\eta = \frac{\omega}{c}\eta, \qquad
          v_f = \frac{\omega}{k} = \frac{c}{\eta}, \qquad
          \lambda_\text{medio} = \frac{\lambda_0}{\eta}. \label{eq:homogenea}
        \end{equation}

  \item[\textbf{Caso 2: $\ve{k}\perp\ve{a}$ (onda evanescente).}]
        La fase se propaga en una dirección ($z$) y la amplitud decae en
        otra ($x$):
        \[
          \ve{E} = \ve{E}_0\,e^{-ax}\,e^{-ikz}\,e^{i\omega t}.
        \]
        Es una \textbf{onda inhomogénea} o evanescente: los frentes de
        amplitud constante no coinciden con los de fase constante.
\end{description}

\paragraph{Medios absorbentes: $\kappa \neq 0$}

\begin{description}
  \item[\textbf{Caso 1: $\ve{k}\parallel\ve{a}$ (onda amortiguada).}]
        $k = \eta k_0$ y $a = \kappa k_0$. El campo vale:
        \begin{equation}
          \ve{E} = \ve{E}_0\,e^{-\kappa k_0 r}\,e^{-i\eta k_0 r}\,e^{i\omega t}.
          \label{eq:absorbente}
        \end{equation}
        La \emph{longitud de penetración} $d$ es la distancia a la que la
        amplitud cae a $1/e$ de su valor inicial:
        \begin{mdframed}[style=resultado]
        \begin{equation}
          d = \frac{1}{k_0\kappa} = \frac{\lambda_0}{2\pi\kappa}.
          \label{eq:long_penetracion}
        \end{equation}
        \end{mdframed}
        donde $k_0 = 2\pi/\lambda_0$ es el número de onda en el vacío y $\kappa$ el coeficiente de extinción del índice complejo $n_c = \eta - i\kappa$.

  \item[\textbf{Caso 2: $\ve{k}\not\parallel\ve{a}$ (caso general).}]
        La atenuación tiene componentes en todas las direcciones. Se usan las
        ecuaciones~\eqref{eq:relaciones_ka} con los valores concretos de
        $\eta$ y $\kappa$ para hallar $\ve{k}$ y $\ve{a}$.
\end{description}

\subsubsection{Vector de Poynting en la materia}

Para una onda con $\ve{H} = \frac{1}{\omega\mu}\ve{k}_c\times\ve{E}$, el
promedio del vector de Poynting es (tomando solo la parte real de $\ve{k}_c$):
\begin{equation}
  \langle\ve{S}\rangle = \frac{\eta}{2\mu c}\,E_0^2\,\uvec{k}.
  \label{eq:poynting_mat}
\end{equation}

\begin{nota}[Velocidad de fase vs.\ velocidad de grupo]
En un medio dispersivo, la velocidad de fase $v_f = c/\eta$ y la velocidad
de grupo $v_g = d\omega/dk$ son en general distintas. La energía y la
información se propagan a la velocidad de grupo:
\[
  v_g = \frac{c}{\eta + \omega\,d\eta/d\omega}.
\]
En dispersión normal ($d\eta/d\omega > 0$), $v_g < v_f < c$. En dispersión
anómala ($d\eta/d\omega < 0$ y de gran magnitud), $v_g$ puede exceder $c$
o incluso volverse negativa, pero esto no viola la relatividad especial
porque la señal óptica útil se propaga a la velocidad del frente de onda,
que siempre es $\leq c$.
\end{nota}

% ══════════════════════════════════════════════════════════════════════════
\subsection{Modelo Microscópico del Índice de Refracción}
% ══════════════════════════════════════════════════════════════════════════

El objetivo es conectar el índice de refracción complejo $n_c = \eta - i\kappa$
con los parámetros microscópicos del medio (Tema~2).

\subsubsection{Medios dieléctricos poco densos}

En un medio diluido, las cargas no interactúan entre sí. El dipolo de cada
átomo en el modelo de Lorentz es (Tema~2):
\[
  \ve{p} = q\ve{r} = \frac{q^2/m}{\omega_0^2-\omega^2+i\gamma\omega}\ve{E}_0.
\]
La polarización total del medio es $\ve{P} = N\ve{p}$, que igualada a la
relación de constitución $\ve{P} = \varepsilon_0\chi_e\ve{E}$ da:
\[
  \chi_e = \frac{Nq^2}{m\varepsilon_0(\omega_0^2-\omega^2+i\gamma\omega)}.
\]
Usando $n_c^2 = 1 + \chi_e$ y separando real e imaginaria:

\begin{mdframed}[style=resultado]
\begin{align}
  n_c^2 &= 1 + \frac{\omega_p^2}{\omega_0^2-\omega^2+i\gamma\omega},
  \label{eq:nc2_diluido}\\[6pt]
  \eta &= 1 + \frac{(\omega_0^2-\omega^2)\,\omega_p^2}
              {(\omega_0^2-\omega^2)^2+\gamma^2\omega^2},
  \label{eq:n_diluido}\\[6pt]
  \kappa &= \frac{\gamma\omega\,\omega_p^2}
              {(\omega_0^2-\omega^2)^2+\gamma^2\omega^2},
  \label{eq:kappa_diluido}
\end{align}
\end{mdframed}
donde la \emph{frecuencia de plasma} es:
\begin{equation}
  \omega_p^2 = \frac{Nq^2}{2m\varepsilon_0}. \label{eq:plasma}
\end{equation}

\paragraph{Bandas de absorción y transparencia}

La función $\kappa(\omega)$ presenta un \textbf{pico de absorción} centrado
en $\omega = \omega_0$ con anchura del orden de $\gamma$. Fuera de las bandas
de absorción ($|\omega-\omega_0|\gg\gamma$), $\kappa\approx 0$ y el medio
es transparente (\emph{zona de transparencia}).

\begin{figure}[H]
  \centering
  \begin{tikzpicture}[>=Stealth, scale=1.1]
    % Ejes
    \draw[->] (0,0) -- (8.5,0) node[right] {$\omega$};
    \draw[->] (0,-0.8) -- (0,2.8) node[above] {};
    % Línea de referencia n=1
    \draw[thin, dotted] (0,1.0) -- (8.5,1.0) node[right] {\small $1$};
    % Línea de resonancia
    \draw[thin, dotted] (4.0,-0.8) -- (4.0,2.8);
    \node[below] at (4.0,-0.9) {$\omega_0$};

    % eta(omega): dispersión normal izquierda, anómala en resonancia, normal derecha
    % Izquierda (omega << omega0): crece lentamente
    \draw[thick] plot[smooth] coordinates {
      (0.2, 1.15) (0.8,1.18) (1.5,1.22) (2.2,1.28) (2.8,1.38)
      (3.2,1.55) (3.6,1.90) (3.85, 2.50)};
    % Salto y parte derecha
    \draw[thick] plot[smooth] coordinates {
      (4.15,-0.55) (4.4, 0.20) (4.7, 0.55) (5.2,0.72)
      (5.8,0.82) (6.4,0.88) (7.5,0.93) (8.2,0.96)};
    \node[left] at (0,1.0) {\small $\eta$};

    % kappa(omega): pico en resonancia
    \draw[thick, dashed] plot[smooth] coordinates {
      (0.2,0.02) (1.5,0.04) (2.5,0.08) (3.2,0.20)
      (3.6,0.55) (3.85,1.40) (4.0,1.80)
      (4.15,1.40) (4.4,0.55) (4.7,0.20)
      (5.2,0.08) (6.0,0.03) (7.5,0.01)};
    \node[align=left, anchor=west] at (8.8,2.0)
      {\footnotesize Línea continua: $\eta(\omega)$\\
       \footnotesize Línea discontinua: $\kappa(\omega)$\\
       \footnotesize Zona B: dispersión normal\\
       \footnotesize Zona A: dispersión anómala};
  \end{tikzpicture}
  \caption{Variación cualitativa de $\eta(\omega)$ (línea continua) y
           $\kappa(\omega)$ (línea discontinua) para un dieléctrico diluido.
           La zona~B (izquierda de $\omega_0$) muestra dispersión normal;
           la zona~A (en $\omega_0$) muestra dispersión anómala y máxima absorción.}
  \label{fig:dispersion}
\end{figure}

\paragraph{Dispersión normal y anómala}

\begin{description}
  \item[\textbf{Dispersión normal.}] $d\eta/d\omega > 0$: el índice de
        refracción aumenta con la frecuencia (disminuye con $\lambda$). Es el
        comportamiento habitual fuera de las bandas de absorción. La mayoría de
        los vidrios ópticos exhiben dispersión normal en el rango visible,
        lo que explica que un prisma disperse el rojo hacia un ángulo menor
        que el violeta.

  \item[\textbf{Dispersión anómala.}] $d\eta/d\omega < 0$: el índice de
        refracción disminuye al aumentar la frecuencia. Ocurre en las
        inmediaciones de $\omega_0$, coincidiendo con la banda de absorción.
        En esta zona, la velocidad de grupo puede superar $c$ o incluso
        volverse negativa; sin embargo, la señal óptica útil sigue
        propagándose a velocidad $\leq c$ de acuerdo con la causalidad.
\end{description}

\begin{observacion}[Relaciones de Kramers-Kronig]
Las partes real $\eta(\omega)$ e imaginaria $\kappa(\omega)$ del índice de
refracción complejo no son independientes: están relacionadas por las
\emph{relaciones de Kramers-Kronig}, que son una consecuencia directa de la
causalidad (la respuesta del medio no puede preceder a la excitación):
\[
  \eta(\omega) - 1 = \frac{2}{\pi}\,\mathcal{P}\!\!\int_0^\infty
  \frac{\omega'\kappa(\omega')}{\omega'^2-\omega^2}\,d\omega',
  \qquad
  \kappa(\omega) = -\frac{2\omega}{\pi}\,\mathcal{P}\!\!\int_0^\infty
  \frac{\eta(\omega')-1}{\omega'^2-\omega^2}\,d\omega',
\]
donde $\mathcal{P}$ denota el valor principal de Cauchy. En consecuencia,
si se conoce el espectro de absorción $\kappa(\omega)$ para todas las
frecuencias, el índice real $\eta(\omega)$ queda completamente determinado,
y viceversa. Estas relaciones se usan ampliamente en espectroscopía para
reconstruir uno de los dos componentes a partir del otro.
\end{observacion}

\paragraph{Casos límite: leyes de Cauchy}

\begin{proposicion}[Ley de Cauchy — $\omega\ll\omega_0$]
Lejos de la resonancia por debajo ($\omega\ll\omega_0$, zona visible para
muchos materiales), $\gamma\omega\ll\omega_0^2$:
\[
  \eta \approx 1 + \frac{\omega_p^2}{\omega_0^2-\omega^2}
     \approx 1 + \frac{\omega_p^2}{\omega_0^2}\!\left(1+\frac{\omega^2}{\omega_0^2}\right),
\]
que en función de $\lambda$ (con $\omega=2\pi c/\lambda$) da:
\begin{equation}
  \eta = A + \frac{B}{\lambda^2} \quad\text{(Ley de Cauchy).}
  \label{eq:cauchy}
\end{equation}
\end{proposicion}

\begin{proof}
Bajo las hipótesis:
\[
  \gamma\omega\ll(\omega_0^2-\omega^2),
  \qquad
  \omega\ll\omega_0,
\]
se desarrolla:
\[
  \left(1-\frac{\omega^2}{\omega_0^2}\right)^{-1}
  \approx
  1+\frac{\omega^2}{\omega_0^2}.
\]
Así:
\[
  \eta \approx 1 + \frac{\omega_p^2}{\omega_0^2-\omega^2}
  = 1 + \frac{\omega_p^2}{\omega_0^2}
      \left(1-\frac{\omega^2}{\omega_0^2}\right)^{-1}
  \approx
  1+\frac{\omega_p^2}{\omega_0^2}
  +\frac{\omega_p^2}{\omega_0^4}\omega^2.
\]
Usando:
\[
  \omega=\frac{2\pi c}{\lambda},
\]
\[
  \eta \approx
  \left(1+\frac{\omega_p^2}{\omega_0^2}\right)
  +\frac{\omega_p^2(2\pi c)^2}{\omega_0^4}\frac{1}{\lambda^2}.
\]
Identificando:
\[
  A=1+\frac{\omega_p^2}{\omega_0^2},
  \qquad
  B=\frac{\omega_p^2(2\pi c)^2}{\omega_0^4},
\]
se obtiene~\eqref{eq:cauchy}.
\end{proof}

\begin{proposicion}[Dispersión por encima de la resonancia — $\omega\gg\omega_0$]
\begin{equation}
  \eta = A - B\lambda^2 - C\lambda^4, \label{eq:cauchy_alta}
\end{equation}
siendo el término $-B\lambda^2$ dominante. El índice decrece con $\omega$
(dispersión anómala).
\end{proposicion}

\subsubsection{Medios dieléctricos densos: ecuación de Clausius-Mossotti}

En medios diluidos era suficiente identificar campo microscópico y
macroscópico. En medios densos, esa identificación deja de ser válida:
cada dipolo siente un \emph{campo local} $\ve{E}_{\text{loc}}$ que incluye la
contribución de sus vecinos. Para simetría cúbica, la corrección local de
Lorentz conduce a:
\[
  \ve{E}_\text{loc} = \ve{E} + \frac{\ve{P}}{3\varepsilon_0}.
\]
Combinando esta expresión con $\ve{P}=N\chi_a\ve{E}_{\text{loc}}$ y comparando
con la ley macroscópica $\ve{P}=\varepsilon_0\chi_e\ve{E}$ se obtiene:

\begin{mdframed}[style=resultado]
\begin{equation}
  \frac{n_c^2 - 1}{n_c^2 + 2} = \frac{N\chi_a}{3\varepsilon_0}
  \qquad\text{(Clausius-Mossotti).} \label{eq:clausius}
\end{equation}
\end{mdframed}

\begin{proof}
Partimos de:
\[
  \ve{P}
  = N\chi_a\ve{E}_\text{loc}
  = N\chi_a\left(\ve{E}+\frac{\ve{P}}{3\varepsilon_0}\right).
\]
Despejando $\ve{P}$:
\[
  \ve{P} - \frac{N\chi_a}{3\varepsilon_0}\ve{P} = N\chi_a\ve{E},
\]
\[
  \ve{P}\left(1-\frac{N\chi_a}{3\varepsilon_0}\right)=N\chi_a\ve{E},
\]
\[
  \ve{P} = \frac{N\chi_a}{1 - N\chi_a/(3\varepsilon_0)}\ve{E}
  \equiv \varepsilon_0\chi_e\ve{E}.
\]
Definimos:
\[
  n_c^2 = 1+\chi_e,
  \qquad
  x = \frac{N\chi_a}{3\varepsilon_0}.
\]
Entonces:
\[
  n_c^2 = 1 + \frac{3x}{1-x} = \frac{1+2x}{1-x}.
\]
Multiplicando en cruz:
\[
  n_c^2(1-x)=1+2x
  \quad\Rightarrow\quad
  n_c^2-1=x(n_c^2+2).
\]
Finalmente:
\[
  x=\frac{n_c^2-1}{n_c^2+2}
  =\frac{N\chi_a}{3\varepsilon_0},
\]
que es~\eqref{eq:clausius}.
\end{proof}

\begin{observacion}
En un medio denso, el campo local modifica apreciablemente el comportamiento
espectral respecto al gas diluido. Dos consecuencias principales:
\begin{itemize}
  \item Las bandas de absorción se \emph{ensanchan} y se \emph{desplazan}
        en frecuencia respecto al resonador aislado, porque los dipolos
        vecinos actúan como perturbaciones sobre la frecuencia de resonancia
        efectiva.
  \item La ecuación de Clausius-Mossotti predice el índice de refracción a
        partir de la polarizabilidad atómica $\chi_a$, que puede medirse
        independientemente. Esto la convierte en un puente entre la respuesta
        microscópica (calculable con mecánica cuántica) y el índice macroscópico
        (observable ópticamente), con aplicaciones en el diseño de vidrios
        ópticos y cristales.
\end{itemize}
\end{observacion}

\subsubsection{Medios conductores ($\sigma\neq 0$)}

En un metal o semiconductor, además de los electrones ligados (que contribuyen
a $\chi_e$ como en un dieléctrico), existen \emph{electrones libres} que pueden
desplazarse macroscópicamente y dan lugar a una corriente de conducción
$\ve{j} = \sigma\ve{E}$. En las ecuaciones de Maxwell, esta corriente aparece
como un término adicional en la ley de Ampère:
\[
  \curl\ve{H} = \ve{j} + \pdv{\ve{D}}{t}
  = \sigma\ve{E} + \varepsilon_0\varepsilon_r\pdv{\ve{E}}{t}.
\]
Para una onda armónica $\ve{E} \propto e^{i\omega t}$, este término produce
una permitividad compleja efectiva:
\[
  \varepsilon_\text{ef} = \varepsilon_r + \frac{i\sigma}{\varepsilon_0\omega},
\]
que absorbe la contribución de los libres en el índice complejo $n_c$.

Modelando el electrón libre como un oscilador de Lorentz sin fuerza recuperadora
($\omega_0 = 0$) con amortiguamiento $\Gamma$ (debido a colisiones con la red),
la ecuación de movimiento $m\ddot{\ve{r}}+m\Gamma\dot{\ve{r}} = q\ve{E}$
conduce, con $\ve{j} = Nq\dot{\ve{r}} = \sigma\ve{E}$, a:
\[
  \sigma = \frac{Nq^2}{m(\Gamma + i\omega)}, \qquad
  \sigma_0 = \frac{Nq^2}{m\Gamma} \quad(\text{conductividad estática}).
\]
El índice de refracción complejo resulta:
\begin{equation}
  n_c^2 = 1 + \chi_e - i\frac{\omega_p^2}{\Gamma\omega + i\omega^2}
        = 1 - \frac{\omega_p^2}{\Gamma^2\omega^2+\omega^4}(\Gamma^2\omega^2+\omega^4)^{-1}
          \cdot(\omega^2 + i\Gamma\omega)\omega_p^2.
  \label{eq:nc_conductor}
\end{equation}

\medskip
Se distinguen dos regímenes:

\medskip
\begin{description}
  \item[\textbf{Caso 1: $\Gamma\ll\omega$ (poco rozamiento).}]
        $n_c^2 \approx 1 - \omega_p^2/\omega^2 - i\,2\omega_p^2\Gamma/\omega^3$.
        \begin{itemize}
          \item Si $\omega>\omega_p$: $\kappa\approx 0$, medio \textbf{transparente}.
          \item Si $\omega<\omega_p$: $\kappa\neq 0$, medio \textbf{absorbente}.
        \end{itemize}
        La frecuencia de plasma $\omega_p$ actúa como frecuencia de corte.

  \item[\textbf{Caso 2: $\Gamma\gg\omega$ (mucho rozamiento — buen conductor).}]
        $\varepsilon_r \ll \sigma/(\varepsilon_0\omega)$. Entonces:
        \[
          n_c^2 \approx -i\frac{\omega_p^2}{\Gamma\omega}
        \]
        y la parte real e imaginaria del índice son iguales:
        \begin{mdframed}[style=resultado]
        \begin{equation}
          \eta = \kappa = \sqrt{\frac{\omega_p^2}{2\Gamma\omega}}
               = \sqrt{\frac{\sigma_0}{2\varepsilon_0\omega}}.
          \label{eq:conductor_grueso}
        \end{equation}
        \end{mdframed}
        La longitud de penetración es la \emph{profundidad de efecto piel}:
        $d = 1/(k_0\kappa)$. A mayor frecuencia, menor penetración.
\end{description}

\begin{center}
\begin{tabularx}{0.98\linewidth}{@{}L{2.6cm}L{2.4cm}L{2.7cm}Y@{}}
  \toprule
  Régimen & Condición & Transparente si & $\eta$ y $\kappa$ \\
  \midrule
  Poco rozamiento & $\Gamma\ll\omega$ & $\omega>\omega_p$ & $\eta\neq\kappa$ \\
  Buen conductor  & $\Gamma\gg\omega$ & nunca & $\eta=\kappa=\sqrt{\sigma_0/(2\varepsilon_0\omega)}$ \\
  \bottomrule
\end{tabularx}
\end{center}

% ══════════════════════════════════════════════════════════════════════════
\section{Tema 4: Reflexión y Refracción en Medios Homogéneos e Isótropos}
% ══════════════════════════════════════════════════════════════════════════

\subsection{Introducción}

Los temas anteriores han estudiado la propagación de la luz en medios
homogéneos ilimitados. Ahora abordamos el problema más rico y de mayor
relevancia práctica: \textbf{¿qué ocurre cuando una onda plana incide sobre
la superficie de separación de dos medios con índices de refracción distintos?}

Este problema —la reflexión y la refracción en una interfaz— es el fundamento
teórico de la óptica instrumental (lentes, espejos, prismas, recubrimientos
antirreflectantes, fibra óptica) y de la espectroscopia de superficies
(elipsometría, reflectometría).

El problema se formula con tres preguntas bien diferenciadas:
\begin{enumerate}
  \item Determinar la \emph{dirección de propagación} de las ondas reflejada
        y transmitida (leyes de reflexión y de Snell).
  \item Determinar el \emph{estado de polarización} de los campos resultantes
        (fórmulas de Fresnel).
  \item Cuantificar la \emph{energía} transmitida y reflejada en función del
        ángulo de incidencia y del estado de polarización (reflectancia y
        transmitancia).
\end{enumerate}

\subsubsection*{Hipótesis de trabajo}

\begin{enumerate}
  \item Los medios son \textbf{isótropos y homogéneos}: el índice de refracción
        $n$ es un escalar independiente de la posición y de la dirección de
        propagación.
  \item La superficie de separación es \textbf{plana}.
  \item Cambio de índice $n_1 \to n_2$; si $n_1 = n_2$ no existe superficie
        óptica.
\end{enumerate}

Estas hipótesis delimitan el alcance del tema: interfaz plana, medios lineales
isótropos y ondas armónicas planas. Los casos anisótropos o con absorción se
tratan posteriormente.

\subsubsection*{Notación geométrica}

\begin{itemize}
  \item El \emph{plano de incidencia} es el plano $XZ$; la interfaz está en
        $z = 0$; la normal exterior apunta en $+\uvec{z}$.
  \item Los subíndices $i$, $r$, $t$ denotan el campo incidente, reflejado
        y transmitido, respectivamente.
  \item $\theta_i$, $\theta_r$, $\theta_t$: ángulos que los vectores de onda
        forman con la normal $\uvec{z}$.
\end{itemize}

\begin{figure}[H]
  \centering
  \begin{tikzpicture}[>=Stealth, scale=1.1]
    % Interfaz y normal
    \draw[thick] (-4,0) -- (4,0);
    \node[left] at (-3.9,0.25) {$n_1$};
    \node[left] at (-3.9,-0.25) {$n_2$};
    \draw[->, thick] (0,0) -- (0,2.0) node[above] {$\uvec{n}$};

    % Rayos
    \draw[->, very thick] (-2.6,1.8) -- (0,0) node[midway, above left] {$\ve{k}_i$};
    \draw[->, very thick] (0,0) -- (2.4,1.5) node[midway, above right] {$\ve{k}_r$};
    \draw[->, very thick] (0,0) -- (2.8,-1.6) node[midway, below right] {$\ve{k}_t$};

    % Angulos
    \draw (0,0.75) arc[start angle=90, end angle=145, radius=0.75];
    \node at (-0.55,0.95) {$\theta_i$};
    \draw (0,0.9) arc[start angle=90, end angle=32, radius=0.9];
    \node at (0.55,1.05) {$\theta_r$};
    \draw (0,-0.9) arc[start angle=-90, end angle=-30, radius=0.9];
    \node at (0.6,-1.0) {$\theta_t$};
  \end{tikzpicture}
  \caption{Geometría de incidencia en una interfaz plana entre dos medios:
           haz incidente, reflejado y transmitido.}
  \label{fig:interfaz_geo}
\end{figure}

Los tres campos son ondas planas armónicas:
\begin{align}
  \ve{E}_i &= \ve{E}\, e^{-i\ve{k}_i \cdot \ve{r}}\, e^{i\omega t},
  & \ve{k}_i &= (k_x, 0, k_z),\quad |\ve{k}_i| = \frac{\omega}{c}n_1, \\
  \ve{E}_r &= \ve{R}\, e^{-i\ve{k}_r \cdot \ve{r}}\, e^{i\omega t}, \\
  \ve{E}_t &= \ve{T}\, e^{-i\ve{k}_t \cdot \ve{r}}\, e^{i\omega t}.
\end{align}

En el plano de incidencia $XZ$, una elección explícita útil es:
\begin{align}
  \ve{k}_i &= \frac{\omega}{c}n_1\left(\sin\theta_i,\,0,\,\cos\theta_i\right), \\
  \ve{k}_r &= \frac{\omega}{c}n_1\left(\sin\theta_r,\,0,\,-\cos\theta_r\right), \\
  \ve{k}_t &= \frac{\omega}{c}n_2\left(\sin\theta_t,\,0,\,\cos\theta_t\right).
\end{align}

La condición de frontera fundamental es la \textbf{continuidad de las
componentes tangenciales} del campo eléctrico (y magnético) a ambos lados
de la interfaz $z=0$.

\begin{nota}
Existen dos enfoques equivalentes para llegar a los mismos resultados.
El enfoque \emph{microscópico} (Teorema de Oseen--Ewald): los campos esparcidos
por los dipolos inducidos en el medio 2 cancelan exactamente el haz incidente
transmitido en el medio 1 y construyen coherentemente la onda refractada en el
medio 2. El enfoque \emph{macroscópico}: imponer las condiciones de contorno de
las ecuaciones de Maxwell (continuidad de componentes tangenciales de $\ve{E}$
y $\ve{H}$) a través de la interfaz. Ambos enfoques producen idénticos
resultados; en este tema se adopta el macroscópico por su concisión algebraica.
\end{nota}

% ══════════════════════════════════════════════════════════════════════════
\subsection{Dieléctrico--Dieléctrico Transparente}
% ══════════════════════════════════════════════════════════════════════════

Consideramos ambos medios \textbf{transparentes}: $n_1, n_2 \in \R$,
$n_1, n_2 > 0$. Los vectores de onda son reales.

\subsubsection{Leyes de reflexión y de Snell}

La condición de contorno en $z = 0$ (plano $XY$, posición $\ve{r} = (x,y,0)$)
puede leerse como \emph{igualación de fase tangencial}: para que la igualdad
de campos sea válida en todo punto de la interfaz, las fases deben coincidir
para todo $(x,y)$:
\[
  \ve{k}_i \cdot \ve{r}\big|_{z=0}
  = \ve{k}_r \cdot \ve{r}\big|_{z=0}
  = \ve{k}_t \cdot \ve{r}\big|_{z=0}
  \quad \forall\, (x,y).
\]

Esto impone la igualdad de las componentes tangenciales:
\[
  k_{ix} = k_{rx} = k_{tx}, \qquad k_{iy} = k_{ry} = k_{ty} = 0.
\]

La segunda condición confirma que \textbf{el plano de incidencia se conserva}
(los tres haces son coplanarios). De la primera, con
$|\ve{k}_i| = |\ve{k}_r| = \frac{\omega}{c}n_1$ y
$|\ve{k}_t| = \frac{\omega}{c}n_2$:

\begin{proof}
Escribimos las componentes tangenciales:
\[
  k_{ix} = \frac{\omega}{c}n_1 \sin\theta_i,
  \qquad
  k_{rx} = \frac{\omega}{c}n_1 \sin\theta_r,
  \qquad
  k_{tx} = \frac{\omega}{c}n_2 \sin\theta_t.
\]
La continuidad de fase en la interfaz impone:
\[
  \frac{\omega}{c}n_1\sin\theta_i = \frac{\omega}{c}n_1\sin\theta_r
  \quad\Rightarrow\quad
  \sin\theta_i=\sin\theta_r.
\]
Como $0\le\theta_i,\theta_r\le\pi/2$:
\[
  \theta_i = \theta_r,
\]
\[
  \frac{\omega}{c}n_1\sin\theta_i = \frac{\omega}{c}n_2\sin\theta_t.
\]
Cancelando $\omega/c$:
\[
  n_1\sin\theta_i=n_2\sin\theta_t.
\]
\end{proof}

\begin{mdframed}[style=resultado]
\textbf{Ley de reflexión:}
\begin{equation}
  \theta_i = \theta_r.
  \label{eq:ley_reflexion}
\end{equation}
\textbf{Ley de Snell (refracción):}
\begin{equation}
  n_1 \sin\theta_i = n_2 \sin\theta_t.
  \label{eq:snell}
\end{equation}
\end{mdframed}

\begin{observacion}
La ley de Snell no es una regla geométrica ad hoc: es la consecuencia directa
de la conservación del vector de onda tangencial en la interfaz.
\end{observacion}

\begin{figure}[H]
  \centering
  \begin{tikzpicture}[>=Stealth, scale=1.0]
    % Prisma
    \draw[thick, fill=gray!10] (-1.8,-0.9) -- (0,1.2) -- (1.8,-0.9) -- cycle;

    % Haz blanco incidente
    \draw[->, very thick] (-4.5,0.2) -- (-1.3,0.2);

    % Haces dispersados
    \draw[->, thick]    (1.0,0.1) -- (4.1,0.7);
    \draw[->, thick] (1.0,0.05) -- (4.0,0.2);
    \draw[->, thick]   (1.0,0.0) -- (4.0,-0.5);
    \node[align=left, anchor=west] at (4.35,0.95)
      {\footnotesize Triángulo: prisma\\
       \footnotesize Flecha izquierda: haz incidente\\
       \footnotesize Flechas derechas: longitudes\\
       \footnotesize de onda separadas angularmente};
  \end{tikzpicture}
  \caption{Dispersión por refracción en prisma (base de la espectroscopia):
           distintas longitudes de onda salen con diferentes ángulos.}
  \label{fig:prisma_dispersion}
\end{figure}

\subsubsection{Fórmulas de Fresnel (amplitudes)}

El campo eléctrico $\ve{E} \perp \ve{k}$ se descompone en dos polarizaciones:

\begin{itemize}
  \item \textbf{Polarización $p$ (TM, paralela):} $E_\parallel$, componente
        contenida en el plano de incidencia $XZ$.
  \item \textbf{Polarización $s$ (TE, perpendicular):} $E_\perp$, componente
        perpendicular al plano de incidencia (dirección $y$).
\end{itemize}

En componentes cartesianas los campos incidente y reflejado son:
\begin{align*}
  \ve{E}_i &= (E_\parallel\cos\theta_i,\; E_\perp,\; -E_\parallel\sin\theta_i), \\
  \ve{R}   &= (-R_\parallel\cos\theta_r,\; R_\perp,\; -R_\parallel\sin\theta_r), \\
  \ve{T}   &= (T_\parallel\cos\theta_t,\; T_\perp,\; -T_\parallel\sin\theta_t).
\end{align*}

\paragraph{Condiciones de contorno}

Continuidad de $\ve{E}$ tangencial (componentes $x$ e $y$) en $z=0$:
\begin{align}
  E_\parallel\cos\theta_i - R_\parallel\cos\theta_r
    &= T_\parallel\cos\theta_t, \tag{a} \label{cc:Etg_p} \\
  E_\perp + R_\perp &= T_\perp. \tag{b} \label{cc:Etg_s}
\end{align}

Continuidad de $\ve{H}$ tangencial
($\ve{H} = \tfrac{n}{\mu_0 c}\,\uvec{k}\times\ve{E}$):
\begin{align}
  n_1 E_\parallel + n_1 R_\parallel
    &= n_2 T_\parallel, \tag{c} \label{cc:Htg_p} \\
  n_1(E_\perp - R_\perp)\cos\theta_i
    &= n_2 T_\perp\cos\theta_t. \tag{d} \label{cc:Htg_s}
\end{align}

Una forma equivalente (y útil para comprobar signos) es escribir:
\begin{align*}
  \ve{H}_i &= \frac{n_1}{\mu_0 c}
    \left(-E_\perp\cos\theta_i,\; E_\parallel,\; E_\perp\sin\theta_i\right), \\
  \ve{H}_r &= \frac{n_1}{\mu_0 c}
    \left(R_\perp\cos\theta_r,\; R_\parallel,\; R_\perp\sin\theta_r\right), \\
  \ve{H}_t &= \frac{n_2}{\mu_0 c}
    \left(-T_\perp\cos\theta_t,\; T_\parallel,\; T_\perp\sin\theta_t\right).
\end{align*}

\paragraph{Resolución para la polarización $p$}

\begin{proof}
Partimos de (\ref{cc:Etg_p}) y (\ref{cc:Htg_p}) con:
\[
  \theta_r = \theta_i.
\]
Las ecuaciones quedan:
\[
  E_\parallel\cos\theta_i - R_\parallel\cos\theta_i
  = T_\parallel\cos\theta_t,
\]
\[
  n_1(E_\parallel+R_\parallel)=n_2T_\parallel.
\]
Eliminamos $T_\parallel$ usando la segunda:
\[
  T_\parallel=\frac{n_1}{n_2}(E_\parallel+R_\parallel).
\]
Sustituyendo en la primera:
\[
  E_\parallel\cos\theta_i - R_\parallel\cos\theta_i
  = \frac{n_1}{n_2}(E_\parallel+R_\parallel)\cos\theta_t.
\]
Reordenando términos en $R_\parallel$:
\[
  R_\parallel\!\left(n_2\cos\theta_i+n_1\cos\theta_t\right)
  =E_\parallel\!\left(n_2\cos\theta_i-n_1\cos\theta_t\right).
\]
Entonces:
\[
  r_\parallel = \frac{R_\parallel}{E_\parallel}
  = \frac{n_2\cos\theta_i - n_1\cos\theta_t}{n_2\cos\theta_i + n_1\cos\theta_t},
\]
Para $t_\parallel$:
\[
  t_\parallel
  =\frac{T_\parallel}{E_\parallel}
  =\frac{n_1}{n_2}\left(1+r_\parallel\right)
  =\frac{2n_1\cos\theta_i}{n_2\cos\theta_i+n_1\cos\theta_t}.
\]
\end{proof}

\paragraph{Resolución para la polarización $s$}

\begin{proof}
Partimos de (\ref{cc:Etg_s}) y (\ref{cc:Htg_s}).
Estas ecuaciones son:
\[
  E_\perp+R_\perp=T_\perp,
\]
\[
  n_1(E_\perp-R_\perp)\cos\theta_i=n_2T_\perp\cos\theta_t.
\]
Sustituyendo:
\[
  T_\perp=E_\perp+R_\perp,
\]
en la segunda ecuación:
\[
  n_1(E_\perp-R_\perp)\cos\theta_i
  =n_2(E_\perp+R_\perp)\cos\theta_t.
\]
Reordenando:
\[
  R_\perp\!\left(n_1\cos\theta_i+n_2\cos\theta_t\right)
  =E_\perp\!\left(n_1\cos\theta_i-n_2\cos\theta_t\right).
\]
\[
  r_\perp=\frac{R_\perp}{E_\perp}
  =\frac{n_1\cos\theta_i-n_2\cos\theta_t}
         {n_1\cos\theta_i+n_2\cos\theta_t}.
\]
Luego:
\[
  t_\perp = \frac{T_\perp}{E_\perp}
  = 1+r_\perp
  = \frac{2n_1\cos\theta_i}{n_1\cos\theta_i+n_2\cos\theta_t}.
\]
\end{proof}

\begin{mdframed}[style=resultado,nobreak=true]
\textbf{Fórmulas de Fresnel (polarización $p$):}
\begin{align}
  r_\parallel &= \frac{n_2\cos\theta_i - n_1\cos\theta_t}
                     {n_2\cos\theta_i + n_1\cos\theta_t}
               = \frac{\tan(\theta_i-\theta_t)}{\tan(\theta_i+\theta_t)},
  \label{eq:rp} \\
  t_\parallel &= \frac{2n_1\cos\theta_i}
                     {n_2\cos\theta_i + n_1\cos\theta_t}
               = \frac{2\cos\theta_i\sin\theta_t}
                      {\sin(\theta_i+\theta_t)\cos(\theta_i-\theta_t)}.
  \label{eq:tp}
\end{align}
\end{mdframed}

\begin{mdframed}[style=resultado,nobreak=true]
\textbf{Fórmulas de Fresnel (polarización $s$):}
\begin{align}
  r_\perp &= \frac{n_1\cos\theta_i - n_2\cos\theta_t}
                  {n_1\cos\theta_i + n_2\cos\theta_t}
           = -\frac{\sin(\theta_i-\theta_t)}{\sin(\theta_i+\theta_t)},
  \label{eq:rs} \\
  t_\perp &= \frac{2n_1\cos\theta_i}
                  {n_1\cos\theta_i + n_2\cos\theta_t}
           = \frac{2\cos\theta_i\sin\theta_t}{\sin(\theta_i+\theta_t)}.
  \label{eq:ts}
\end{align}
\end{mdframed}

\begin{observacion}
Estas expresiones cuantifican, para cada polarización, cómo reparte la
interfaz la amplitud entre onda reflejada y transmitida. Las diferencias entre
$s$ y $p$ explican la polarización parcial en reflexión oblicua.
\end{observacion}

\begin{nota}
Las formas trigonométricas (columna derecha) son consecuencia directa de
aplicar la ley de Snell a los cocientes de índices.
\end{nota}

\subsubsection{Incidencia normal}

Para $\theta_i = \theta_r = \theta_t = 0^{\circ}$ las fórmulas de Fresnel se
simplifican considerablemente. Tomando el límite $\theta_i \to 0$ en
(\ref{eq:rp})--(\ref{eq:ts}):

\begin{mdframed}[style=resultado,nobreak=true]
\textbf{Incidencia normal:}
\begin{align}
  r_\parallel &= \frac{n_2 - n_1}{n_1 + n_2} = -r_\perp,
  \label{eq:r_normal} \\
  t_\parallel &= t_\perp = \frac{2n_1}{n_1 + n_2}.
  \label{eq:t_normal}
\end{align}
\end{mdframed}

El campo reflejado resultante, en términos del campo incidente
$(E_\parallel, E_\perp)^T$, es:
\[
  \ve{R} = r_\parallel \begin{pmatrix} E_\parallel \\ -E_\perp \end{pmatrix},
\]
es decir, la componente perpendicular adquiere un cambio de fase de
$\pi$ respecto a $E_\perp$.

\begin{itemize}
  \item Si $n_1 < n_2$: $r_\parallel > 0$, no hay cambio de fase adicional
        en la componente $p$.
  \item Si $n_1 > n_2$: $r_\parallel < 0$, el campo reflejado experimenta
        un desfase de $\pi$ (equivalente a inversión del campo). Este es el
        motivo por el que una lámina fina (espesor $\ll \lambda$) depositada
        sobre vidrio muestra destrucción en reflexión: el haz reflejado en la
        cara superior (aire$\to$lámina) invierte fase, el reflejado en la cara
        inferior (lámina$\to$vidrio, si $n_{\text{lám}} < n_{\text{vid}}$)
        también invierte fase; si el espesor óptico es $\lambda/4$, el desfase
        extra de $\pi/2 + \pi/2 = \pi$ hace que ambos reflejos interfieran
        destructivamente: principio de los \emph{recubrimientos antirreflectantes}.
\end{itemize}

La reflectancia (fracción de intensidad reflejada) vale:
\[
  \mathcal{R} = |r_\parallel|^2 = \left(\frac{n_1 - n_2}{n_1 + n_2}\right)^2.
\]

\begin{ejemplo}[Aire--Vidrio]
Para $n_1 = 1$ (aire) y $n_2 = 1{,}5$ (vidrio):
\[
  \mathcal{R} = \left(\frac{0{,}5}{2{,}5}\right)^2 = 0{,}04 = 4\%.
\]
El $4\%$ de la intensidad se refleja en cada interfaz. Una cámara fotográfica
con 10 superficies aire-vidrio perdería $(1-0.04)^{10} \approx 66\%$ de la
luz sin recubrimientos antirreflectantes; con recubrimientos de $\mathcal{R}
< 0.1\%$ por superficie, la pérdida total es menor del $1\%$.
\end{ejemplo}

\subsubsection{Ángulo de Brewster}

De la forma trigonométrica de $r_\parallel$ [ec.~(\ref{eq:rp})]:
\[
  r_\parallel = \frac{\tan(\theta_i - \theta_t)}{\tan(\theta_i + \theta_t)}.
\]
Esta expresión se anula cuando $\tan(\theta_i + \theta_t) \to \infty$, es
decir, cuando:
\[
  \theta_i + \theta_t = 90^{\circ}.
\]

\begin{proof}
En el ángulo de Brewster:
\[
  \theta_t = 90^{\circ} - \theta_i.
\]
Sustituyendo en la ley de Snell~(\ref{eq:snell}):
\[
  n_1\sin\theta_i = n_2\sin(90^{\circ} - \theta_i) = n_2\cos\theta_i
  \quad\Rightarrow\quad
  \tan\theta_B = \frac{n_2}{n_1}.
\]
Es decir:
\[
  \theta_B=\arctan\!\left(\frac{n_2}{n_1}\right).
\]
\end{proof}

\begin{mdframed}[style=resultado]
\textbf{Ángulo de Brewster:}
\begin{equation}
  \tan\theta_B = \frac{n_2}{n_1}.
  \label{eq:brewster}
\end{equation}
En $\theta_i = \theta_B$: $r_\parallel = 0$, el haz reflejado está
completamente polarizado en $s$ (polarización perpendicular).
\end{mdframed}

Geométricamente, en Brewster los rayos reflejado y transmitido son
ortogonales; energéticamente, toda la componente $p$ pasa al medio transmitido
sin reflexión especular.

\begin{ejemplo}[Aire--Vidrio]
Para $n_1 = 1$, $n_2 = 1{,}5$:
\[
  \theta_B = \arctan(1{,}5) \approx 56{,}3^{\circ}.
\]
\end{ejemplo}

Los coeficientes de transmisión en el ángulo de Brewster son:
\[
  t_\parallel = \frac{n_1}{n_2}, \qquad
  t_\perp = \frac{2n_1^2}{n_1^2 + n_2^2}.
\]

\begin{observacion}
El fenómeno de Brewster tiene importantes aplicaciones:
\begin{itemize}
  \item \emph{Láminas polarizadoras tipo Brewster}: una pila de láminas de
        vidrio inclinadas al ángulo de Brewster transmite prácticamente toda
        la polarización $p$ y refleja progresivamente la polarización $s$, de
        modo que el haz transmitido queda muy bien polarizado en $p$.
  \item \emph{Ventanas de Brewster} en láseres: láminas inclinadas al ángulo
        de Brewster en los extremos de la cavidad láser permiten extraer un haz
        perfectamente polarizado sin pérdidas de reflexión para la polarización $p$.
  \item El ángulo de Brewster se usa en elipsometría para medir el índice de
        refracción de superficies con gran precisión.
\end{itemize}
\end{observacion}

\subsubsection{Reflexión total interna y onda evanescente}

Si $n_1 > n_2$, la ley de Snell $n_1\sin\theta_i = n_2\sin\theta_t$ exige
$\sin\theta_t = (n_1/n_2)\sin\theta_i > \sin\theta_i$.
Existe un \textbf{ángulo crítico} $\theta_c$ para el que $\sin\theta_t = 1$:

\begin{mdframed}[style=resultado]
\textbf{Ángulo crítico:}
\begin{equation}
  \sin\theta_c = \frac{n_2}{n_1} \qquad (n_1 > n_2).
  \label{eq:angulo_critico}
\end{equation}
\end{mdframed}

Para $\theta_i > \theta_c$ no existe $\theta_t$ real. Se analizan tres
hipótesis para el campo transmitido:

\begin{enumerate}
  \item \textbf{H1}: $\ve{E}_t = 0$. Incompatible con las condiciones de
        contorno (no se puede satisfacer la continuidad de campos si el
        transmitido es nulo).
  \item \textbf{H2}: $\ve{k}_t$ real con $|\ve{k}_t| = (\omega/c)n_2$.
        Incompatible con la igualación de las fases tangenciales
        $k_{tx} = k_{ix} = (\omega/c)n_1\sin\theta_i > (\omega/c)n_2$.
  \item \textbf{H3}: $\ve{k}_t \in \C$. Solución correcta.
\end{enumerate}

La tercera opción es la única compatible con contorno y causalidad: la
transmisión no desaparece, sino que cambia de naturaleza (campo no propagante
normal a la interfaz).

\paragraph{Onda evanescente}

Se escribe $\ve{k}_{tc} = \ve{k}_t - i\,\ve{a}$, con $\ve{k}_t \perp \ve{a}$
(medio transparente, $\kappa = 0$). Las condiciones de contorno imponen:
\[
  k_{tx} = k_{ix} = \frac{\omega}{c}n_1\sin\theta_i \in \R, \qquad
  a_x = 0.
\]

Así $\ve{a} = (0, 0, a_z)$ con:
\[
  a_z = \frac{\omega}{c}\sqrt{n_1^2\sin^2\theta_i - n_2^2} > 0
  \quad (\theta_i > \theta_c).
\]

\begin{mdframed}[style=resultado]
\textbf{Onda evanescente:}
\begin{equation}
  \ve{E}_t = \ve{T}\,e^{-a_z z}\,e^{-ik_{tx} x}\,e^{i\omega t}.
  \label{eq:evanescente}
\end{equation}
La onda se \emph{propaga} a lo largo de $x$ (paralelo a la interfaz) y se
\emph{atenúa exponencialmente} en $z$ (dentro del medio 2). La profundidad de
penetración es $\delta = 1/a_z$.
\end{mdframed}

Interpretación física: no hay flujo neto de energía hacia el interior del medio
2 en dirección normal, pero sí existe campo superficial capaz de acoplar energía
si aparece un segundo medio cercano (túnel óptico).

La ley de reflexión $\theta_i = \theta_r$ sigue siendo válida. Los
coeficientes de Fresnel resultan ser números complejos de módulo unitario:
\[
  r_\parallel = e^{i\delta_\parallel}, \qquad r_\perp = e^{i\delta_\perp}.
\]

El desfase relativo entre polarizaciones satisface:
\begin{equation}
  \tan\!\left(\frac{\delta_\perp - \delta_\parallel}{2}\right)
  = \frac{\cos\theta_i\,\sqrt{\sin^2\theta_i - n_2^2/n_1^2}}{\sin^2\theta_i}.
  \label{eq:desfase_total}
\end{equation}

\begin{nota}
En los extremos: $\delta_\perp - \delta_\parallel = 0$ tanto para
$\theta_i = \theta_c$ como para incidencia rasante ($\theta_i \to 90^{\circ}$).
El desfase máximo entre polarizaciones se produce para un ángulo intermedio,
y es la base del \emph{rombo de Fresnel}: dos reflexiones totales sucesivas
a ángulos apropiados introducen un desfase acumulado de $\pi/2$, convirtiendo
luz linealmente polarizada a $45^{\circ}$ en luz circularmente polarizada sin
ningún elemento absorbente.
\end{nota}

\begin{observacion}
La onda evanescente generada en la reflexión total interna no transporta energía
neta en la dirección $z$ (normal a la interfaz): su componente $z$ del vector de
Poynting es puramente imaginaria. Sin embargo, si otro medio dieléctrico se sitúa
a distancia menor que $\lambda/(2\pi)$ de la interfaz, la onda evanescente puede
acoplarse a él y se produce la \emph{reflexión total frustrada}, principio de
funcionamiento de los divisores de haz por efecto túnel óptico y de la
microscopía de campo cercano (SNOM).
\end{observacion}

\subsubsection{Reflectancia y Transmitancia}

La energía transportada por la onda se cuantifica con la componente $z$ del
vector de Poynting promediado:
\[
  S_{i,z} = \frac{n_1}{2\mu_0 c}|\ve{E}|^2\cos\theta_i, \quad
  S_{r,z} = \frac{n_1}{2\mu_0 c}|\ve{R}|^2\cos\theta_r, \quad
  S_{t,z} = \frac{n_2}{2\mu_0 c}|\ve{T}|^2\cos\theta_t.
\]

La conservación de energía $S_{i,z} = S_{r,z} + S_{t,z}$ implica
$1 = \mathcal{R} + \mathcal{T}$.

\begin{mdframed}[style=resultado]
\textbf{Reflectancia y Transmitancia} (para cada polarización $j = \parallel, \perp$):
\begin{align}
  \mathcal{R}_j &= |r_j|^2,
  \label{eq:reflectancia} \\
  \mathcal{T}_j &= \frac{n_2\cos\theta_t}{n_1\cos\theta_i}\,|t_j|^2.
  \label{eq:transmitancia}
\end{align}
Se cumple $\mathcal{R}_j + \mathcal{T}_j = 1$ para cada polarización.
\end{mdframed}

\begin{nota}
El factor $n_2\cos\theta_t/(n_1\cos\theta_i)$ en $\mathcal{T}$ no es un simple
módulo cuadrado de $t_j$ porque la energía se calcula con el flujo de Poynting,
que contiene tanto los índices (velocidad de propagación) como el factor
geométrico $\cos\theta$ (proyección sobre la normal de la interfaz). La relación
$\mathcal{T} = 1-\mathcal{R}$ es consecuencia de la conservación de la energía
y puede verificarse directamente sustituyendo las fórmulas de Fresnel.
\end{nota}

La reflectancia y transmitancia totales son:
\[
  \mathcal{R} = \frac{|E_\parallel|^2\,\mathcal{R}_\parallel + |E_\perp|^2\,\mathcal{R}_\perp}
                     {|E_\parallel|^2 + |E_\perp|^2},
  \qquad
  \mathcal{T} = \frac{|E_\parallel|^2\,\mathcal{T}_\parallel + |E_\perp|^2\,\mathcal{T}_\perp}
                     {|E_\parallel|^2 + |E_\perp|^2}.
\]

\begin{table}[H]
  \centering
  \caption{Resumen de casos especiales para dieléctrico--dieléctrico}
  \label{tab:casos_especiales}
  \begin{tabularx}{0.98\linewidth}{@{}L{5.2cm}Y Y@{}}
    \toprule
    \textbf{Caso} & $\mathcal{R}_\parallel$ & $\mathcal{R}_\perp$ \\
    \midrule
    Incidencia normal ($\theta_i = 0$)
      & $\left(\dfrac{n_1-n_2}{n_1+n_2}\right)^2$
      & $\left(\dfrac{n_1-n_2}{n_1+n_2}\right)^2$ \\
    Ángulo de Brewster ($\theta_i = \theta_B$)
      & $0$
      & $\sin^2(\theta_B - \theta_t)/\sin^2(\theta_B + \theta_t)$ \\
    Reflexión total ($\theta_i > \theta_c$, $n_1>n_2$)
      & $1$
      & $1$ \\
    \bottomrule
  \end{tabularx}
\end{table}

% ══════════════════════════════════════════════════════════════════════════
\subsection{Dieléctrico--Medio Absorbente}
% ══════════════════════════════════════════════════════════════════════════

El Medio 1 es un dieléctrico transparente con índice $n_1 \in \R$.
El Medio 2 es absorbente, descrito por un \textbf{índice complejo}:
\[
  n_c = n_2 - i\kappa, \qquad n_2, \kappa > 0.
\]

Las condiciones de contorno se mantienen formalmente iguales, pero ahora
$k_{tcz} \in \C$. Se escribe:
\[
  \ve{k}_{tc} = \ve{k}_t - i\,\ve{a}, \qquad \ve{k}_t \perp \ve{a}.
\]

Las condiciones de fase tangencial en $z = 0$ imponen de nuevo:
\[
  k_{tx} = k_{ix} = \frac{\omega}{c}n_1\sin\theta_i \in \R,
  \qquad a_x = 0.
\]

Por tanto $\ve{a} = (0, 0, a_z)$ con $a_z > 0$ (atenuación en la dirección
de penetración).

\subsubsection{Campo transmitido (onda inhomogénea)}

El campo transmitido en el medio 2 es una \textbf{onda inhomogénea}:
\begin{equation}
  \ve{E}_t = \ve{T}\,e^{-a_z z}\,e^{-ik_{tx} x}\,e^{i\omega t}.
  \label{eq:transmitida_metal}
\end{equation}

Los planos de fase constante (perpendiculares a $\ve{k}_t = (k_{tx}, 0, 0)$)
son \textbf{verticales} (paralelos a $z$), mientras que los planos de
amplitud constante (perpendiculares a $\ve{a} = (0, 0, a_z)$) son
\textbf{horizontales} (paralelos a la interfaz). Estos dos conjuntos de
planos no coinciden, de ahí el término \emph{inhomogénea}.

La relación de dispersión en el medio 2 determina $k_{tz}$ a partir de:
\[
  k_{tx}^2 + k_{tz}^2 = \left(\frac{\omega}{c}\right)^2 n_c^2,
\]
con $k_{tx}$ real ya fijado, por lo que $k_{tz}$ resulta complejo en general.

\subsubsection{Campo reflejado (coeficientes de Fresnel complejos)}

La onda reflejada queda íntegramente en el Medio 1 (transparente). Aplicando
las condiciones de contorno con $k_{tcz}$ complejo se obtienen los
coeficientes de Fresnel generalizados:

\begin{mdframed}[style=resultado]
\textbf{Coeficientes de Fresnel para medio absorbente:}
\begin{align}
  r_\parallel &= \frac{n_c^2\, k_{iz} - n_1^2\, k_{tcz}}
                     {n_c^2\, k_{iz} + n_1^2\, k_{tcz}}
               = \rho_\parallel\, e^{i\Phi_\parallel},
  \label{eq:rp_metal} \\
  r_\perp     &= \frac{k_{iz} - k_{tcz}}
                     {k_{iz} + k_{tcz}}
               = \rho_\perp\, e^{i\Phi_\perp},
  \label{eq:rs_metal}
\end{align}
donde $k_{iz} = (\omega/c)n_1\cos\theta_i$ es real y
$k_{tcz}$ es en general complejo.
\end{mdframed}

El campo reflejado resultante puede escribirse:
\[
  \ve{R} = \begin{pmatrix} \rho_\parallel e^{i\Phi_\parallel} E_\parallel \\[4pt]
                            \rho_\perp e^{i\Phi_\perp} E_\perp \end{pmatrix}
         = \rho_\parallel e^{i\Phi_\parallel}
           \begin{pmatrix} E_\parallel \\[4pt]
                           \dfrac{\rho_\perp}{\rho_\parallel}
                           e^{i(\Phi_\perp - \Phi_\parallel)} E_\perp
           \end{pmatrix}.
\]

\begin{observacion}
A diferencia del caso dieléctrico--dieléctrico por debajo del ángulo crítico,
los módulos $\rho_\parallel$, $\rho_\perp$ pueden ser distintos, y los
desfases $\Phi_\parallel$, $\Phi_\perp$ son en general no nulos y diferentes
entre sí. Esto modifica el estado de polarización del haz reflejado.
\end{observacion}

\subsubsection{Caso particular: incidencia normal sobre un metal}

Para $\theta_i = 0$ los vectores de onda son verticales:
$\ve{k}_i = (0, 0, k_{iz})$, $\ve{a} = (0, 0, a_z)$.
Los coeficientes de Fresnel (\ref{eq:rs_metal}) se reducen a:

\begin{proof}
Para incidencia normal:
\[
  \theta_i = 0.
\]
Luego:
\[
  k_{iz} = \frac{\omega}{c}n_1,
  \qquad
  k_{tcz} = \frac{\omega}{c}n_c
          = \frac{\omega}{c}(n_2 - i\kappa).
\]
Sustituyendo:
\[
  r_\perp = \frac{n_1 - (n_2 - i\kappa)}{n_1 + (n_2 - i\kappa)}
           = \frac{(n_1 - n_2) + i\kappa}{(n_1 + n_2) - i\kappa}.
\]
Multiplicando por el conjugado del denominador:
\[
  |r_\perp|^2
  =\frac{(n_1-n_2)^2+\kappa^2}{(n_1+n_2)^2+\kappa^2}.
\]
\end{proof}

\begin{mdframed}[style=resultado]
\textbf{Reflectancia en incidencia normal sobre metal:}
\begin{equation}
  r_\perp = \frac{(n_1 - n_2) + i\kappa}{(n_1 + n_2) - i\kappa},
  \qquad
  \mathcal{R} = |r_\perp|^2 = \frac{(n_1 - n_2)^2 + \kappa^2}
                                    {(n_1 + n_2)^2 + \kappa^2}.
  \label{eq:R_metal_normal}
\end{equation}
\end{mdframed}

En el límite de metal muy absorbente ($\kappa \gg n_2$):
\[
  r_\perp \approx \frac{n_1 + i\kappa}{n_1 - i\kappa}
  = e^{2i\arctan(\kappa/n_1)},
  \qquad
  \mathcal{R} \approx 1 - \frac{4n_1}{n_1^2/\kappa + \kappa} \to 1.
\]

Los metales buenos conductores son por tanto casi perfectamente reflectantes
en el rango óptico, con $\mathcal{R} \lesssim 1$.

\begin{observacion}
Para el aluminio ($n_2 \approx 1.4$, $\kappa \approx 7.5$ a $\lambda =
550\,\text{nm}$), la fórmula~\eqref{eq:R_metal_normal} con $n_1 = 1$ da
$\mathcal{R} \approx 92\%$. La plata tiene la reflectividad más alta de
todos los metales en el visible ($\mathcal{R} \approx 96{-}98\%$), lo que
la hace el material preferido para espejos astronómicos de alta
reflectividad. El oro tiene $\mathcal{R} < 50\%$ en el azul y violeta
(por eso su color característico), pero $> 98\%$ en el infrarrojo.
\end{observacion}

\begin{table}[H]
  \centering
  \caption{Comparación entre los dos tipos de interfaz estudiados}
  \label{tab:comparacion_interfaces}
  \begin{tabularx}{0.98\linewidth}{@{}L{2.9cm}Y Y@{}}
    \toprule
    \textbf{Magnitud} & \textbf{Diél.--Diél. ($n_1, n_2 \in \R$)}
                      & \textbf{Diél.--Absorbente ($n_c \in \C$)} \\
    \midrule
    Ley de Snell & $n_1\sin\theta_i = n_2\sin\theta_t$ & Válida con $n_c$ \\
    $r_j$ & Real (debajo de $\theta_c$) & Complejo en general \\
    $|r_j|$ & $\leq 1$, $= 1$ solo si $\theta_i > \theta_c$ & $\leq 1$ \\
    Desfase & 0 o $\pi$ (debajo de $\theta_c$) & Continuo ($\Phi_j \neq 0, \pi$) \\
    Onda transmitida & Homogénea (planos coinciden) & Inhomogénea \\
    Reflexión total & Si $n_1 > n_2$ y $\theta_i > \theta_c$ & No existe \\
    \bottomrule
  \end{tabularx}
\end{table}

\section{Tema 4.2: Introducción a los Medios Anisótropos}

\subsection{Introducción}

Hasta ahora se han tratado medios isótropos, en los que el tensor de
permitividad es escalar: la respuesta del material al campo eléctrico es la
misma en cualquier dirección. Sin embargo, muchos materiales cristalinos
—cuarzo, calcita, niobato de litio, mica— tienen una estructura atómica
que rompe esta simetría. En ellos, la relación constitutiva toma la forma
matricial:
\begin{equation}
  \ve{D} = \boldsymbol{\varepsilon}\,\ve{E},
  \qquad
  \boldsymbol{\varepsilon} =
  \begin{pmatrix}
    \varepsilon_x & 0 & 0 \\
    0 & \varepsilon_y & 0 \\
    0 & 0 & \varepsilon_z
  \end{pmatrix}
  \quad\text{(ejes dieléctricos principales),}
\end{equation}
de modo que $\ve{D}$ y $\ve{E}$ ya no son paralelos en general. Esta
propiedad —la \emph{anisotropía óptica}— tiene consecuencias profundas y
de gran utilidad tecnológica: la luz de diferentes polarizaciones se propaga
con velocidades distintas (\emph{birrefringencia}), el plano de polarización
puede rotarse (\emph{actividad óptica}), y se pueden construir dispositivos
para controlar el estado de polarización de la luz con precisión arbitraria.

En el modelo microscópico de Lorentz, la anisotropía aparece cuando la
fuerza recuperadora sobre el electrón ligado es diferente según la dirección
de desplazamiento: la fuerza no es proporcional a $\ve{r}$ sino que actúa
mediante un tensor de rigidez $\hat{c}$,
\begin{equation}
  \ve{F}_{\text{rec}} = -\hat{c}\,\ve{r},
\end{equation}
lo que conduce a frecuencias de resonancia $\omega_{0,x}$, $\omega_{0,y}$,
$\omega_{0,z}$ distintas para cada dirección. En consecuencia, la
susceptibilidad —y la permitividad— son tensoriales incluso a la frecuencia
óptica.

En términos físicos, esto implica que ``índice de refracción'' deja de ser un
único escalar y pasa a depender de la dirección de propagación y de la
polarización excitada.

\begin{mdframed}[style=resultado]
\textbf{Clasificación por simetría del tensor $\boldsymbol{\varepsilon}$:}
\begin{itemize}
  \item $\varepsilon_x=\varepsilon_y=\varepsilon_z$: medio \textbf{isótropo}
        (cúbico o amorfo). Un único índice de refracción.
  \item $\varepsilon_x=\varepsilon_y\neq\varepsilon_z$: medio \textbf{uniáxico}
        (simetría trigonal, hexagonal o tetragonal). Dos índices: ordinario $n_o$
        y extraordinario $n_e$.
  \item $\varepsilon_x\neq\varepsilon_y\neq\varepsilon_z$: medio \textbf{biáxico}
        (simetría ortorrómbica, monoclínica o triclínica). Tres índices principales.
\end{itemize}
\end{mdframed}

\begin{nota}
La calcita ($n_o=1.658$, $n_e=1.486$ para $\lambda=589\,\text{nm}$) es el
cristal uniáxico negativo más estudiado históricamente; el cuarzo
($n_o=1.544$, $n_e=1.553$) es positivo. Los cristales biáxicos como la mica
muscovita se emplean en láminas retardadoras de alta precisión.
\end{nota}

\subsection{Medios Uniáxicos}

\subsubsection{Conceptos básicos}

En un medio uniáxico existe una única dirección privilegiada —el
\emph{eje óptico}— a lo largo de la cual la simetría de rotación se
conserva y el medio se comporta ópticamente como isótropo. Para toda
otra dirección de propagación, el campo eléctrico se descompone en dos
modos con índices de refracción distintos:

\begin{itemize}
  \item \textbf{Eje óptico}: dirección en la que el medio se comporta como
        isótropo; a lo largo de ella, ambas polarizaciones se propagan con
        el mismo índice $n_o$ y no hay separación de haces.
  \item \textbf{Sección principal}: plano que contiene el eje óptico y el
        vector de onda $\ve{k}$; sirve de referencia para descomponer el
        campo.
  \item \textbf{Componente ordinaria} ($o$): componente de $\ve{E}$
        perpendicular a la sección principal. Se propaga con índice $n_o$
        independientemente de la dirección de $\ve{k}$ (de ahí el nombre
        \emph{ordinario}: su velocidad de fase no depende de la dirección
        de propagación).
  \item \textbf{Componente extraordinaria} ($e$): componente de $\ve{E}$
        contenida en la sección principal. Su índice de fase efectivo varía
        con el ángulo $\theta$ entre $\ve{k}$ y el eje óptico:
        \[
          \frac{1}{n_e^2(\theta)}
          = \frac{\cos^2\theta}{n_o^2}+\frac{\sin^2\theta}{n_e^2}.
        \]
        Solo en la dirección perpendicular al eje ($\theta=90^\circ$) alcanza
        el valor extremo $n_e$.
\end{itemize}

En un medio uniáxico los parámetros fundamentales son:
\[
  \varepsilon_x=\varepsilon_y=\varepsilon_o,\qquad
  \varepsilon_z=\varepsilon_e,
\]
\[
  n_o=\sqrt{\varepsilon_o},\quad
  n_e=\sqrt{\varepsilon_e},\quad
  v_o=\frac{c}{n_o},\quad
  v_e=\frac{c}{n_e},\quad
  k_o=\frac{2\pi}{\lambda}n_o,\quad
  k_e=\frac{2\pi}{\lambda}n_e.
\]

\begin{figure}[H]
  \centering
  \begin{tikzpicture}[>=Stealth, scale=1.1]
    \draw[->, thick] (0,0) -- (0,2.2) node[above] {$E_y$};
    \draw[->, thick] (0,0) -- (2.4,0) node[right] {$E_x$};
    \draw[thick, dashed] (0,0) -- (1.8,1.4);
    \draw[very thick] (2.8,-0.2) -- (2.8,2.2);
    \node[align=left, anchor=west] at (3.3,1.4)
      {\footnotesize Línea discontinua: sección principal\\
       \footnotesize Línea vertical: eje óptico\\
       \footnotesize $\ve{E}=\ve{E}_o+\ve{E}_e$};
  \end{tikzpicture}
  \caption{Descomposición del campo en componentes ordinaria y extraordinaria
           respecto a la sección principal y al eje óptico.}
  \label{fig:uniaxico_componentes}
\end{figure}

\subsubsection{Medio uniáxico positivo y negativo}

\begin{itemize}
  \item \textbf{Uniáxico positivo}: $n_o < n_e \Rightarrow v_o > v_e$;
        el haz extraordinario es más lento que el ordinario. Ejemplo:
        cuarzo ($n_o = 1.544$, $n_e = 1.553$ a $\lambda = 589\,\text{nm}$).
  \item \textbf{Uniáxico negativo}: $n_o > n_e \Rightarrow v_o < v_e$;
        el haz extraordinario es más rápido que el ordinario. Ejemplo:
        calcita ($n_o = 1.658$, $n_e = 1.486$ a $\lambda = 589\,\text{nm}$),
        con una birrefringencia $|\Delta n| = 0.172$ excepcionalmente alta.
\end{itemize}

\begin{nota}
La birrefringencia $\Delta n = n_e - n_o$ es positiva para cristales
positivos y negativa para los negativos. Su valor determina directamente
el espesor de lámina retardadora necesario para lograr un desfase dado
$\delta = 2\pi|\Delta n|d/\lambda$. A mayor $|\Delta n|$, menor espesor
necesario (útil para miniaturizar dispositivos) pero también mayor
sensibilidad a la temperatura.
\end{nota}

\subsection{Birrefringencia o Doble Refracción}

La birrefringencia es la propiedad por la que un cristal anisótropo separa un
haz de luz no polarizado en dos haces con polarizaciones ortogonales que se
propagan en direcciones (ligeramente) distintas. Fue observada por primera vez
en la calcita por Erasmus Bartholinus (1669) y explicada mecánicamente por
Huygens (1690) mediante la construcción de superficies de ondas de formas
distintas para cada polarización: esférica para el rayo ordinario y
elipsoidal para el rayo extraordinario.

Cuando una onda incide sobre un cristal uniáxico, la ley de Snell se aplica
por separado a cada componente, pero con índices distintos. Para una interfaz
con medio incidente de índice $n_1$:
\begin{equation}
  n_1\sin\theta_i = n_o\sin\theta_o,
  \qquad
  n_1\sin\theta_i = n_e(\theta)\sin\theta_e.
  \label{eq:snell_birrefringencia}
\end{equation}

Si $n_o\neq n_e$, entonces $\theta_o\neq\theta_e$: aparecen dos haces
transmitidos con polarizaciones ortogonales. Este fenómeno es la base de los
polarizadores de prisma (Nicol, Glan-Taylor, Wollaston), que explotan la
diferencia de ángulo para separar espacialmente las dos polarizaciones.

\begin{figure}[H]
  \centering
  \begin{tikzpicture}[>=Stealth, scale=1.05]
    \draw[thick] (-3,0) -- (3,0);
    \node[left] at (-2.9,0.25) {$n_1$};
    \node[left] at (-2.9,-0.25) {cristal uniáxico};
    \draw[->, thick] (0,0) -- (0,2) node[above] {$\uvec{n}$};

    \draw[->, very thick] (-2.0,1.8) -- (0,0) node[midway, above left] {$\ve{k}_i$};
    \draw[->, very thick] (0,0) -- (1.5,-1.4) node[midway, below right] {$\ve{k}_o$};
    \draw[->, very thick, dashed] (0,0) -- (2.3,-1.0) node[midway, above right] {$\ve{k}_e$};
    \node[align=left, anchor=west] at (3.3,-0.3)
      {\footnotesize Haz ordinario: $\ve{k}_o$\\
       \footnotesize Haz extraordinario: $\ve{k}_e$\\
       \footnotesize Resultado: $\theta_o\neq\theta_e$};
  \end{tikzpicture}
  \caption{Birrefringencia: separación angular entre el haz ordinario y el
           extraordinario.}
  \label{fig:birrefringencia}
\end{figure}

\subsection{Láminas Retardadoras}

Una lámina retardadora es una lámina de cristal uniáxico de espesor $d$
cortada de modo que el eje óptico yace en su plano (ejes neutros $x$ e $y$
del laboratorio). Cuando una onda plana la atraviesa con incidencia normal,
las componentes ordinaria y extraordinaria recorren el mismo camino
geométrico $d$ pero con índices $n_o$ y $n_e$ distintos. Acumulan así una
diferencia de fase:
\begin{equation}
  \delta = \frac{2\pi}{\lambda}(n_e-n_o)d = k_0\,\Delta n\,d,
  \label{eq:retardador_delta}
\end{equation}
donde $\Delta n = n_e - n_o$ es la birrefringencia de la lámina. El campo
de salida es:
\begin{equation}
  \ve{E}_s
  =
  e^{i\omega t}e^{-ik_0 n_o d}
  \begin{pmatrix}
    E_x\\
    E_y\,e^{-i\delta}
  \end{pmatrix}.
\end{equation}
La fase global $e^{-ik_0 n_o d}$ es irrelevante para el estado de
polarización; lo que importa es el \emph{retardo} $\delta$.

\begin{nota}
Para una longitud de onda dada, ajustar el espesor $d$ permite conseguir
cualquier retardo $\delta$ deseado. En la práctica, láminas de orden cero
($\delta = \pi/2$ obtenido con el mínimo $d$ posible) son preferibles porque
su retardo varía menos con la temperatura y con la longitud de onda.
\end{nota}

El desfase depende de la birrefringencia $\Delta n=n_e-n_o$ y del espesor $d$.

\subsubsection{Lámina de cuarto de onda}

Para $\delta=\pi/2 + m\pi$ (casos $m = 0, 1, 2, \ldots$, siendo el de
orden cero el más utilizado):
\[
  L_{\lambda/4}
  =
  \begin{pmatrix}
    1 & 0\\
    0 & i
  \end{pmatrix}
  \quad\text{(equivalente salvo fase global).}
\]
Aplicaciones típicas:
\begin{itemize}
  \item Lineal a $45^\circ \to$ circular (la lámina introduce un retardo
        de $90^\circ$ entre las componentes igualmente amplias, generando
        la trayectoria circular).
  \item Circular $\to$ lineal (inversión de la operación anterior).
  \item Permite determinar la helicidad (dextro o levo) de la luz
        circularmente polarizada: al pasar por $L_{\lambda/4}$, la luz
        circular se convierte en lineal y el ángulo resultante indica la
        helicidad.
\end{itemize}

\subsubsection{Lámina de media onda}

Para $\delta=\pi + 2m\pi$:
\[
  L_{\lambda/2}
  =
  \begin{pmatrix}
    1 & 0\\
    0 & -1
  \end{pmatrix}.
\]
Rota la polarización lineal de ángulo $\theta$ a $-\theta$ respecto a su
eje neutro (o, equivalentemente, rota el plano de polarización en $2\theta$
si se gira la lámina un ángulo $\theta$). Invierte la helicidad de la
polarización circular.

\begin{observacion}
En láminas ideales no hay absorción: la intensidad total se conserva
($I_s=I_i$). Además, la combinación de dos láminas de media onda giradas
a ángulos $\alpha$ y $\beta$ equivale a un rotador de polarización con
ángulo de rotación $2(\beta - \alpha)$. Esto es la base de los
\emph{rotadores de Fresnel} en sistemas láser de alta potencia, donde la
absorción de un polaroide sería inaceptable.
\end{observacion}

\subsection{Matrices de Jones}

El formalismo de Jones (R. C. Jones, 1941) es el marco algebraico para
describir la transformación del estado de polarización de luz monocromática
al pasar por elementos ópticos. Cada elemento se representa por una matriz
$2\times 2$ de números complejos que actúa sobre el vector de Jones de la
onda:
\[
  \ve{E}_s = L\,\ve{E}_i,
\]
donde $L$ es la \emph{matriz de Jones} del elemento. El estado de
polarización resultante al pasar por una secuencia de $N$ elementos es:
\[
  \ve{E}_s = L_N \cdots L_2\, L_1\,\ve{E}_i.
\]
El orden importa (producto de matrices no conmutativo en general). Las
propiedades físicas relevantes —elipse de polarización, intensidad, desfase—
se extraen directamente de $\ve{E}_s$.

\begin{nota}
El formalismo de Jones es equivalente al formalismo de Stokes--Mueller para
luz completamente polarizada. El formalismo de Stokes--Mueller es más general
porque también maneja luz parcialmente polarizada, pero es más complicado
(matrices $4\times 4$ reales).
\end{nota}

Para un retardador general alineado con ejes neutros:
\[
  L(\delta)=
  \begin{pmatrix}
    1 & 0\\
    0 & e^{-i\delta}
  \end{pmatrix}.
\]
Si la lámina se gira un ángulo $\alpha$:
\begin{equation}
  L_\alpha(\delta)=R^{-1}(\alpha)\,L(\delta)\,R(\alpha),
  \qquad
  R(\alpha)=
  \begin{pmatrix}
    \cos\alpha & -\sin\alpha\\
    \sin\alpha & \cos\alpha
  \end{pmatrix}.
  \label{eq:jones_rotada}
\end{equation}

\subsection{Polarizadores}

Un polarizador ideal es un elemento que transmite solo la componente del
campo eléctrico paralela a su eje de transmisión. Físicamente se realiza
mediante:
\begin{itemize}
  \item \emph{Polaroides}: películas con moléculas alineadas que absorben
        selectivamente una polarización (dicroísmo).
  \item \emph{Prismas de Glan o Nicol}: explotan la birrefringencia y la
        reflexión total interna para separar espacialmente las dos
        polarizaciones.
  \item \emph{Rejillas metálicas}: en el infrarrojo, la componente paralela
        a los conductores queda reflejada; la perpendicular se transmite.
\end{itemize}

Si el eje del polarizador forma ángulo $\alpha$ con $x$, su matriz de Jones es:
\begin{equation}
  P_\alpha=
  \begin{pmatrix}
    \cos^2\alpha & \cos\alpha\sin\alpha\\
    \cos\alpha\sin\alpha & \sin^2\alpha
  \end{pmatrix}.
  \label{eq:polarizador_jones}
\end{equation}

Para una onda lineal incidente
$\ve{E}_i=E_0(\cos\theta,\sin\theta)^T$:
\[
  \ve{E}_s = P_\alpha \ve{E}_i
  = E_0\cos(\theta-\alpha)
    \begin{pmatrix}\cos\alpha\\ \sin\alpha\end{pmatrix}.
\]
Por tanto:
\begin{equation}
  I_s = I_0\cos^2(\theta-\alpha).
  \label{eq:malus}
\end{equation}
Esta es la \textbf{ley de Malus} (Étienne-Louis Malus, 1809): la intensidad
transmitida por un polarizador varía con el coseno al cuadrado del ángulo
entre la polarización incidente y el eje del polarizador. Su descubrimiento
experimental (en la reflexión de luz solar en las ventanas del Palais du
Luxembourg) fue la primera evidencia directa de la naturaleza transversal de
la luz.

Si la entrada es circular
$\ve{E}_i=\frac{E_0}{\sqrt{2}}(1,\pm i)^T$, entonces
$I_s=I_0/2$, independiente de $\alpha$: la intensidad transmitida no depende
de la orientación del polarizador, reflejo de la simetría azimutal de la
polarización circular.

\section{Tema 5: Difracción}

\subsection{Introducción}

En óptica geométrica, la luz se modela como rayos que viajan en línea recta y
la sombra de un obstáculo tiene bordes perfectamente nítidos. Sin embargo, esta
aproximación falla cuando la longitud de onda es comparable al tamaño de la
abertura u obstáculo ($a \sim \lambda$): aparecen efectos ondulatorios
irreducibles, colectivamente denominados \emph{difracción}, en los que la
energía se redistribuye en un patrón espacial de máximos y mínimos de
intensidad.

La difracción es una consecuencia directa del principio de Huygens-Fresnel y
de la naturaleza ondulatoria de la luz. Se distinguen dos regímenes principales:

\begin{itemize}
  \item \textbf{Difracción de Fresnel (campo cercano):} la curvatura del frente
        de onda en el plano de observación es relevante; la geometría es más
        compleja. Se observa a distancias intermedias de la abertura
        ($N_F = a^2/\lambda z \gtrsim 1$).
  \item \textbf{Difracción de Fraunhofer (campo lejano):} los frentes de onda
        incidente y difractado pueden aproximarse como planos; el patrón de
        intensidad observado es, matemáticamente, la transformada de Fourier del
        campo en la abertura. Se obtiene para $z \gg a^2/\lambda$ (o, en
        la práctica, en el plano focal de una lente).
\end{itemize}

La frontera entre ambos regímenes no es nominal: determina qué términos de fase
se conservan en el núcleo integral y, por tanto, qué aproximación es
matemáticamente consistente.

\subsection{Principio de Huygens-Fresnel}

El principio de Huygens (1678) establece que cada punto de un frente de onda puede
considerarse fuente de ondas secundarias esféricas cuya envolvente constituye el
siguiente frente. Fresnel (1818) reformuló este principio de forma cuantitativa
al añadir la suma coherente de estas fuentes secundarias —con amplitudes y fases—
para calcular el campo en cualquier punto. La formulación matemática rigurosa fue
dada por Kirchhoff (1882) a partir de las ecuaciones de onda, confirmando el
principio heurístico de Huygens-Fresnel como una consecuencia exacta de la teoría
ondulatoria escalar.

Cada punto de la abertura actúa como una fuente secundaria coherente. En la aproximación
escalar, el campo complejo en un punto de observación $P$ se escribe:
\begin{equation}
  U(P)=\frac{1}{i\lambda}\iint_A
  U_0(x_0,y_0)\,\frac{e^{-ikr}}{r}\,\cos\theta\,dx_0\,dy_0.
  \label{eq:hf_integral_t5}
\end{equation}
donde $U$ es la amplitud compleja del campo óptico, $(x_0,y_0)$ las coordenadas en el plano de la abertura, $r$ la distancia entre el punto emisor y el punto de observación $P$, y $\theta$ el ángulo de oblicuidad respecto a la normal.

\begin{itemize}
  \item $U_0(x_0,y_0)$: amplitud en la abertura.
  \item $\frac{e^{-ikr}}{r}$: propagación esférica desde cada punto.
  \item $\cos\theta$: factor de oblicuidad.
\end{itemize}

\begin{figure}[ht]
  \centering
  \begin{tikzpicture}[scale=0.95]
    \draw[thick] (-3.1,-1.9) -- (-3.1,1.9);
    \draw[thick] (3.2,-1.9) -- (3.2,1.9);
    \filldraw[black] (-3.1,0.4) circle (1.4pt);
    \filldraw[black] (3.2,0.9) circle (1.4pt);
    \draw[->,thick] (-2.95,0.45) -- (3.0,0.86);
    \draw[->] (-4.4,-1.5) -- (-4.4,1.5) node[above,opticslabel] {$y_0$};
    \draw[->] (-4.4,-1.5) -- (-2.7,-1.5) node[right,opticslabel] {$x_0$};
    \draw[->] (4.5,-1.5) -- (4.5,1.5) node[above,opticslabel] {$y$};
    \draw[->] (4.5,-1.5) -- (2.8,-1.5) node[left,opticslabel] {$x$};
    \draw[dashed] (-2.5,-1.9) -- (2.6,-1.9) node[midway,below,opticslabel] {$z$};
    \node[align=left, anchor=west, opticslabel] at (4.95,1.25)
      {$P_0(x_0,y_0)$: punto fuente en abertura\\
       $P(x,y)$: punto de observación\\
       Flecha oblicua: distancia $r$\\
       Línea izquierda: abertura\\
       Línea derecha: pantalla};
  \end{tikzpicture}
  \caption{Geometría básica de difracción por una abertura.}
  \label{fig:geometria_difraccion_t5}
\end{figure}

\subsection{Aproximación de Fresnel (campo cercano)}

Para propagación paraxial:
\[
  r=\sqrt{(x-x_0)^2+(y-y_0)^2+z^2}
  \approx
  z+\frac{(x-x_0)^2+(y-y_0)^2}{2z}.
\]
Sustituyendo en \eqref{eq:hf_integral_t5}:
\begin{equation}
  U(x,y)=\frac{e^{-ikz}}{i\lambda z}
  \iint_A U_0(x_0,y_0)
  \exp\!\left[-i\frac{k}{2z}\left((x-x_0)^2+(y-y_0)^2\right)\right]
  dx_0\,dy_0.
  \label{eq:fresnel_t5}
\end{equation}
donde $z$ es la distancia abertura--pantalla y $(x,y)$ son coordenadas en el
plano de observación.

\subsection{Aproximación de Fraunhofer (campo lejano)}

Si $x_0,y_0\ll z$, en el término cuadrático se desprecia $x_0^2,y_0^2$ y queda:
\begin{equation}
  U(x,y)\approx
  \frac{e^{-ikz}}{i\lambda z}\,
  e^{-i\frac{k}{2z}(x^2+y^2)}
  \iint_{-\infty}^{+\infty}
  t(x_0,y_0)\,U_0(x_0,y_0)\,
  e^{\,i\frac{k}{z}(xx_0+yy_0)}\,dx_0\,dy_0.
  \label{eq:fraunhofer_t5}
\end{equation}
Es decir, el campo en observación es la transformada de Fourier del campo
transmitido por la abertura.

Usando una lente y observando en su plano focal imagen ($z=f'$):
\begin{equation}
  U(x',y')=
  \frac{e^{-ikf'}}{i\lambda f'}
  e^{-i\frac{k}{2f'}(x'^2+y'^2)}
  \iint
  t(x_0,y_0)U_0(x_0,y_0)
  e^{\,i\frac{k}{f'}(x'x_0+y'y_0)}dx_0dy_0.
  \label{eq:fraunhofer_lente_t5}
\end{equation}
donde $f'$ es la focal imagen de la lente y $(x',y')$ son coordenadas en su
plano focal.

\subsubsection{Criterio de validez y número de Fresnel}

Una forma práctica de separar regímenes de propagación es mediante el
\textbf{número de Fresnel}:
\begin{equation}
  N_F=\frac{a^2}{\lambda z},
  \label{eq:fresnel_number_t5}
\end{equation}
donde $a$ es una escala transversal característica de la abertura.
\begin{itemize}
  \item $N_F\gtrsim 1$: régimen de Fresnel (campo cercano).
  \item $N_F\ll 1$: régimen de Fraunhofer (campo lejano), equivalente a
  $z\gg a^2/\lambda$.
\end{itemize}
Esta distinción ayuda a escoger directamente el modelo integral adecuado.

\subsection{Principio de Babinet}

El principio de Babinet establece una simetría elegante entre la difracción
por una abertura y por su obstáculo complementario. Un obstáculo es
\emph{complementario} de una abertura si ocupa exactamente el espacio que
la abertura deja libre, es decir:
\begin{equation}
  t_a(x_0,y_0)+t_o(x_0,y_0)=1.
  \label{eq:babinet_transmision_t5}
\end{equation}
Por linealidad de la integral de Huygens-Fresnel, en el régimen de Fraunhofer:
\[
  U_a(x,y)+U_o(x,y)=U_{\text{inc}}(x,y).
\]
Para una onda plana incidente ($U_{\text{inc}} = \mathrm{cte}$ en el plano de
observación, salvo en el eje), los patrones de intensidad difractada por la
abertura y por el obstáculo son idénticos fuera del eje central. Esto implica,
por ejemplo, que un disco circular opaco y un agujero circular del mismo
radio producen exactamente el mismo patrón de difracción lejos del eje (salvo
el punto central del eje óptico).

Este principio tiene aplicaciones en la teoría de antenas, en la caracterización
de partículas en suspensión y en holografía.

\subsection{Difracción por abertura rectangular}

Para una abertura $a\times b$:
\[
  t(x_0,y_0)=
  \begin{cases}
    1, & |x_0|\le a/2,\ |y_0|\le b/2,\\
    0, & \text{resto.}
  \end{cases}
\]
En Fraunhofer:
\begin{equation}
  U(x,y)=C\;ab\;
  \operatorname{sinc}\!\left(\frac{kax}{2z}\right)
  \operatorname{sinc}\!\left(\frac{kby}{2z}\right),
  \qquad
  \operatorname{sinc}(u)=\frac{\sin u}{u}.
  \label{eq:rect_amplitud_t5}
\end{equation}
donde $C$ es una constante compleja global (fase y factor de escala).
Luego la intensidad:
\begin{equation}
  I(x,y)=I_0\,
  \operatorname{sinc}^2\!\left(\frac{kax}{2z}\right)
  \operatorname{sinc}^2\!\left(\frac{kby}{2z}\right).
  \label{eq:rect_intensidad_t5}
\end{equation}
Mínimos:
\begin{equation}
  x_m=\frac{m\lambda z}{a},\qquad
  y_m=\frac{m\lambda z}{b},
  \qquad m\in\mathbb{Z}\setminus\{0\}.
  \label{eq:rect_minimos_t5}
\end{equation}

\subsubsection{Caso límite: rendija simple}

Si $b\to\infty$:
\begin{equation}
  I(x)=I_{\max}\,
  \operatorname{sinc}^2\!\left(\frac{kax}{2z}\right).
  \label{eq:rendija_simple_t5}
\end{equation}
Aquí $I_{\max}$ es la intensidad del máximo central.
Con incidencia oblicua $\alpha$, el patrón se desplaza:
\begin{equation}
  I(x)=I_{\max}\,
  \operatorname{sinc}^2\!\left[\frac{ka}{2z}\left(x-z\sin\alpha\right)\right],
  \qquad
  x_m=\frac{m\lambda z}{a}+z\sin\alpha.
  \label{eq:rendija_oblicua_t5}
\end{equation}

\subsection{Difracción por abertura circular}

Para abertura circular (radio $a$), la simetría azimutal del problema hace que el
patrón en Fraunhofer solo dependa del radio $r$ en el plano de observación. La
transformada de Fourier bidimensional de la función pupila circular —una función
de sombrero de copa (\emph{top-hat} circular)— es la transformada de Hankel de
orden cero, que resulta proporcional a $J_1(\xi)/\xi$, donde $J_1$ es la función
de Bessel de primera especie y primer orden. Definiendo:
\[
  \xi=\frac{k\,a\,r}{f'},
\]
la amplitud e intensidad toman la forma:
\begin{equation}
  U(r)=U(0)\,\frac{2J_1(\xi)}{\xi},
  \qquad
  I(r)=I(0)\left[\frac{2J_1(\xi)}{\xi}\right]^2,
  \label{eq:airy_t5}
\end{equation}
Este patrón se conoce como \textbf{patrón de Airy} (George Biddell Airy, 1835), y
la zona central brillante como \emph{disco de Airy}.

El primer cero de $J_1(\xi)$ ocurre en $\xi_1\simeq 3.8317 \approx 1.22\pi$:
\begin{equation}
  r_{\text{Airy}}=\frac{1.22\,\lambda f'}{D},
  \qquad D=2a.
  \label{eq:radio_airy_t5}
\end{equation}

\begin{figure}[ht]
  \centering
  \begin{tikzpicture}[scale=1.0]
    \draw[fill=gray!8,draw=black] (0,0) circle (1.4);
    \draw[fill=gray!20,draw=black] (1.4,0) circle (1.4);
    \filldraw[black] (0,0) circle (1.2pt);
    \filldraw[black] (1.4,0) circle (1.2pt);
    \draw[<->,thick] (0,-1.95) -- (1.4,-1.95)
      node[midway,below=2pt,opticslabel] {$\Delta\theta_{\min}$};
    \node[align=left, anchor=west, opticslabel] at (3.0,1.3)
      {Disco 1 y Disco 2:\\
       máximos de Airy adyacentes\\
       separados en el límite de Rayleigh};
  \end{tikzpicture}
  \caption{Criterio de Rayleigh: separación angular límite entre dos máximos.}
  \label{fig:rayleigh_t5}
\end{figure}

\subsection{Poder de resolución}

Usando el criterio de Rayleigh para una abertura circular:
\begin{equation}
  \Delta\theta_{\min}\simeq\frac{1.22\,\lambda}{D}.
  \label{eq:rayleigh_angular_t5}
\end{equation}
En un plano focal imagen:
\begin{equation}
  d_{\min}\simeq f'\,\Delta\theta_{\min}
  =\frac{1.22\,\lambda f'}{D}.
  \label{eq:rayleigh_lineal_t5}
\end{equation}

\begin{observacion}
El criterio de Rayleigh establece que dos puntos fuente se consideran
\emph{apenas resueltos} cuando el máximo central de uno cae sobre el
primer mínimo del otro. Para el ojo humano con pupila $D \approx 3\,\text{mm}$
y $\lambda = 550\,\text{nm}$, esto da $\Delta\theta_{\min} \approx 0.3'$
(minutos de arco), coherente con la separación de los conos en la fóvea.
Para el Hubble Space Telescope ($D = 2.4\,\text{m}$), el límite de Rayleigh
para luz visible es $\Delta\theta_{\min} \approx 0.05''$ (milisegundos de arco).

El criterio de Sparrow (alternativa al de Rayleigh) define la resolución
como el punto en que el doble máximo se fusiona en un perfil plano, y es
más apropiado para sistemas ruidosos. En microscopía, el criterio de Abbe
($d_{\min} = \lambda/(2\,\text{NA})$ con NA el número de apertura numérica)
es el más utilizado.
\end{observacion}

\section{Tema 6: Interferencia}

\subsection{Introducción y definición}

La interferencia es el fenómeno que se produce cuando dos o más ondas se
superponen en una región del espacio: la intensidad resultante no es la suma
de las intensidades individuales, sino que depende también del desfase relativo.
Esto es una consecuencia directa del principio de superposición (válido para
la ecuación de onda lineal) y de la naturaleza cuadrática de la intensidad
respecto al campo.

Para dos ondas escalares $E_1$ y $E_2$, la intensidad total es:
\begin{equation}
  I \propto \left|E_1+E_2\right|^2
  =|E_1|^2+|E_2|^2+2\,\Re\!\left(E_1E_2^*\right).
  \label{eq:def_interferencia_t6}
\end{equation}
El tercer término, $2\,\Re(E_1 E_2^*)$, es el \emph{término de interferencia}
o \emph{término cruzado}: puede ser positivo (interferencia constructiva) o
negativo (interferencia destructiva). Para dos ondas monocromáticas de
intensidades $I_1, I_2$ y desfase relativo constante $\Delta\phi$:
\begin{equation}
  I=I_1+I_2+2\sqrt{I_1I_2}\cos\Delta\phi.
  \label{eq:intensidad_2ondas_t6}
\end{equation}
La intensidad oscila entre $I_\text{máx} = (\sqrt{I_1}+\sqrt{I_2})^2$
(interferencia totalmente constructiva, $\Delta\phi = 2m\pi$) e
$I_\text{mín} = (\sqrt{I_1}-\sqrt{I_2})^2$ (interferencia totalmente
destructiva, $\Delta\phi = (2m+1)\pi$).

Para dos ondas de la misma intensidad $I_1=I_2=I_0$, estos extremos son
$I_\text{máx} = 4I_0$ e $I_\text{mín} = 0$: la energía se redistribuye
espacialmente, sin violar la conservación energética.

Esta distinción entre superposición lineal de campos y carácter cuadrático de
la intensidad será el hilo conductor de los interferómetros del tema.

\subsection{Condiciones de interferencia estable}

Para observar un patrón de franjas estable en el tiempo —requisito necesario
para cualquier medida interferométrica— deben cumplirse tres condiciones
prácticas:
\begin{itemize}
  \item \textbf{Polarizaciones no ortogonales:} dos ondas con polarizaciones
        ortogonales ($\ve{E}_1 \perp \ve{E}_2$) tienen $\Re(E_1 E_2^*) = 0$
        y no interfieren. Solo interfieren las componentes con la misma
        dirección de polarización.
  \item \textbf{Frecuencias iguales o muy próximas:} si $\omega_1 \neq \omega_2$,
        el desfase $\Delta\phi = (\omega_1-\omega_2)t + \Delta\phi_0$ varía en
        el tiempo, haciendo que las franjas se muevan y se borren al promediar.
  \item \textbf{Desfase relativo estable en el tiempo (coherencia):} las fuentes
        deben ser coherentes entre sí; fuentes independientes tienen diferencias
        de fase aleatorias que promedian a cero el término de interferencia.
\end{itemize}

Las magnitudes que cuantifican la coherencia de una fuente son:
\begin{equation}
  \tau_c:\ \text{tiempo de coherencia},\qquad
  l_c=c\,\tau_c\quad (\text{longitud de coherencia}),\qquad
  \Delta\nu\,\tau_c\sim 1.
  \label{eq:coherencia_t6}
\end{equation}
Para que exista interferencia observable, la diferencia de camino óptico entre
los dos brazos debe ser menor que $l_c$. Un láser monomodo puede tener $l_c$
de varios metros; una lámpara espectral, de unos pocos centímetros.

La \emph{visibilidad de franjas} cuantifica el contraste del patrón:
\begin{equation}
  \mathcal{V}=\frac{I_{\max}-I_{\min}}{I_{\max}+I_{\min}}
  =\frac{2\sqrt{I_1I_2}}{I_1+I_2},
  \label{eq:visibilidad_t6}
\end{equation}
con $0 \leq \mathcal{V} \leq 1$. La visibilidad máxima ($\mathcal{V}=1$) se
obtiene cuando $I_1 = I_2$ y las ondas son perfectamente coherentes.

\subsection{Interferómetro de Young}

El experimento de Young (1801) fue la primera demostración concluyente de la
naturaleza ondulatoria de la luz. En él, una fuente puntual ilumina una
pantalla con dos ranuras separadas una distancia $d$. Las dos ranuras actúan
como fuentes secundarias coherentes (mismo origen, mismo desfase inicial) y
sus ondas se superponen en una segunda pantalla a distancia $D \gg d$.

En la aproximación paraxial, la diferencia de camino óptico entre las dos
trayectorias que llegan al punto $x$ de la pantalla vale
$\Delta \approx d\sin\theta \approx d\,x/D$, por lo que:
\begin{equation}
  I(x)=I_1+I_2+2\sqrt{I_1I_2}
  \cos\!\left(\frac{2\pi}{\lambda}\frac{d\,x}{D}\right).
  \label{eq:young_intensidad_t6}
\end{equation}
Si $I_1=I_2=I_0$:
\begin{equation}
  I(x)=2I_0\!\left[1+\cos\!\left(\frac{2\pi d}{\lambda D}x\right)\right]
  =4I_0\cos^2\!\left(\frac{\pi d}{\lambda D}x\right).
  \label{eq:young_simetrico_t6}
\end{equation}

Posición de máximos y mínimos:
\begin{equation}
  d\sin\theta_m=m\lambda,\qquad
  x_m=\frac{m\lambda D}{d},\qquad
  \Delta x=\frac{\lambda D}{d}.
  \label{eq:young_interfranja_t6}
\end{equation}

\begin{figure}[ht]
  \centering
  \begin{tikzpicture}[scale=0.95]
    \draw[thick] (-3.5,-1.8) -- (-3.5,1.8);
    \draw[thick] (3.6,-1.8) -- (3.6,1.8);
    \filldraw[black] (-3.5,0.6) circle (1.2pt);
    \filldraw[black] (-3.5,-0.6) circle (1.2pt);
    \filldraw[black] (3.6,0.25) circle (1.4pt);
    \draw[->,thick] (-3.45,0.6) -- (3.45,0.25);
    \draw[->,thick] (-3.45,-0.6) -- (3.45,0.25);
    \draw[<->] (-3.9,-0.6) -- (-3.9,0.6);
    \draw[dashed] (-3.2,0) -- (3.4,0);
    \draw[<->] (4.05,0) -- (4.05,0.25);
    \draw[dashed] (-2.8,-1.95) -- (3.0,-1.95);
    \node[align=left, anchor=west, opticslabel] at (4.55,1.35)
      {$O_1,O_2$: fuentes coherentes\\
       $P(x)$: punto en pantalla\\
       Segmentos oblicuos: $r_1,r_2$\\
       Separación de fuentes: $d$\\
       Distancia a pantalla: $D$};
  \end{tikzpicture}
  \caption{Geometría de Young e interfranja $\Delta x=\lambda D/d$.}
  \label{fig:young_t6}
\end{figure}

\subsubsection{Biprisma de Fresnel}

El biprisma de Fresnel es una lente divisora de haz que crea dos imágenes
virtuales coherentes de la misma fuente puntual, funcionando como equivalente
óptico del experimento de Young. Consiste en dos prismas delgados de índice
$n$ y ángulo de vértice $\alpha$ (tipicamente $\alpha \lesssim 1^\circ$) unidos
por su base. La refracción en los dos prismas desplaza el rayo en sentidos
opuestos, haciendo que los haces refractados provengan aparentemente de dos
fuentes virtuales separadas una distancia:
\begin{equation}
  d_{\mathrm{eq}}=(n-1)\,\alpha\,a,
  \label{eq:biprisma_sep_t6}
\end{equation}
donde $a$ es la distancia de la fuente real al biprisma. La interfranja se
obtiene de nuevo con $\Delta x=\lambda D/d_{\mathrm{eq}}$, donde $D$ es la
distancia total fuente-pantalla.

\begin{nota}
La ventaja práctica del biprisma respecto a las dos rendijas de Young es que
no bloquea energía: toda la luz de la fuente contribuye a las franjas,
aumentando la intensidad. Es útil en espectroscopía de precisión y para medir
longitudes de onda de fuentes de baja intensidad.
\end{nota}

\subsection{Doble rendija: difracción e interferencia}

Cuando se iluminan dos rendijas de ancho $a$ y separación entre centros $d$,
el patrón de intensidad en campo lejano resulta de la superposición de dos
efectos: la difracción por cada rendija individual (que da la envolvente de
amplitud, con mínimos en $x = m\lambda z/a$) y la interferencia entre las
dos rendijas (que modula la envolvente con una estructura de coseno, con
máximos en $x = m\lambda z/d$). El resultado:
\begin{equation}
  I(x)\propto
  \operatorname{sinc}^2\!\left(\frac{kax}{2z}\right)\,
  4\cos^2\!\left(\frac{kdx}{2z}\right).
  \label{eq:doble_rendija_t6}
\end{equation}
donde $a$ es el ancho de cada rendija, $d$ la separación entre centros,
$z$ la distancia al plano de observación y $k=2\pi/\lambda$.
El primer factor es la envolvente de difracción (debida al ancho finito de
la rendija); el segundo, la modulación interferencial (debida a la
separación entre rendijas).

\begin{nota}
Cuando la razón $d/a$ es un entero $m$, el $m$-ésimo máximo de interferencia
cae exactamente sobre el primer mínimo de difracción, desapareciendo del
patrón (\emph{órdenes perdidos}). Por ejemplo, si $d = 3a$, los órdenes $m
= 3, 6, 9, \ldots$ están ausentes. Este fenómeno es una manifestación directa
de la dualidad onda-corpúsculo: la cual rendija produce cada fotón no es
determinable, y la interferencia de los caminos posibles determina el patrón
observado.
\end{nota}

\subsection{Red de difracción}

Una red de difracción es un elemento periódico que divide el frente de onda
en $N$ fuentes secundarias igualmente espaciadas con separación $d$ (el
\emph{paso de red}). La interferencia de estas $N$ fuentes da lugar a
máximos muy agudos (\emph{máximos principales}) en las direcciones que
satisfacen la ecuación de red, y a $N-2$ máximos secundarios de
intensidad mucho menor entre cada par de principales. Cuanto mayor es $N$,
más agudos son los picos y mayor es el poder de resolución espectral.

Para $N$ rendijas idénticas:
\begin{equation}
  I(\theta)=I_1(\theta)
  \left[\frac{\sin(N\beta)}{\sin\beta}\right]^2,
  \qquad
  \beta=\frac{\pi d}{\lambda}\sin\theta.
  \label{eq:red_intensidad_t6}
\end{equation}
donde $I_1(\theta)$ es el patrón de difracción de una sola rendija (envolvente),
$N$ el número de rendijas iluminadas, $d$ el periodo de red (paso entre
rendijas) y $\theta$ el ángulo de observación.

El factor $[\sin(N\beta)/\sin\beta]^2$ alcanza el valor $N^2$ en los
\textbf{máximos principales}, que obedecen la ecuación de red:
\begin{equation}
  d\sin\theta_m=m\lambda, \qquad m = 0, \pm 1, \pm 2, \ldots
  \label{eq:red_ecuacion_t6}
\end{equation}
Entre dos máximos principales consecutivos hay exactamente $N-1$ ceros y
$N-2$ máximos secundarios con intensidad $\sim 1/N^2$ respecto al principal.

El \textbf{poder resolutivo} (criterio de Rayleigh para red) es:
\begin{equation}
  R=\frac{\lambda}{\Delta\lambda}=mN.
  \label{eq:red_resolucion_t6}
\end{equation}
donde $m$ es el orden de difracción y $\Delta\lambda$ la mínima separación
espectral resoluble alrededor de $\lambda$. El resultado $R = mN$ tiene una
interpretación física limpia: para resolver dos longitudes de onda en el
orden $m$, hace falta que la red ilumine al menos $R/m = N$ rendijas.

La red también permite cuantificar la \textbf{dispersión angular}:
\begin{equation}
  \frac{d\theta}{d\lambda}=\frac{m}{d\cos\theta},
  \label{eq:red_dispersion_t6}
\end{equation}
de modo que órdenes altos y pasos de red pequeños separan mejor longitudes
de onda cercanas en la pantalla.

\begin{observacion}
Las redes de difracción son el instrumento dispersivo más utilizado en
espectroscopía. Una red echelle de $N \sim 10^5$ ranuras trabajando en orden
$m \sim 50$ puede alcanzar $R \sim 5 \times 10^6$, suficiente para resolver
líneas espectrales separadas $\Delta\lambda \sim 10^{-4}\,\text{nm}$ en el
visible. La \emph{ordenación por blaze} (redes en ángulo) concentra la
energía difractada en el orden deseado.
\end{observacion}

\subsection{Interferómetro de Michelson}

El interferómetro de Michelson (1881) divide un haz mediante un
divisor semitransparente en dos brazos ortogonales. Cada haz se refleja
en un espejo y regresa al divisor, donde ambas ondas se recombinan e
interfieren. El diseño es especialmente versátil porque el desfase entre
brazos se puede variar de forma continua y controlada desplazando uno de
los espejos.

La lectura operacional es directa: el dispositivo convierte diferencias de
camino óptico en modulación de intensidad medible con alta sensibilidad.

\begin{nota}
El experimento de Michelson y Morley (1887) intentó medir el movimiento
de la Tierra respecto al éter luminífero comparando la velocidad de la
luz en direcciones perpendiculares. El resultado nulo fue uno de los
pilares experimentales que motivaron la teoría de la relatividad especial
de Einstein (1905). El mismo principio se usa hoy en los detectores de
ondas gravitacionales LIGO/Virgo, que miden variaciones relativas de
longitud del orden de $10^{-21}$.
\end{nota}

Para diferencia de longitudes de brazo $\Delta d=d_2-d_1$, el haz
reflejado en el brazo más largo recorre un camino extra $2\Delta d$,
por lo que la diferencia de fase es $\phi = 4\pi\Delta d/\lambda$:
\begin{equation}
  I=I_0\left[1+\cos\!\left(2k\Delta d\right)\right]
  =I_0\left[1+\cos\!\left(\frac{4\pi}{\lambda}\Delta d\right)\right].
  \label{eq:michelson_int_t6}
\end{equation}
donde $\Delta d=d_2-d_1$ es el desbalance entre brazos ópticos y
$I_0$ es la intensidad media.
Si se introduce una lámina de espesor $L$ e índice $n$ en un brazo, la fase
extra equivale a $\Delta_{\!L}=2(n-1)L$, de modo que:
\begin{equation}
  I=I_0\left[1+\cos\!\left(\frac{4\pi}{\lambda}(\Delta d+\Delta_{\!L})\right)\right].
  \label{eq:michelson_lamina_t6}
\end{equation}
donde $\Delta_{\!L}=2(n-1)L$ representa el incremento de camino óptico
introducido por la lámina.
Para franjas circulares:
\begin{equation}
  2d\cos\theta=m\lambda.
  \label{eq:michelson_anillos_t6}
\end{equation}
Aquí $d$ es la longitud de brazo efectiva del interferómetro y $\theta$ el
ángulo de observación del rayo.

\subsubsection{Conteo de franjas y sensibilidad}

Si un espejo se desplaza una cantidad $\Delta d$, la diferencia de camino óptico
varía en $2\Delta d$, y el número de franjas que atraviesan un punto fijo es:
\begin{equation}
  m=\frac{2\Delta d}{\lambda}.
  \label{eq:michelson_fringe_count_t6}
\end{equation}
Análogamente, al insertar una lámina de espesor $L$ e índice $n$ en un brazo:
\begin{equation}
  m=\frac{2(n-1)L}{\lambda}.
  \label{eq:michelson_lamina_count_t6}
\end{equation}
Estas relaciones explican por qué la interferometría permite medidas submicrométricas
de desplazamiento, espesor e índice.

\begin{figure}[ht]
  \centering
  \begin{tikzpicture}[scale=0.9]
    \draw[thick,fill=gray!10] (-0.22,-0.22) rectangle (0.22,0.22);

    \draw[thick] (-3,0) -- (-0.22,0);

    \draw[thick] (0.22,0) -- (3,0);
    \draw[thick] (3,-0.8) -- (3,0.8);

    \draw[thick] (0,0.22) -- (0,2.4);
    \draw[thick] (-0.8,2.4) -- (0.8,2.4);

    \draw[thick] (0,-0.22) -- (0,-2.4);
    \draw[thick] (-1.3,-2.4) -- (1.3,-2.4);

    \draw[->,thick] (-2.8,0.06) -- (-0.24,0.06);
    \draw[->,thick] (0.24,0.06) -- (2.8,0.06);
    \draw[->,thick] (2.8,-0.06) -- (0.24,-0.06);
    \draw[->,thick] (0.06,0.24) -- (0.06,2.25);
    \draw[->,thick] (-0.06,2.25) -- (-0.06,0.24);
    \draw[->,thick] (0,-0.24) -- (0,-2.2);
    \node[align=left, anchor=west, opticslabel] at (3.55,1.8)
      {BS: divisor de haz\\
       Brazo horizontal: espejo $E_1$\\
       Brazo vertical: espejo $E_2$\\
       Salida inferior: pantalla\\
       Flechas: trayectorias ópticas};
  \end{tikzpicture}
  \caption{Esquema básico del interferómetro de Michelson.}
  \label{fig:michelson_t6}
\end{figure}

\subsection{Interferómetro de Fabry-Perot}

El interferómetro de Fabry-Perot (1899) está formado por dos espejos planos
paralelos con reflectividad alta $R \lesssim 1$, separados una distancia $d$.
La luz que entra en la cavidad da lugar a múltiples reflexiones y transmisiones
parciales que interfieren en el plano de salida. Cuanto mayor sea $R$, mayor
es el número de reflexiones que contribuyen y más agudos son los picos de
transmisión, lo que permite resolver detalles espectrales muy finos.

Desde el punto de vista físico, Fabry-Perot actúa como un filtro resonante en
frecuencia: solo se transmiten con alta eficiencia las longitudes de onda que
satisfacen coherencia de fase tras múltiples recorridos.

La intensidad transmitida es la suma de todas las contribuciones parciales
(serie geométrica), que conduce a la \emph{función de Airy}:
\begin{equation}
  I_t=\frac{I_0}{1+F\sin^2\!\left(\frac{\phi}{2}\right)},
  \qquad
  F=\frac{4R}{(1-R)^2}.
  \label{eq:fabry_airy_t6}
\end{equation}
donde $R$ es la reflectividad de cada espejo, $F$ el coeficiente de Airy e $I_t$ la intensidad transmitida.
con fase:
\begin{equation}
  \phi=\frac{4\pi n d\cos\varepsilon'}{\lambda}.
  \label{eq:fabry_fase_t6}
\end{equation}
donde $n$ es el índice del medio interno, $d$ la separación entre espejos y
$\varepsilon'$ el ángulo interno respecto a la normal.
Máximos de transmisión:
\begin{equation}
  2nd\cos\varepsilon'=m\lambda.
  \label{eq:fabry_max_t6}
\end{equation}
Para una lente de focal $f'$, el radio de anillos puede escribirse como:
\begin{equation}
  \rho_m=f'\sqrt{1-\left(\frac{m\lambda}{2nd}\right)^2}.
  \label{eq:fabry_radio_t6}
\end{equation}

\subsubsection{Rango espectral libre, fineza y resolución}

Además de la forma de Airy, en Fabry-Perot son clave tres figuras de mérito:

\begin{description}
  \item[\textbf{Rango espectral libre (FSR):}] separación en frecuencia entre
        dos máximos principales sucesivos ($m$ y $m+1$):
\begin{equation}
  \Delta\nu_{\mathrm{FSR}}=\frac{c}{2nd\cos\varepsilon'}.
  \label{eq:fabry_fsr_t6}
\end{equation}
Si el espectro de la fuente se extiende más allá del FSR, los órdenes se
superponen y la interpretación se vuelve ambigua.

  \item[\textbf{Fineza:}] cociente entre FSR y anchura del pico a media
        altura $\delta\nu$:
\begin{equation}
  \mathcal{F}=\frac{\Delta\nu_{\mathrm{FSR}}}{\delta\nu}
  \simeq \frac{\pi\sqrt{R}}{1-R}.
  \label{eq:fabry_finesse_t6}
\end{equation}
La fineza cuantifica la agudeza relativa de los picos; depende solo de $R$
(no de $d$). Para $R = 0.99$, $\mathcal{F} \approx 312$.

  \item[\textbf{Poder resolutivo:}] para el orden interferencial $m$:
\begin{equation}
  \frac{\lambda}{\delta\lambda}\simeq m\,\mathcal{F}.
  \label{eq:fabry_res_t6}
\end{equation}
\end{description}

En consecuencia, aumentar $R$ estrecha los picos de transmisión y mejora la
resolución espectral del interferómetro.

\begin{observacion}
La fineza real está limitada por dos factores: (i) la fineza de reflexión
$\mathcal{F}_R = \pi R^{1/2}/(1-R)$, que aumenta con $R$ y en la práctica
requiere espejos de muy alta calidad; (ii) la fineza de planitud
$\mathcal{F}_P \simeq \lambda/(2\Delta_\lambda)$, donde $\Delta_\lambda$ es
la desviación máxima de planitud de cada espejo, que impone un límite
práctico $\mathcal{F}_P \lesssim D/(2\Delta_\lambda)$ para un espejo de
diámetro $D$.

Un etalón Fabry-Perot de laboratorio con $d = 10\,\text{mm}$, $n = 1$, y
$R = 0.99$ tiene $\text{FSR} \approx 15\,\text{GHz}$, $\mathcal{F} \approx
312$, $\delta\nu \approx 48\,\text{MHz}$, y en el orden $m \approx 30\,000$
(para $\lambda = 633\,\text{nm}$) un poder resolutivo $\approx 9.4 \times
10^6$.
\end{observacion}

\subsection{Marco de coherencia (complementario)}

La formulación que se ha usado en este tema —ondas monocromáticas con desfase
constante— es un idealization. Las fuentes reales emiten paquetes de onda de
duración finita $\tau_c$; distintos paquetes tienen diferencias de fase
aleatorias entre sí. El \emph{grado complejo de coherencia mútua}
$\gamma_{12}(\tau)$ generaliza la descripción al caso cuasi-monocromático:
\begin{equation}
  \gamma_{12}(\tau)
  = \frac{\langle E_1(t+\tau)\,E_2^*(t)\rangle}
         {\sqrt{\langle|E_1|^2\rangle\,\langle|E_2|^2\rangle}},
  \quad |\gamma_{12}| \in [0,1].
\end{equation}
Para $\tau = \Delta d/c$ (retardo asociado a la diferencia de camino):
\begin{equation}
  I=I_1+I_2+2\sqrt{I_1I_2}\,|\gamma_{12}|\cos(\Delta\phi+\arg\gamma_{12}),
  \qquad 0\le|\gamma_{12}|\le 1.
  \label{eq:gamma_interferencia_t6}
\end{equation}
El factor $|\gamma_{12}|$ reduce la visibilidad respecto al caso coherente:
cuando $|\tau| \gg \tau_c$, el promedio temporal hace que $\gamma_{12} \to 0$
y las franjas desaparecen completamente.

La visibilidad queda entonces:
\begin{equation}
  \mathcal{V}
  =\frac{2\sqrt{I_1I_2}}{I_1+I_2}\,|\gamma_{12}|.
  \label{eq:vis_gamma_t6}
\end{equation}
Para intensidades equilibradas ($I_1=I_2$), se obtiene
$\mathcal{V}=|\gamma_{12}|$: la visibilidad medida es directamente el módulo
del grado de coherencia. Esta conexión permite medir espectros de emisión
con el interferómetro de Michelson (transformada de Fourier de la función de
correlación), técnica conocida como espectroscopía FTIR.

\begin{observacion}
El grado de coherencia $|\gamma_{12}|$ depende simultáneamente de la
\emph{coherencia temporal} (relacionada con el ancho espectral $\Delta\nu$
y la longitud de coherencia $l_c = c/\Delta\nu$) y de la \emph{coherencia
espacial} (relacionada con el tamaño angular de la fuente). Un telescopio
que resuelve una estrella doble mide implícitamente $|\gamma_{12}|$ en
función de la separación de sus espejos (experimento de Michelson estelar,
1920).
\end{observacion}

\section{Formulario General}

\begin{mdframed}[style=resultado]
\textbf{Tema 1: Ondas en vacío}
\formulaname{Ecuaciones de Maxwell en vacío}
\[
\begin{aligned}
\div\ve{E} &= 0, &
\div\ve{B} &= 0, \\
\curl\ve{E} &= -\pdv{\ve{B}}{t}, &
\frac{1}{\mu_0}\curl\ve{B} &= \varepsilon_0\pdv{\ve{E}}{t}.
\end{aligned}
\]
\formulaname{Ecuación de onda electromagnética}
\[
\nabla^2\ve{A}=\frac{1}{v^2}\pdv[2]{\ve{A}}{t}
\]
\formulaname{Relaciones de propagación en vacío}
\[
k=\frac{\omega}{c}=\frac{2\pi}{\lambda},
\qquad
\lambda\nu=c,
\qquad
E=cB
\]
\formulaname{Polarización e intensidad}
\[
\ve{J}=
\begin{pmatrix}
E_{0x}\\E_{0y}e^{i\delta}
\end{pmatrix},
\qquad
I=\frac{E_0^2}{2\mu_0 c},
\qquad
\langle\ve{S}\rangle=\frac{1}{2\mu_0}\Re(\ve{E}_0\times\ve{B}_0^*)
\]
\end{mdframed}

\begin{mdframed}[style=resultado]
\textbf{Tema 2: Interacción radiación-materia}
\formulaname{Fuerza de Lorentz en régimen no relativista}
\[
\ve{F}\approx q\ve{E},
\qquad
|\ve{v}|\ll c
\]
\formulaname{Dipolo inducido y fuerza recuperadora}
\[
\mathbb{p}=q\ve{r},
\qquad
\ve{F}_{\text{rec}}=-m\omega_0^2\ve{r}
\]
\formulaname{Oscilador armónico amortiguado y forzado}
\[
m\ddot{\ve{r}}+m\Gamma\dot{\ve{r}}+m\omega_0^2\ve{r}
=q\ve{E}_0e^{i\omega t}
\]
\formulaname{Potencia transferida y potencia radiada}
\[
\langle P_{\text{trans}}\rangle=\left\langle q\ve{E}\cdot\dot{\ve{r}}\right\rangle,
\qquad
\langle P_{\text{esp}}\rangle\propto\omega^4|\mathbb{p}_0|^2
\]
\formulaname{Patrón dipolar y ley de Rayleigh}
\[
\frac{dP}{d\Omega}\propto\sin^2\theta,
\qquad
I_{\text{Rayleigh}}\propto\frac{1}{\lambda^4}
\]
\end{mdframed}

\begin{mdframed}[style=resultado]
\textbf{Tema 3: Propagación en medios materiales}
\formulaname{Criterio de dilución del medio}
\[
\frac{N_{\text{átomos}}}{\lambda^3}\le 1\ (\text{diluido}),
\qquad
\frac{N_{\text{átomos}}}{\lambda^3}>1\ (\text{denso})
\]
\formulaname{Fuentes macroscópicas promedio}
\[
\langle\rho\rangle=\rho_{\text{libre}}-\div\ve{P},
\qquad
\langle\ve{j}\rangle=\ve{j}_{\text{libre}}+\curl\ve{M}+\pdv{\ve{P}}{t}
\]
\formulaname{Relaciones constitutivas lineales}
\[
\ve{D}=\varepsilon_0\varepsilon_r\ve{E},
\qquad
\ve{B}=\mu_0\mu_r\ve{H},
\qquad
\ve{j}_{\text{libre}}=\sigma\ve{E}
\]
\formulaname{Índice complejo y permitividad efectiva}
\[
n_c^2=\varepsilon_{\text{gen}}
=\varepsilon_r\mu_r-i\frac{\sigma\mu_r}{\varepsilon_0\omega},
\qquad
n_c=\eta-i\kappa
\]
\formulaname{Vector de onda complejo y flujo de energía}
\[
\ve{k}_c=k_0n_c=\ve{k}-i\ve{a},
\qquad
\ve{S}=\ve{E}\times\ve{H}
\]
\end{mdframed}

\begin{mdframed}[style=resultado]
\textbf{Tema 4 y 4.2: Reflexión, refracción y anisotropía}
\formulaname{Leyes de reflexión y Snell}
\[
\theta_i=\theta_r,
\qquad
n_1\sin\theta_i=n_2\sin\theta_t
\]
\formulaname{Coeficientes de Fresnel (amplitud)}
\[
\begin{aligned}
r_\parallel&=\frac{n_2\cos\theta_i-n_1\cos\theta_t}{n_2\cos\theta_i+n_1\cos\theta_t},
&
r_\perp&=\frac{n_1\cos\theta_i-n_2\cos\theta_t}{n_1\cos\theta_i+n_2\cos\theta_t},
\\
t_\parallel&=\frac{2n_1\cos\theta_i}{n_2\cos\theta_i+n_1\cos\theta_t},
&
t_\perp&=\frac{2n_1\cos\theta_i}{n_1\cos\theta_i+n_2\cos\theta_t}.
\end{aligned}
\]
\formulaname{Ángulos característicos}
\[
\tan\theta_B=\frac{n_2}{n_1},
\qquad
\sin\theta_c=\frac{n_2}{n_1}\ (n_1>n_2)
\]
\formulaname{Reflectancia y transmitancia energética}
\[
\mathcal{R}_j=|r_j|^2,
\qquad
\mathcal{T}_j=\frac{n_2\cos\theta_t}{n_1\cos\theta_i}|t_j|^2,
\qquad
\mathcal{R}_j+\mathcal{T}_j=1
\]
\formulaname{Tensor dieléctrico e índices ordinario/extraordinario}
\[
\boldsymbol{\varepsilon}=
\begin{pmatrix}
\varepsilon_x&0&0\\
0&\varepsilon_y&0\\
0&0&\varepsilon_z
\end{pmatrix},
\qquad
n_o=\sqrt{\varepsilon_o},
\qquad
n_e=\sqrt{\varepsilon_e}
\]
\formulaname{Retardo de fase y formalismo de Jones}
\[
\delta=\frac{2\pi}{\lambda}(n_e-n_o)d,
\qquad
L_\alpha(\delta)=R^{-1}(\alpha)L(\delta)R(\alpha)
\]
\formulaname{Polarizador ideal y ley de Malus}
\[
P_\alpha=
\begin{pmatrix}
\cos^2\alpha & \cos\alpha\sin\alpha\\
\cos\alpha\sin\alpha & \sin^2\alpha
\end{pmatrix},
\qquad
I_s=I_0\cos^2(\theta-\alpha)
\]
\end{mdframed}

\begin{mdframed}[style=resultado]
\textbf{Tema 5: Difracción}
\formulaname{Integral de Huygens-Fresnel}
\[
U(P)=\frac{1}{i\lambda}\iint_A U_0(x_0,y_0)\frac{e^{-ikr}}{r}\cos\theta\,dx_0dy_0
\]
\formulaname{Aproximaciones de Fresnel y Fraunhofer}
\[
U_{\text{Fresnel}}(x,y)=\frac{e^{-ikz}}{i\lambda z}\iint_A U_0\,
e^{-i\frac{k}{2z}\left((x-x_0)^2+(y-y_0)^2\right)}dx_0dy_0
\]
\[
U_{\text{Fraunhofer}}(x,y)\propto\iint t(x_0,y_0)U_0(x_0,y_0)
e^{i\frac{k}{z}(xx_0+yy_0)}dx_0dy_0
\]
\formulaname{Número de Fresnel y principio de Babinet}
\[
N_F=\frac{a^2}{\lambda z},
\qquad
t_a+t_o=1
\]
\formulaname{Abertura rectangular (intensidad y mínimos)}
\[
I_{\text{rect}}(x,y)=I_0\operatorname{sinc}^2\!\left(\frac{kax}{2z}\right)
\operatorname{sinc}^2\!\left(\frac{kby}{2z}\right)
\]
\[
x_m=\frac{m\lambda z}{a},
\qquad
y_m=\frac{m\lambda z}{b}
\]
\formulaname{Rendija simple}
\[
I_{\text{rendija}}(x)=I_{\max}\operatorname{sinc}^2\!\left(\frac{kax}{2z}\right)
\]
\formulaname{Abertura circular y criterio de Rayleigh}
\[
I_{\text{circular}}(r)=I(0)\left[\frac{2J_1(\xi)}{\xi}\right]^2,
\qquad
\xi=\frac{kar}{f'}
\]
\[
r_{\text{Airy}}=\frac{1.22\,\lambda f'}{D},
\qquad
\Delta\theta_{\min}\simeq\frac{1.22\,\lambda}{D}
\]
\end{mdframed}

\begin{mdframed}[style=resultado]
\textbf{Tema 6: Interferencia}
\formulaname{Interferencia de dos haces y visibilidad}
\[
I=I_1+I_2+2\sqrt{I_1I_2}\cos\Delta\phi
\]
\[
\mathcal{V}=\frac{I_{\max}-I_{\min}}{I_{\max}+I_{\min}}
=\frac{2\sqrt{I_1I_2}}{I_1+I_2}
\]
\formulaname{Longitud y tiempo de coherencia}
\[
l_c=c\tau_c,
\qquad
\Delta\nu\,\tau_c\sim 1
\]
\formulaname{Doble rendija}
\[
d\sin\theta_m=m\lambda,
\qquad
x_m=\frac{m\lambda D}{d},
\qquad
\Delta x=\frac{\lambda D}{d}
\]
\[
I_{\text{doble rendija}}\propto
\operatorname{sinc}^2\!\left(\frac{kax}{2z}\right)\,
4\cos^2\!\left(\frac{kdx}{2z}\right)
\]
\formulaname{Red de difracción}
\[
I_{\text{red}}(\theta)=I_1(\theta)\left[\frac{\sin(N\beta)}{\sin\beta}\right]^2,
\qquad
\beta=\frac{\pi d}{\lambda}\sin\theta
\]
\[
\frac{\lambda}{\Delta\lambda}=mN,
\qquad
\frac{d\theta}{d\lambda}=\frac{m}{d\cos\theta}
\]
\formulaname{Interferómetro de Michelson}
\[
I_{\text{Michelson}}=I_0\left[1+\cos\!\left(\frac{4\pi}{\lambda}\Delta d\right)\right]
\]
\[
m_{\text{franjas}}=\frac{2\Delta d}{\lambda},
\qquad
m_{\text{lámina}}=\frac{2(n-1)L}{\lambda}
\]
\formulaname{Interferómetro de Fabry-Perot}
\[
I_t=\frac{I_0}{1+F\sin^2(\phi/2)},
\qquad
F=\frac{4R}{(1-R)^2},
\qquad
\phi=\frac{4\pi nd\cos\varepsilon'}{\lambda}
\]
\[
\Delta\nu_{\mathrm{FSR}}=\frac{c}{2nd\cos\varepsilon'},
\qquad
\mathcal{F}\simeq\frac{\pi\sqrt{R}}{1-R},
\qquad
\frac{\lambda}{\delta\lambda}\simeq m\mathcal{F}
\]
\formulaname{Interferencia parcial con coherencia compleja}
\[
I=I_1+I_2+2\sqrt{I_1I_2}\,|\gamma_{12}|
\cos(\Delta\phi+\arg\gamma_{12})
\]
\end{mdframed}

\end{document}
